\documentclass[11pt]{amsart}
\IfFileExists{lmodern.sty}%
  {\usepackage[T1]{fontenc}\usepackage{lmodern}\usepackage{microtype}}%
  {\usepackage[expansion=false]{microtype}}
\usepackage{amssymb,mathtools}
\usepackage{booktabs,array}
\usepackage{enumitem}
\emergencystretch=2.5em
\usepackage[colorlinks=true,linkcolor=blue,citecolor=blue,urlcolor=blue]{hyperref}

% ---------------------------------------------------------------------------
% Location of the Lean formalization.
\newcommand{\FormalizationURL}{https://github.com/slitsey/catalan-irrationality-measure}
\newcommand{\FormalizationLink}{\url{\FormalizationURL}}
% ---------------------------------------------------------------------------

\newtheorem{theorem}{Theorem}[section]
\setcounter{tocdepth}{1}
\newtheorem{lemma}[theorem]{Lemma}
\newtheorem{proposition}[theorem]{Proposition}
\newtheorem{corollary}[theorem]{Corollary}
\theoremstyle{definition}
\newtheorem{definition}[theorem]{Definition}
\theoremstyle{remark}
\newtheorem{remark}[theorem]{Remark}

\numberwithin{equation}{section}

\newcommand{\Z}{\mathbb Z}
\newcommand{\Q}{\mathbb Q}
\newcommand{\R}{\mathbb R}
\newcommand{\N}{\mathbb N}
\newcommand{\C}{\mathbb C}
\newcommand{\lean}[1]{\texttt{\detokenize{#1}}}
\newcommand{\den}{\operatorname{den}}
\renewcommand{\Re}{\operatorname{Re}}
\renewcommand{\Im}{\operatorname{Im}}
\newcommand{\Lm}{\Lambda}      % content constant 313/105
\newcommand{\Gram}{\mathcal G} % Gram matrix
\newcommand{\Rk}{\mathsf R}
\newcommand{\Fu}{\mathrm{Full}}
\newcommand{\F}{\mathbb F}

\begin{document}

\title{Irrationality of $G\pm5\pi^2/48$ and irrationality exponents}

\author{Sean Litsey}
\email{sean.litsey.math@gmail.com}

\date{October 18, 2026}

\subjclass[2020]{Primary 11J72; Secondary 11J82, 11Y60, 68V20}

\keywords{Catalan's constant, irrationality, irrationality exponent, L\"owner matrices,
Andr\'eief identity, $p$-adic nonvanishing, computer-assisted proof, Lean}

\begin{abstract}
Let $G$ be Catalan's constant and put $\lambda=G/2-5\pi^2/96$, so that
$G-5\pi^2/48=2\lambda$. We prove that $\lambda$ is irrational and that
$|\lambda-a/b|>b^{-46.45}$ for every integer $a$ and every sufficiently large integer $b$;
thus $\mu(\lambda)\le46.45$ (earlier versions of this paper gave $126$, $79.35$ and $57.85$). The mirror constant $\lambda'=G/2+5\pi^2/96$, so that $G+5\pi^2/48=2\lambda'$, is irrational with $\mu(\lambda')\le57.85$ (the hybrid improvement for $\lambda'$ has not been done; the two constants satisfy $\lambda+\lambda'=G$, but the two measures do not combine to a measure for $G$). The proof uses determinants of Löwner matrices of $\psi(z+1)-\psi(z+3/4)$ at integer nodes. In the square frame ($h\times h$ nodes, $\mu\le57.85$) the determinant $P_h(\lambda)$ of $T_h=\lambda I+M_0$ has common denominator with $\log\den P_h\le(313/105+o(1))h^2$, a contour representation and Andr\'eief's identity give $\log|P_h(\lambda)|\le-(3.1113\ldots-o(1))h^2$ (the energy step uses the conditionally positive definite kernel $-\frac12\log(|z-w|^2+\rho^2)$ and an interval-arithmetic certificate), a node-removal bound gives an exponentially wide window around $\lambda$, and an $\ell$-adic Schur-complement argument with a closed-form evaluation of a family of Krylov determinants gives nonvanishing. The improvement to $46.45$ uses a \emph{rectangular hybrid frame}: a common block of $r$ nodes together with two interleaved blocks of $m=r/8$ extra row and column nodes; the new ingredients are a sharper $p$-adic content bound (a rank saving at primes in two intervals, on top of an asymmetric tropical certificate), a decay and window certificate for the rectangular frame, and a bordered Krylov determinant, which is also evaluated in closed form. The Lean~4 formalization proves the irrationality and $\mu(\lambda)\le77.28$ (exponent $31682445459/409965000$, with the simpler choice of $h$), unconditionally; the mirror theorem is formalized with the exponents $78$ for $\lambda'$ and $79$ for $G+5\pi^2/48$. The determinant evaluation of \S\ref{sec:krylov}, the bound $57.85$ and the hybrid theorem $46.45$ are \emph{not} formalized.
\end{abstract}

\maketitle

\section{Introduction}\label{sec:intro}

OpenAI \cite{OpenAI2026} recently proved that Catalan's constant
\[
 G=\sum_{j=0}^\infty\frac{(-1)^j}{(2j+1)^2}=0.91596\,55941\ldots
\]
is irrational, and formalized the proof in Lean. (An earlier claimed proof
\cite{Sun2026} was found to contain a gap \cite{Wachs2026}.) Their argument, and the
irrationality-measure results built on it \cite{Jha2026,Carreira2026,Litsey2026}, use a
determinant of large size whose denominators and decay are balanced very closely.

In the present paper we study a different, much smaller, family of determinants, attached to the
number
\[
 \lambda \;=\; \frac G2-\frac{5\pi^2}{96}\;=\;\sum_{k\ge1}\Bigl(\frac1{16k^2}-\frac1{(4k-1)^2}\Bigr)
 \;=\;-0.05605\ldots
\]
(see \S\ref{sec:matrix}). Equivalently $G-5\pi^2/48=2\lambda$. We also treat the ``mirror'' constant $\lambda'=G/2+5\pi^2/96$ (\S\ref{sec:mirror}), for which $G+5\pi^2/48=2\lambda'$ and $\lambda+\lambda'=G$.

\begin{theorem}\label{thm:irr}
The number $\lambda=G/2-5\pi^2/96$ is irrational. Equivalently, $G-5\pi^2/48\notin\Q$.
\end{theorem}

\begin{theorem}\label{thm:main}
For every integer $a$ and every sufficiently large integer $b$,
\[
 \left|\lambda-\frac ab\right|>b^{-57.85}.
\]
In particular $\mu(\lambda)\le57.85$ (the computed value is $57.841$).
\end{theorem}

\begin{theorem}[the hybrid bound; Part~\ref{part:hybrid}]\label{thm:hybrid}
For every integer $a$ and every sufficiently large integer $b$,
\[
 \left|\lambda-\frac ab\right|>b^{-46.45}.
\]
In particular $\mu(\lambda)\le46.45$ (the computed value is $46.447$) and $\mu(G-5\pi^2/48)\le46.45$. This is proved in Part~\ref{part:hybrid} (\S\ref{sec:hframe}--\S\ref{sec:hfinal}) with a rectangular frame. It is \textbf{not formalized in Lean}, and it is established only for $\lambda$ (primes $\ell\equiv3\pmod4$): the hybrid for $\lambda'$ (primes $\ell\equiv1\pmod4$) has not been done, so Theorem~\ref{thm:mirror} below remains at $57.85$.
\end{theorem}
A simpler choice of $h$ (without the determinant evaluation of \S\ref{sec:krylov}) gives the weaker bound $77.13$; see Remark~\ref{rem:77}. The Lean formalization proves, unconditionally, the irrationality and $\mu(\lambda)\le77.28$ (exponent $31682445459/409965000$; with the simpler choice of $h$; the value $77.28$ rather than $77.13$ comes from rounding the certified constants to short rationals); the determinant evaluation of \S\ref{sec:krylov} and the bound $57.85$ are not formalized.

\begin{theorem}[the mirror constant]\label{thm:mirror}
The number $\lambda'=G/2+5\pi^2/96$ is irrational, and for every integer $a$ and every sufficiently large integer $b$, $|\lambda'-a/b|>b^{-57.85}$; thus $\mu(\lambda')\le57.85$. Consequently $G+5\pi^2/48\notin\Q$ and $\mu(G+5\pi^2/48)\le57.85$.
\end{theorem}
The proof of Theorem~\ref{thm:mirror} (\S\ref{sec:mirror}) is the same architecture with the frame $\psi(z+1)-\psi(z+\frac14)$; the constants (content $313/105$, the decay certificate, $c_{\rm cert}$) are identical, and only the $\ell$-adic part and some $o(1)$ bookkeeping change. Since $\lambda+\lambda'=G$ it is natural to ask whether Theorems~\ref{thm:main} and~\ref{thm:mirror} give a measure for $G$; they do not (Remark~\ref{rem:nocombine}). The Lean formalization does not cover Theorem~\ref{thm:mirror}.

Here the \emph{irrationality exponent} $\mu(\xi)$ of a real number $\xi$ is the infimum of the $\mu$ such that
$|\xi-a/b|>b^{-\mu}$ for all integers $a$ and all sufficiently large $b$.
We stress that Theorem~\ref{thm:irr} concerns one specific linear form in
$1,G,\pi^2$; it is not a statement about $G$ alone, whose irrationality is OpenAI's theorem.

A search of the literature (including work of Rivoal and Zudilin and of Calegari, Dimitrov and Tang on linear forms involving $G$) did not turn up a previous proof of the irrationality of $G-5\pi^2/48$ or of $G+5\pi^2/48$.

\emph{Related work on determinants.} Irrationality proofs built on determinants of large size with positivity (Hankel determinants of moment sequences) have a precedent in W.~Zudilin, \emph{A determinantal approach to irrationality} \cite{Zudilin}: that paper proves an irrationality criterion for sequences of integer approximations that carry Hankel structure, uses it to reprove the irrationality of $\pi$ and of $\zeta(3)$, and in its Section~5 ("Commentary") explains a limit of the method for integer-valued data: the integer transfinite diameter of the cube $[0,\varepsilon]^n$ tends to $\varepsilon/4$, so that the denominator--decay balance obtained from a Vandermonde-squared construction cannot be improved by that route. Its last section treats Catalan's constant and $q$-analogues of Apéry's constants. We do not use Zudilin's criterion; our matrices are Löwner matrices of a digamma difference rather than moment matrices, and the balance between denominator and decay is obtained from the $p$-adic and energy estimates below. We mention \cite{Zudilin} as prior art for the determinantal point of view, not as a source of any step of our proofs.

\subsection{Method}
For $h\ge1$ let $M_0$ be the explicit rational $h\times h$ matrix of \S\ref{sec:matrix} and $T_h=\lambda I+M_0$;
thus $P_h(X)=\det(XI+M_0)\in\Q[X]$ and $P_h(\lambda)=\det T_h$.
\begin{enumerate}[label=(\roman*),leftmargin=*]
\item \emph{Content} (\S\ref{sec:content}). A prime-by-prime tropical (assignment) bound, with an explicit
 dual certificate, gives $\log\den P_h\le(313/105+o(1))h^2$.
\item \emph{Decay} (\S\ref{sec:decay}). $T_h$ is $\tfrac1{16}$ times a L\"owner matrix. Moving the Hermite contour integral
 to a two-ray contour produces finitely many signed atoms plus an integral against an exponentially small
 analytic weight, and Andr\'eief's identity then bounds $|P_h(\lambda)|\le16^{-h}\det\Gram_h$ with $\Gram_h$ a
 positive definite Gram matrix. The energy of the resulting point configuration is bounded below with the kernel
 $k_\rho(z,w)=-\frac12\log(|z-w|^2+\rho^2)$, $\rho=1/1024$, which is conditionally positive definite on finite sums
 (Frullani's integral and Gaussian sums), against a certificate measure $\sigma$ of $300$ dyadic point masses; the
 resulting finite inequalities are checked in interval arithmetic and give $\det\Gram_h\le\exp(-(3.1113-o(1))h^2)$. The decay beats the content by the margin
 $3.1113790-313/105>0.1304$.
\item \emph{Nonvanishing} (\S\ref{sec:nonvanishing}). For a prime $\ell\equiv3\pmod4$, $\ell\ge7$, and any rational $r$
 with $\ell\nmid\den r$, one of $P_{\ell-1}(r),P_{\ell-2}(r),P_{(3\ell-1)/4}(r)$ is nonzero, because
 the $\ell$-adic Schur complement of $D(rI+M_0)D$ has a transparent Cauchy-matrix form.
 This gives irrationality with no information on the smallest eigenvalue of $T_h$.
\item \emph{A determinant evaluation and the choice of $h$} (\S\ref{sec:krylov}). Lemma~CC says that two consecutive
 Gram-type matrices $N_K,N_{K+1}$ over $\mathbb F_\ell$ have no common eigenvalue. It is reduced to the non-vanishing modulo $\ell$ of a Krylov determinant
 $u(K,K)=\det[w,N_Kw,\dots,N_K^{K-1}w]$, and Theorem~A evaluates $|u|$ as a product of powers of integers $\le2K<\ell$
 (zeros and poles in two shift parameters by Weyl's inequality and a cluster lemma, the constant by a recursion). Consequently, for any rational $r$
 with $\ell\nmid\den r$, one of $P_{\ell-K}(r),P_{\ell-K-1}(r)$ is nonzero, and the two values of $h$ are within a factor $1+O(1/\ell)$ of each other.
\item \emph{Measure} (\S\ref{sec:nodeline}--\ref{sec:measure}). The proof of the decay bound extends to every principal
 minor of $T_h$: removing $j=\beta h$ nodes costs at most $e^{c\beta h^2}$, with $c=2(m'-g_{\min})\le6.8717994$ and
 $g_{\min}$ certified in interval arithmetic (Theorem~\ref{thm:N}). Expanding $\det(T_h+tI)$ in principal minors, this gives
 $|P_h(\lambda+t)|\le e^{-(3.1113-o(1))h^2}$ for all $|t|\le e^{-ch}$. A rational $p/q$ very close to $\lambda$ lies in this window,
 and the arithmetic of the choice of $\ell$ and of $h\in\{\ell-K,\ell-K-1\}$ yields
 $\mu(\lambda)\le c\,(1/\eta+3/4)<57.85$, $\eta=3.1113790\ldots-313/105=0.1304266\ldots$ (the computed value is $57.841$; with the simpler choice of $h$ of Theorem~\ref{thm:nonvanish} one gets $c\,(4/(3\eta)+1)<77.13$, Remark~\ref{rem:77}; an earlier circle-smoothed energy step gave $79.35$ with that choice, Remark~\ref{rem:circle}).
\item[(vi)] \emph{The hybrid frame} (Part~\ref{part:hybrid}). Replace the square $h\times h$ L\"owner matrix by a rectangular one: rows $C\sqcup X'$, columns $C\sqcup Y'$, with a long common block $C$ of $r$ nodes and two interleaved blocks $X',Y'$ of $m=r/8$ nodes. The content constant rises to $\Lambda_F\le3.7211401$ (it is $313/105$ for $m=0$), but the decay constant rises further to $D_F\ge3.9062695$; the margin is $\eta_{\rm hyb}\ge0.1851294$ instead of $0.1304266$. The window constant is $c_{\rm hyb}\le7.6540647$ and the $\ell$-adic window contributes $\kappa=2/3$ instead of $3/4$: $\mu(\lambda)\le c_{\rm hyb}(1/\eta_{\rm hyb}+\kappa)\le46.45$.
\end{enumerate}

\subsection{Status of the proofs}
The proofs of Theorems~\ref{thm:irr}, \ref{thm:main} and~\ref{thm:mirror} were derived by Claude (Anthropic) agents
working with the author, and each part was re-derived and spot-checked numerically by separate Claude audit
agents; no human referee has checked them. Two steps are computer-assisted (interval-arithmetic certificates, \S\ref{ssec:cert} and
\S\ref{ssec:cert2}). The Krylov determinant evaluation (Theorem~A) and the choice of $h$ of \S\ref{sec:krylov} were derived and then audited by separate Claude agents (exact rational checks; see \S\ref{sec:formal}). The Lean formalization proves the irrationality and the weaker bound $77.13$, conditional on the numerical certificate facts; Theorem~A and the deduction of $57.85$ are not yet formalized, and the mirror theorem, Theorem~\ref{thm:mirror}, is formalized with the same weaker exponent ($78$ for $\lambda'$, $79$ for $G+5\pi^2/48$). See
\S\ref{sec:formal} for an exact account.

The hybrid theorem, Theorem~\ref{thm:hybrid}, and the whole of Part~\ref{part:hybrid} were derived by Claude (Anthropic) agents working with the author, in the same way, and each part was then audited by separate Claude agents (adversarial passes with independent exact computations; see \S\ref{sec:formal}); no human referee has checked them. One step is computer-assisted (the interval-arithmetic certificate of \S\ref{ssec:hcert}). \textbf{The hybrid theorem is not formalized in Lean}; the Lean results are the $57.85$-line results $77.28$ (for $\lambda$), $78$ (for $\lambda'$) and $79$ (for $G+5\pi^2/48$). A list of the statements for which our written proof is weaker than the rest is in \S\ref{ssec:htodo}.

\subsection{Notation}
$\psi$ is the digamma function, $H_n=\sum_{k\le n}1/k$, $H^{(2)}_n=\sum_{k\le n}1/k^2$, and
$\vartheta(x)$, $\vartheta(x;4,a)=\sum_{p\le x,\;p\equiv a\,(4)}\log p$ are Chebyshev functions.
For $P\in\Q[X]$, $\den P$ is the least positive integer $d$ with $dP\in\Z[X]$.
Steps resting on numerical certificates are labelled ``computer-assisted''.

%%% --------------------------------------------------------------------
\section{The matrices}\label{sec:matrix}

\subsection{The number $\lambda$}
Since $\sum_{n\equiv1(4)}n^{-2}-\sum_{n\equiv3(4)}n^{-2}=G$ and the sum of both is $\pi^2/8$ (sums over $n\ge1$),
\[
 \sum_{n\equiv3\,(4)}n^{-2}=\frac{\pi^2}{16}-\frac G2,\qquad
 \sum_{k\ge1}\frac1{16k^2}=\frac{\pi^2}{96},
\]
so
\begin{equation}\label{eq:lambda}
 \lambda:=\sum_{k\ge1}\Bigl(\frac1{16k^2}-\frac1{(4k-1)^2}\Bigr)=\frac{\pi^2}{96}-\frac{\pi^2}{16}+\frac G2=\frac G2-\frac{5\pi^2}{96}.
\end{equation}

\subsection{L\"owner structure}
Let
\[
 F(z)=\psi(z+1)-\psi(z+\tfrac34)=\sum_{j\ge0}\Bigl[\frac1{z+j+\frac34}-\frac1{z+j+1}\Bigr],
\]
meromorphic with simple poles at $-(j+\frac34)$ (residue $+1$) and $-(j+1)$ (residue $-1$), $j\ge0$.
For nodes $x_0,\dots,x_{h-1}$ the \emph{L\"owner matrix} of $F$ is
$L_F[a][b]=(F(x_b)-F(x_a))/(x_b-x_a)$ for $a\ne b$ and $L_F[a][a]=F'(x_a)$.
Take the nodes to be $x_a=a$, $a=0,\dots,h-1$.

From $F(k)-F(k-1)=\frac1k-\frac1{k-1/4}=-\frac1{k(4k-1)}$ and the corresponding identity for $F'$,
\begin{align}
 L_F[a][b]&=-\frac1{b-a}\sum_{a<k\le b}\frac1{k(4k-1)}\qquad(a<b),\label{eq:Loff}\\
 L_F[a][a]&=F'(0)+\sum_{k\le a}\Bigl(\frac{16}{(4k-1)^2}-\frac1{k^2}\Bigr),\label{eq:Ldiag}
\end{align}
and $F'(0)=\sum_{j\ge0}\bigl[(j+1)^{-2}-(j+\frac34)^{-2}\bigr]=\frac{\pi^2}6-16\sum_{n\equiv3(4)}n^{-2}=8G-\frac{5\pi^2}6=16\lambda$.

\begin{definition}
Let $M_0=M_0^{(h)}$ be the symmetric rational $h\times h$ matrix with, for $0\le a<b<h$,
\[
 M_0[a][b]=-\frac1{4(b-a)}\sum_{k=a+1}^{b}\frac1{4k(4k-1)},\qquad
 M_0[a][a]=-\sum_{k=1}^{a}\Bigl(\frac1{16k^2}-\frac1{(4k-1)^2}\Bigr),
\]
and set $T_h=\lambda I+M_0$ and $P_h(X)=\det(XI+M_0)\in\Q[X]$, monic of degree $h$.
\end{definition}

\begin{proposition}\label{prop:lowner}
$T_h=\frac1{16}L_F$ at the nodes $0,\dots,h-1$. In particular $P_h(\lambda)=16^{-h}\det L_F$, and
$P_h(\lambda+t)=\det(T_h+tI)$ for $t\in\R$.
\end{proposition}
\begin{proof}
Compare \eqref{eq:Loff}, \eqref{eq:Ldiag} with the definition, using $F'(0)=16\lambda$ and
$\frac1{16k^2}-\frac1{(4k-1)^2}=\frac1{16}\bigl(\frac1{k^2}-\frac{16}{(4k-1)^2}\bigr)$.
\end{proof}

\begin{remark}
The coefficients of $P_h$ are sums of principal minors of $M_0$, whose entries are
harmonic-type partial sums; there is no factorial closed form. Also $L_F$ is not
a Gram matrix of a positive measure, since the ``Stieltjes measure'' of $F$ consists of the signed
atoms $+1$ at $-(j+\frac34)$ and $-1$ at $-(j+1)$; this is the reason for the contour step in \S\ref{sec:decay}.
\end{remark}

\begin{lemma}[integrality]\label{lem:integrality}
Let $Q_h=\den(P_h)\cdot P_h\in\Z[X]$. If $a/b$ is in lowest terms and $Q_h(a/b)\ne0$, then $|Q_h(a/b)|\ge b^{-h}$.
\end{lemma}
\begin{proof}
$Q_h$ has degree $h$ and integer coefficients, so $b^hQ_h(a/b)\in\Z\setminus\{0\}$.
\end{proof}

%%% --------------------------------------------------------------------
\section{The denominator bound}\label{sec:content}

\begin{theorem}\label{thm:content}
$\log\den P_h\le\bigl(\tfrac{313}{105}+o(1)\bigr)h^2$ as $h\to\infty$.
\end{theorem}

We write $\Lm=313/105=2.98095\ldots$. The proof is prime by prime. For a prime $p$ put
$W_{ab}=\max(0,-v_p(M_0[a][b]))$.

\subsection{Tropical bound}
\begin{lemma}\label{lem:trop}
Let $\tau_p$ be the maximum of $\sum_aW_{a,\sigma(a)}$ over permutations $\sigma$ of $\{0,\dots,h-1\}$. Then
$v_p(\den P_h)\le\tau_p$. If $u_a,v_b\ge0$ satisfy $W_{ab}\le u_a+v_b$ for all $a,b$, then $\tau_p\le\sum_au_a+\sum_bv_b$.
\end{lemma}
\begin{proof}
The coefficient of $X^{h-j}$ in $\det(XI+M_0)$ is a sum of $j\times j$ principal minors, each a signed sum of
products $\prod_{a\in S}M_0[a][\sigma(a)]$ over permutations $\sigma$ of a $j$-set $S$. Each product has
$-v_p\le\sum_{a\in S}W_{a\sigma(a)}$, and since $W\ge0$, extending $\sigma$ by the identity off $S$ shows this is at most
$\tau_p$. For the second statement, $\sum_aW_{a\sigma(a)}\le\sum_a(u_a+v_{\sigma(a)})$.
\end{proof}

\subsection{Valuation majorant}
Let $p$ be odd and $\nu(k)=v_p(k(4k-1))=\max(v_p(k),v_p(4k-1))$ (the factors $k$ and $4k-1$ are coprime). Let $\mu(a)=\max_{1\le k\le a}\nu(k)$. By the ultrametric inequality applied to the formulas
of the Definition,
\begin{equation}\label{eq:major}
 W_{ab}\le v_p(b-a)+\max_{a<k\le b}\nu(k)\;(a<b),\qquad W_{aa}\le2\mu(a).
\end{equation}
For $q=p^e$ set
\begin{align*}
 \mathrm{div}_q(a,b)&=\mathbf 1[q\mid b-a,\ a\ne b],\\
 \mathrm{supp}_q(a,b)&=\mathbf 1[\exists k\in(\min(a,b),\max(a,b)]:\ q\mid k(4k-1)]\ (a\ne b),\qquad
 \mathrm{supp}_q(a,a)=2\cdot\mathbf 1[\mu(a)\ge e].
\end{align*}
Because $v_p(n)=\sum_{e\ge1}\mathbf 1[p^e\mid n]$ and $\max_k\nu(k)=\sum_e\mathbf 1[\exists k:\ p^e\mid k(4k-1)]$,
\eqref{eq:major} says $W\le\sum_{e\ge1}w^{(e)}$ entrywise with $w^{(e)}=\mathrm{div}_{p^e}+\mathrm{supp}_{p^e}\in\{0,1,2\}$.
Since a maximum of a sum is at most the sum of maxima, $\tau_p\le\sum_e\tau(w^{(e)})$.

\subsection{Per-level dual certificate}
Let $q=p^e$ be an odd prime power. The least $k\ge1$ with $q\mid k(4k-1)$ is $k_q=(j_qq+1)/4$,
where $j_q=1$ if $q\equiv3\pmod4$ and $j_q=3$ if $q\equiv1\pmod4$ (indeed $4k-1=jq$ forces $jq\equiv3\pmod 4$, and the branch $q\mid k$ gives $k\ge q>k_q$).
Hence $\mathrm{supp}_q(a,b)\ne0$ forces $\max(a,b)\ge k_q$. Put $n_q=\min(h,k_q)$ and define potentials
\[
 u_a=v_a=\begin{cases}1&a\ge k_q\ (\text{``high''}),\\ \mathbf 1[a\le h-1-q]&a<k_q\ (\text{``low''}).\end{cases}
\]

\begin{lemma}\label{lem:dual}
$w^{(e)}_{ab}\le u_a+v_b$ for all $a,b$. Consequently
\[
 \tau(w^{(e)})\le\delta(q):=\min\bigl(2h,\ 2(h-n_q)+2\min(n_q,(h-q)_+)\bigr),
\]
and $\delta(q)=0$ if $q\ge h$ and $k_q\ge h$; in particular $\tau_p=0$ for $p>4h$.
\end{lemma}
\begin{proof}
Note $k_q<q$. \emph{Both low:} $|b-a|<k_q<q$ so $\mathrm{div}=0$, and $(a,b]\subset[1,k_q-1]$ so $\mathrm{supp}=0$ (diagonal likewise).
\emph{Both high:} $w\le2=u_a+v_b$. \emph{$a$ low, $b$ high:} $\mathrm{supp}\le1=v_b$; if $\mathrm{div}=1$ then $b\ge a+q$, $b\le h-1$, so $a\le h-1-q$ and $u_a=1$; the case $a$ high, $b$ low is symmetric.
The diagonal is $0$ at low and at most $2$ at high indices. Now apply Lemma~\ref{lem:trop}: $\sum u+\sum v=2\#\{\text{high}\}+2\#\{a<n_q:\,a\le h-1-q\}$, and the trivial potentials $u=v\equiv1$ give $2h$.
\end{proof}
By Lemma~\ref{lem:trop}, $v_p(\den P_h)\le\sum_{e\ge1}\delta(p^e)$ for odd $p$.

\subsection{Summation}
\textbf{Prime $2$.} Every entry of $M_0$ has $-v_2\le4+2\log_2h$ (from the factors $1/(4(b-a))$, $1/(4k(4k-1))$, $1/(16k^2)$), so the prime $2$ contributes $O(h\log h)$ to $\log\den P_h$.

\textbf{Levels $e\ge2$.} Here $\delta\le2h$, so these contribute at most $2h\,(\sum_{p^e\le4h,e\ge2}\log p)=2h(\psi_{\rm Ch}(4h)-\vartheta(4h))=O(h^{3/2})$ (Chebyshev: $\psi_{\rm Ch}(x)-\vartheta(x)=O(\sqrt x)$).

\textbf{Level $e=1$.} For $q=p$ a direct computation from Lemma~\ref{lem:dual} (using $|k_q-j_qq/4|\le1/4$ and that $\delta$ is $2$-Lipschitz in $n_q$) gives $\delta(p)\le h\,g_j(p/h)+\tfrac12$, $j=j_p$, where
\[
 g_j(\theta)=\begin{cases}
 2&0\le\theta\le\frac4{4+j},\\
 4-\frac{(4+j)\theta}2&\frac4{4+j}\le\theta\le1,\\
 2-\frac{j\theta}2&1\le\theta\le\frac4j,\\
 0&\theta\ge\frac4j.\end{cases}
\]
Indeed for $\theta\le4/(4+j)$ one has $n_q\le h-q$ and $\delta=2h$; for $4/(4+j)\le\theta\le1$, $\delta=2(h-n_q)+2(h-q)$; for $\theta\ge1$, $\delta=2(h-n_q)$.
The function $g_j$ is continuous, non-increasing and piecewise linear, with $\int_0^\infty g_j=\frac4j+\frac j{4+j}$, that is
\[
 \int g_1=\frac{21}5\ (p\equiv3\bmod4),\qquad \int g_3=\frac{37}{21}\ (p\equiv1\bmod4).
\]
By partial summation, since $-g_j'\ge0$ is supported on $\theta\ge4/7$,
\[
 \frac1h\sum_{p\equiv a\,(4)}\log p\;g_j(p/h)=\frac1h\int(-g_j'(\theta))\,\vartheta(\theta h;4,a)\,d\theta\le\frac{1+\varepsilon_h}2\int g_j,
\]
where $\varepsilon_h=\sup_{4h/7\le x\le4h}\bigl(2\vartheta(x;4,a)/x-1\bigr)\to0$ by the prime number theorem in arithmetic progressions mod~$4$
(de la Vall\'ee Poussin; see e.g.\ \cite[Ch.~20]{Davenport} or \cite[Cor.~11.19]{MV}); only an upper bound for $\vartheta(x;4,a)$ is needed because the weights are non-negative.
Collecting terms,
\[
 \frac{\log\den P_h}{h^2}\le(1+\varepsilon_h)\Bigl(\frac12\cdot\frac{21}5+\frac12\cdot\frac{37}{21}\Bigr)+O\Bigl(\frac1{\sqrt h}\Bigr)
 =(1+\varepsilon_h)\frac{313}{105}+o(1).
\]
This proves Theorem~\ref{thm:content}. \qed

\begin{remark}[computer-assisted supplement; not used]
Evaluating the per-$h$ bound $\sum_q\delta(q)\log p$ (plus a bound for $p=2$) exactly for every $h\le2\cdot10^5$ gives
$\log\den P_h\le2.98635\,h^2$, with maximum at $h=1844$ (script \texttt{scan.py}). This is not needed for the
theorems. The dual certificate is also essentially sharp: for odd primes and $h=64,\dots,256$ the model value equals the exact tropical optimum except at $p=3$.
\end{remark}

%%% --------------------------------------------------------------------
\section{The decay bound}\label{sec:decay}

\subsection{A contour representation of $L_F$}
Hermite's formula gives, for any contour $C$ enclosing the nodes and no pole of $F$,
\[
 L_F=\frac1{2\pi i}\oint_CF(z)\,v(z)v(z)^{\!\top}dz,\qquad v_a(z)=\frac1{z-a}
\]
(the residue at $z=a\neq b$ gives $(F(a)-F(b))/(a-b)$, the double pole gives $F'(a)$).
By the reflection formulas $\psi(1+z)=\psi(-z)-\pi\cot\pi z$ and $\psi(\frac34+z)=\psi(\frac14-z)+\pi\cot(\frac\pi4-\pi z)$,
\begin{equation}\label{eq:reflect}
 F=\tilde F+E,\qquad \tilde F(z)=\psi(-z)-\psi(\tfrac14-z),\quad
 E(z)=-\pi\bigl[\cot\pi z+\cot(\tfrac\pi4-\pi z)\bigr]=\frac{-\pi\sin(\pi/4)}{\sin\pi z\,\sin(\pi/4-\pi z)}.
\end{equation}
Here $\tilde F$ is analytic in $\Re z<0$ and $O(1/|z|)$ there, while $E$ is $1$-periodic, has poles only on the real axis, and $|E(x+iy)|\le C e^{-2\pi|y|}$ for $|y|\ge1$.

Fix $r_0=7/50$, $N_0=\lfloor r_0h\rfloor$, $x^*=-N_0-\frac38$ (which lies in the gap $(-N_0-\frac34,-N_0)$ between poles of $F$). Let $\Gamma_h$ be the two rays
$x^*+u(-1\pm3i)$, $u\ge0$ (the \emph{rational direction} $(-1,\pm3)$, at angle $\pi-\arctan3=0.6024\pi$ from the positive axis), traversed downward, and let $\mathcal A=\{j+\frac34,\;j+1:\ 0\le j<N_0\}$, $|\mathcal A|=2N_0$. For $\mu\in\mathcal A$ let $c_\mu=+1$ for $\mu=j+\frac34$ and $-1$ for $\mu=j+1$, and $u_\mu(a)=1/(a+\mu)$.
The direction enters only through three facts: the left wedge bounded by $\Gamma_h$ has half-angle $\arctan3<\pi/2$, so $\Re z\le x^*<0$ on it; $E$ and $x^*\equiv5/8\pmod1$ are periodic; and on the part of $\Gamma_h$ with $|\Im z|<1$ (that is, $u<\frac13$) one has $\Re z\in(x^*-\frac13,x^*]=(-N_0-0.7083\ldots,-N_0-\frac38]$, strictly inside the gap between the poles $-N_0-\frac34$ and $-N_0$.

\begin{proposition}[contour identity]\label{prop:contour}
For every $h$,
\[
 L_F=-\sum_{\mu\in\mathcal A}c_\mu\,u_\mu u_\mu^{\!\top}+\frac1{2\pi i}\int_{\Gamma_h}E(z)\,v(z)v(z)^{\!\top}dz.
\]
\end{proposition}
\begin{proof}
Move the contour in Hermite's formula from around the nodes to $\Gamma_h$, closing on the right where the integrand is $O(|z|^{-3})$; this picks up the residues of $F$ at the $2N_0$ poles $-\mu$, $\mu\in\mathcal A$, which lie to the right of $x^*$, each contributing $\pm c_\mu u_\mu u_\mu^\top$ according to orientation. Then split $F=\tilde F+E$; the $\tilde F$ part integrates to $0$ over $\Gamma_h$, by closing in the left wedge where $\tilde Fvv^\top$ is analytic and $O(|z|^{-3})$.
\end{proof}
(The identity, with the stated orientation and signs, was checked entrywise to $10^{-17}$ with 25-digit arithmetic at $N_0=2$ and $N_0=7$, for the direction $(-1,3)$ and for the earlier angle $3\pi/5$.)

\subsection{Absolute-value Andr\'eief bound}
Let $|m|=\sum_{\mu\in\mathcal A}\delta_{-\mu}+\frac1{2\pi}|E(z)||dz|$ on $\Gamma_h$, and
\[
 \Gram_h=\sum_{\mu\in\mathcal A}u_\mu u_\mu^{\!\top}+\frac1{2\pi}\int_{\Gamma_h}|E|\,v\,v^*|dz|,
\]
a Hermitian positive definite matrix (the $v_a$ are linearly independent functions).
\begin{proposition}\label{prop:andreief}
$|\det L_F|\le\det\Gram_h$; hence $|P_h(\lambda)|\le16^{-h}\det\Gram_h$. More generally, for real $t$,
$|P_h(\lambda+t)|\le\det\bigl(\Gram_h/16+|t|I\bigr)$.
Moreover
\[
 \det\Gram_h=\frac1{h!}\int\prod_{i<j}|z_i-z_j|^2\prod_{a<b}(b-a)^2\prod_{i,a}|z_i-a|^{-2}\prod_i d|m|(z_i).
\]
\end{proposition}
\begin{proof}
By Proposition~\ref{prop:contour}, $L_F[a][b]=\int v_av_b\,dm$ for a complex measure $m$ with total variation $|m|$.
Andr\'eief's identity \cite{Forrester2018Andreief} gives
$\det L_F=\frac1{h!}\int\det[v_a(z_i)]^2\prod dm(z_i)$, so $|\det L_F|\le\frac1{h!}\int|\det[v_a(z_i)]|^2\prod d|m|(z_i)=\det\int v_a\overline{v_b}\,d|m|=\det\Gram_h$,
and the displayed formula is Cauchy's determinant $\det[1/(z_i-a)]=\pm\prod_{i<j}(z_i-z_j)\prod_{a<b}(b-a)\big/\prod(z_i-a)$.
For the last assertion, $T_h+tI$ is the Gram integral $\int\Phi_a\Phi_b\,dM$ on the disjoint union of $\Gamma_h\cup(-\mathcal A)$ and the node set, with $\Phi_a=v_a/4$ on the first part and $\Phi_a=\mathbf 1_{\{a\}}$ on the second, and $M=m\sqcup t\cdot\#$; the same argument bounds $|\det|$ by $\det(\Gram_h/16+|t|I)$.
\end{proof}

\subsection{The kernel $k_\rho$}\label{ssec:krho}
Fix $\rho=1/1024$ and put
\[
 k_\rho(z,w):=-\tfrac12\log\bigl(|z-w|^2+\rho^2\bigr)\qquad(z,w\in\C),\qquad k_\rho(z,z)=-\log\rho .
\]
The energy step below uses $k_\rho$ in place of the logarithmic kernel. Its virtue is that the conditional positive definiteness it needs is a statement about \emph{finite} sums and is proved by three elementary integrals (no Fourier transform and no measure on $\C$), and that no regularity constant for the potential of the certificate measure is ever needed.

\begin{lemma}[comparison]\label{lem:K0}
For $z\ne w$, $k_\rho(z,w)\le-\log|z-w|$.
\end{lemma}
\begin{proof}
$|z-w|^2<|z-w|^2+\rho^2$ and $\log$ is increasing.
\end{proof}

\begin{lemma}[Frullani form of the logarithm]\label{lem:F}
For $a>0$, $\displaystyle\log a=\int_0^\infty\frac{e^{-t}-e^{-at}}t\,dt$, and the integrand is integrable on $(0,\infty)$.
\end{lemma}
\begin{proof}
For $t>0$, $(e^{-t}-e^{-at})/t=\int_1^ae^{-tu}\,du$ (oriented integral); the integrand is bounded by $|a-1|$ near $0$ and by $2e^{-\min(1,a)t}$ on $[1,\infty)$. If $a\ge1$ the function $(t,u)\mapsto e^{-tu}$ on $(0,\infty)\times[1,a]$ is non-negative and $\int_1^a(\int_0^\infty e^{-tu}\,dt)\,du=\int_1^a du/u=\log a<\infty$, so Tonelli lets us exchange the integrals. If $0<a<1$, interchange with reversed orientation (the integrand is then $-\int_a^1e^{-tu}du$).
\end{proof}

\begin{lemma}[finite Gaussian sums]\label{lem:G}
For $t\ge0$, finitely many points $p_\alpha=(x_\alpha,y_\alpha)\in\R^2$ and reals $w_\alpha$, $\sum_{\alpha,\beta}w_\alpha w_\beta e^{-t|p_\alpha-p_\beta|^2}\ge0$.
\end{lemma}
\begin{proof}
Put $v_\alpha=w_\alpha e^{-t|p_\alpha|^2}$; the sum is $\sum v_\alpha v_\beta e^{2t\langle p_\alpha,p_\beta\rangle}=\sum_{n\ge0}\frac{(2t)^n}{n!}\sum_{\alpha,\beta}v_\alpha v_\beta\langle p_\alpha,p_\beta\rangle^n$. By the binomial theorem $\langle p,p'\rangle^n=\sum_{k=0}^n\binom nk(x^ky^{n-k})(x'^ky'^{n-k})$, hence
\[\sum_{\alpha,\beta}v_\alpha v_\beta\langle p_\alpha,p_\beta\rangle^n=\sum_{k=0}^n\binom nk\Bigl(\sum_\alpha v_\alpha x_\alpha^ky_\alpha^{n-k}\Bigr)^2\ge0,\]
a finite identity. A finite sum of convergent series is the series of the finite sums, and a series of non-negative terms is non-negative.
\end{proof}

\begin{lemma}[conditional positive definiteness]\label{lem:K}
Let $p_1,\dots,p_N\in\C$ (repetitions allowed) and $q_1,\dots,q_N\in\R$ with $\sum q_\alpha=0$. Then $\sum_{\alpha,\beta}q_\alpha q_\beta\,k_\rho(p_\alpha,p_\beta)\ge0$.
\end{lemma}
\begin{proof}
Put $a_{\alpha\beta}=|p_\alpha-p_\beta|^2+\rho^2>0$. By Lemma~\ref{lem:F} and linearity,
$\sum q_\alpha q_\beta(-\log a_{\alpha\beta})=\int_0^\infty\sum_{\alpha,\beta}q_\alpha q_\beta\frac{e^{-ta_{\alpha\beta}}-e^{-t}}t\,dt$.
The $e^{-t}$ part is $(\sum q_\alpha)^2e^{-t}/t=0$ pointwise. What remains is $\frac{e^{-t\rho^2}}t\sum q_\alpha q_\beta e^{-t|p_\alpha-p_\beta|^2}\ge0$ by Lemma~\ref{lem:G}. Multiply by $\frac12$.
\end{proof}

For a finite set $X\subset\C$ with $|X|=h$ and a certificate measure $\sigma=\sum_{k\le n_\sigma}m_k\delta_{c_k}$ ($m_k>0$, $\sum m_k=1$) put
\[
 U^\sigma_\rho(z)=\sum_km_k\,k_\rho(z,c_k),\qquad I_\rho(\sigma)=\sum_{k,l}m_km_l\,k_\rho(c_k,c_l).
\]
\begin{corollary}\label{cor:Kp}
$\displaystyle\frac1{h^2}\sum_{x\ne y\in X}k_\rho(x,y)\ \ge\ \frac2h\sum_{x\in X}U^\sigma_\rho(x)-I_\rho(\sigma)+\frac{\log\rho}h.$
\end{corollary}
\begin{proof}
Apply Lemma~\ref{lem:K} to the list $(X,c_1,\dots,c_{n_\sigma})$ with weights $(\frac1h,\dots,\frac1h,-m_1,\dots,-m_{n_\sigma})$, which sum to $0$. Expanding, $\frac1{h^2}\sum_{x,y\in X}k_\rho(x,y)-\frac2h\sum_xU^\sigma_\rho(x)+I_\rho(\sigma)\ge0$, and the diagonal terms $x=y$ contribute $\frac1{h^2}\cdot h\cdot(-\log\rho)$ to the left, which we move across.
\end{proof}

\subsection{The energy functional and the certificate}\label{ssec:cert}
Scale $z=hs$ and put $\tau=$ Lebesgue measure on $[0,1]$ and $U^\tau(s)=-\int_0^1\log|s-t|\,dt$, so that $I(\tau)=3/2$. For $\Re z<0$ and $y=\Im z\ge0$ one has the closed form
\[
 U^\tau(x+iy)=1-x\log|z|-(1-x)\log|1-z|-y\arctan\frac y{x^2+y^2-x},
\]
symmetric in $y\mapsto-y$; on the real line $U^\tau(t)=1-t\log|t|-(1-t)\log|1-t|$. Write $\Gamma$ for the rays $\{-r_0+u(-1\pm3i):u\ge0\}$ and $\kappa=10^{-3}$. For a fixed \emph{certificate measure} $\sigma$ define
\[
 g(s)=U^\sigma_\rho(s)-U^\tau(s)+(1-\kappa)\pi|\Im s|.
\]
The only requirement on $\sigma$ is that its quantities are computed rigorously; its relation to an optimizer is irrelevant.

\textbf{The measure $\sigma$.} $\sigma=\sum_{k=1}^{n_\sigma}m_k\delta_{c_k}$ is a sum of $n_\sigma=300$ point masses (84 real points and 108 conjugate pairs, so $\sigma$ and $g$ are conjugation-symmetric), with dyadic data: coordinates of $c_k$ of the form $X/2^{24}$, $Y/2^{24}$ with integers $X,Y$ of at most 23 bits, and masses $n_k/2^{32}$ with $\sum m_k=1$ \emph{exactly}. The data are in the file \texttt{sigma300.json}. Everything below is a finite computation with these rationals; in particular there are no circles, no Lipschitz constant and no continuum energy $I(\sigma)$.

\textbf{What is certified.} Put $\omega=2^{-30}$ (a box half-width absorbing the horizontal shift of $\Gamma_h/h$) and $u_{\max}=\frac32$.
\begin{enumerate}[label=(\textbf{\arabic*}),leftmargin=*]
\item[(R)] A cover of $u\in[0,u_{\max}]$ by intervals $[u_a,u_b]$ such that $g\ge m'$ on each box $[-r_0-u_b-\omega,-r_0-u_a+\omega]\times[3u_a,3u_b]$, together with the tail bound $T(u_{\max})\ge m'$: for $u\ge u_{\max}$, $|z|\le r_0+\omega+u\sqrt{10}$, $U^\sigma_\rho(z)\ge-\frac12\log((|z|+R_\sigma)^2+\rho^2)$ with $R_\sigma=\max|c_k|$, $U^\tau(z)\le-\log(|z|-1)$, and the field term is $\ge(1-\kappa)3\pi u$; the resulting lower bound $T(u)$ is increasing.
\item[(S)] A partition of $[-r_0,0]$ into cells $C_\ell=[a_\ell,b_\ell]$ with $g\ge\gamma_\ell$ on $C_\ell$. Put $S_{\rm cert}:=2\sum_\ell(b_\ell-a_\ell)\min(\gamma_\ell-m',0)$.
\item[(L)] A cover of $[0,1]$ (in fact of $[-\frac1{64},1]$) by cells on which $U^\sigma_\rho-U^\tau\ge g_{\min}$.
\end{enumerate}
Then put
\[
 D_{\rm cert}:=2(m'+S_{\rm cert})-I_\rho(\sigma)+\tfrac32,\qquad c:=2(m'-g_{\min}).
\]
On a box or cell the lower bound for $U^\sigma_\rho$ is $-\frac12\sum_km_k\log(d_{x,k}^2+d_{y,k}^2+\rho^2)$ with $d_{x,k},d_{y,k}$ the largest horizontal and vertical distances from $c_k$ to the box, which makes every logarithm argument as large as possible on the box (the two coordinates decouple, so the maximum is attained); $U^\tau$ is bounded above on ray boxes by ball arithmetic on the closed form, on segment cells by its value at the right end (it is increasing on $(-\infty,0]$), and on cells of $[0,1]$ by monotonicity on each side of $\frac12$ (the supremum is $U^\tau(b)$ if $b\le\frac12$, $U^\tau(a)$ if $a\ge\frac12$, and $1+\log2$ otherwise).

\textbf{Certified numbers (computer-assisted).} The script \texttt{kcert.py} (Arb \cite{Arb}, 96-bit precision; ray boxes bisected adaptively) was run on \texttt{sigma300.json} with the following results; $S_{\rm cert}$ is assembled from exact rationals, and all other entries are lower or upper bounds as indicated.
\begin{center}
\begin{tabular}{p{0.5\textwidth}p{0.3\textwidth}l}
\toprule
quantity & certified value & cells\\
\midrule
$I_\rho(\sigma)$ (exact finite sum) & $\le2.1107199269$ & $\binom{n_\sigma+1}2$ terms\\
$m'$: $\inf g$ over the ray boxes, $u\in[0,\frac32]$ & $\ge1.9718000062$ & $26{,}977$ boxes\\
tail $T(\frac32)$ & $\ge13.735$ (analytic) & --\\
$S_{\rm cert}$ & $\ge-0.1107505186$ & $14{,}000$ cells\\
$g_{\min}$ on $[-\frac1{64},1]$ & $\ge-1.4640996549$ & $4{,}423$ cells\\
\midrule
$D_{\rm cert}$ & $\ge3.1113790$ & \\
$c=2(m'-g_{\min})$ & $\le6.8717994$ & \\
$\eta=D_{\rm cert}-\Lm$ & $\ge0.1304266$ & \\
$c\,(1/\eta+3/4)$ (with Lemma~\ref{lem:3}) & $\le57.85$ (computed: $57.841$) & \\
$c\,(1/\eta+3/2)$ (without Lemma~\ref{lem:3}) & $<63$ (computed: $62.995$) & \\
$c\,(4/(3\eta)+1)$ (simpler $h$, Remark~\ref{rem:77}) & $\le77.1217$ (unrounded: $77.1212$) & \\
\bottomrule
\end{tabular}
\end{center}
(The rational data of the certificate were also re-derived in exact arithmetic by \texttt{mkdata.py}, with the masses rounded down to $2^{-32}$; this gave $D\ge3.11137905$ and $c\le6.87179932$.)
The certificate is valid for $h\ge1/\omega=2^{30}$. It is a computation about the fixed function $g$ and does not depend on $h$.

\subsection{The decay theorem}
\begin{lemma}[pointwise energy inequality $A_\rho$]\label{lem:A}
Let $w_\kappa=|E|^{1-\kappa}$ on $\Gamma_h$ and $w_\kappa=1$ at the atoms. For every configuration $z_1,\dots,z_h\in\Gamma_h\cup(-\mathcal A)$ of pairwise distinct points (otherwise the Cauchy determinant vanishes), writing $s_i=z_i/h$,
\begin{multline*}
 -\log\Bigl(\prod_{i<j}|z_i-z_j|^2\prod_{a<b}(b-a)^2\prod_{i,a}|z_i-a|^{-2}\prod w_\kappa(z_i)\Bigr)\\
 \ge\ h^2\Bigl[\tfrac2h\textstyle\sum_ig(s_i)-I_\rho(\sigma)+\tfrac32-\eta_1(h)\Bigr],
\end{multline*}
with $\eta_1(h)=\frac{|\log\rho|}h+\frac{2\log(1+8h/3)}h+\varepsilon_2(h)+\frac{\log^+C_E}h\to0$, where $C_E=\sup_{\Gamma_h}|E|e^{2\pi|\Im z|}<\infty$ is independent of $h$ and $\varepsilon_2(h)=O((\log h)/h)$.
\end{lemma}
\begin{proof}[Proof (this is the case $S=\{0,\dots,h-1\}$ of the proof of Theorem~\ref{thm:N})]
The exponent is $2h\log h+h^2H$, and the term $2h\log h\ge0$ is dropped. We write $H=H_1+H_2+H_3+H_4$.
(H$_1$) The $h$ points $s_i$ are distinct, so by Lemma~\ref{lem:K0} and Corollary~\ref{cor:Kp}, $\frac1{h^2}\sum_{i\ne j}-\log|s_i-s_j|\ge\frac1{h^2}\sum_{i\ne j}k_\rho(s_i,s_j)\ge\frac2h\sum_iU^\sigma_\rho(s_i)-I_\rho(\sigma)-\frac{|\log\rho|}h$.
(H$_2$) Since $\Re z_i\le-3/8$, $t\mapsto\log|s-t|$ is increasing on $[0,1]$ and the left Riemann sum for $U^\tau$ loses at most $\log(1+1/|s|)/h\le\log(1+8h/3)/h$; this gives $\frac2h\sum_i(-U^\tau(s_i))-\frac{2\log(1+8h/3)}h$.
(H$_3$) With $t_a=a/h$, $\frac1{h^2}\sum_{a\ne b}-\log|t_a-t_b|=\frac1{h^2}[h(h-1)\log h-2\sum_{n<h}\log n!]\ge\frac32-\varepsilon_2(h)$, using $\log n!\le(n+\frac12)\log n-n+1$.
(H$_4$) $-\log|E|^{1-\kappa}\ge(1-\kappa)(2\pi|\Im z|-\log^+C_E)$, with $C_E$ independent of $h$ because $E$ is $1$-periodic and $x^*\equiv5/8\pmod1$. Since $|\Im z|=h|\Im s|$, the sum of the field terms is $h^2\cdot\frac2h\sum_i(1-\kappa)\pi|\Im s_i|-h\log^+C_E$. Adding (H$_1$)--(H$_4$) gives the claim.
\end{proof}

\begin{lemma}[finite-cell bound $B_\rho$]\label{lem:B}
Let $h\ge2^{30}$ and let $s_1,\dots,s_k$ be the scaled points of any admissible configuration of $k\le h$ points of $\Gamma_h\cup(-\mathcal A)$ (atoms pairwise distinct). Then
\[
 \frac1h\sum_{i\le k}g(s_i)\ \ge\ \frac kh\,m'+S_{\rm cert}-\frac2h\sum_\ell(m'-\gamma_\ell)^+ .
\]
\end{lemma}
\begin{proof}
Write $\frac1h\sum_ig(s_i)=\frac khm'+\frac1h\sum_i(g(s_i)-m')$.
\emph{Contour points.} Write $x^*/h=-r_0+\delta_h$, $\delta_h=(r_0h-N_0-\frac38)/h\in[-\frac3{8h},\frac5{8h})$, so $|\delta_h|<1/h\le\omega$ and $s_i=-r_0+\delta_h+u(-1\pm3i)$. For $u\le u_{\max}$ the point $s_i$ (or its conjugate; $g$ is conjugation-symmetric) lies in a box of (R), and for $u>u_{\max}$ the tail bound applies; in either case $g(s_i)\ge m'$.
\emph{Atoms.} Atoms lie in $-(\Z_{\ge0}+\frac34)/h\,\cup\,-\Z_{\ge1}/h\subset[-r_0,0]$ and are distinct. Assign each to one cell containing it (an atom on a shared endpoint is assigned once). The two progressions have spacing $1/h$, so a closed cell $C_\ell$ receives $N_\ell\le2h(b_\ell-a_\ell)+2$ atoms. Since $g\ge\gamma_\ell$ on $C_\ell$ and $\min(\gamma_\ell-m',0)\le0$,
$\frac1h\sum_{\rm atoms}(g-m')\ge\sum_\ell\frac{N_\ell}h\min(\gamma_\ell-m',0)\ge S_{\rm cert}-\frac2h\sum_\ell(m'-\gamma_\ell)^+$.
Here $N_\ell/h\le2(b_\ell-a_\ell)+2/h$ and $2\sum_\ell(b_\ell-a_\ell)\min(\gamma_\ell-m',0)=S_{\rm cert}$.
\end{proof}

\begin{theorem}[decay; computer-assisted]\label{thm:decay}
$\displaystyle\det\Gram_h\le\exp\bigl(-(D^+-o(1))h^2\bigr)$ with $D^+=D_{\rm cert}\ge3.1113790$. Hence $|P_h(\lambda)|\le\exp(-(D^+-o(1))h^2)$.
\end{theorem}
\begin{proof}
Split $d|m|=w_\kappa\,d\Omega_\kappa$ with $\Omega_\kappa=\sum_{\mathcal A}\delta+\frac1{2\pi}|E|^\kappa|dz|$, of total mass $\le2h+c_\kappa$, $c_\kappa=\frac1{2\pi}\int_{\Gamma_0}|E|^\kappa|dz|<\infty$ independent of $h$ by periodicity. By Lemmas~\ref{lem:A} and~\ref{lem:B} (with $k=h$), the integrand of Proposition~\ref{prop:andreief} against $\prod w_\kappa$ is at most $\exp(-h^2[D_{\rm cert}-\eta_1-2\eta_2])$ pointwise, where $\eta_2(h)=\frac2h\sum_\ell(m'-\gamma_\ell)^+$ (the cell number is fixed, so this is $O(1/h)$). Therefore
\[
 \det\Gram_h\le\frac1{h!}(2h+c_\kappa)^h\exp\bigl(-h^2[D_{\rm cert}-\eta_1-2\eta_2]\bigr),
\]
and $\log(2h+c_\kappa)/h\to0$. Finally $|P_h(\lambda)|\le16^{-h}\det\Gram_h$ by Proposition~\ref{prop:andreief}.
\end{proof}

\begin{remark}
The error terms $\eta_1,\eta_2$ are $O((\log h)/h)$ with moderate explicit constants; no explicit $h_0$ is claimed (the certificate itself requires $h\ge2^{30}$). The exponent $D^+\ge3.1113790$ is certified; numerics suggest that the true limit is $\frac92\log2=3.11916\ldots$ (not used). The margin for the deductions below is
\[
 \eta:=D^+-\Lm\ \ge\ 3.1113790-2.9809524=0.1304266,
\]
and the argument of \S\ref{sec:measure} uses the decay theorem only through this margin.
\end{remark}

\begin{remark}[cross-check by circle smoothing]\label{rem:circle}
An earlier version of the energy step replaced each point by a uniform measure on a circle of radius $1/h$ and used a certificate measure $\sigma$ supported on $448$ disjoint circles together with the inequality $I(\nu-\sigma)\ge0$ for signed measures of mass zero \cite[Lemma~I.1.8]{SaffTotik} and a Lipschitz bound for $U^\sigma$ ($L_\sigma\le1522.04$). That route, with an independent interval-arithmetic certificate, gives $D^+\ge3.1072419$, $m'\ge1.9681060$, $S\ge-0.1101022$, $g_{\min}\ge-1.4643531$, $c\le6.8650$ and $\mu(\lambda)\le79.35$ (the value $79.3439$). We keep it only as an independent cross-check of the decay exponent: the two certificates use different kernels, different measures $\sigma$ and a different ray angle ($3\pi/5$), and agree on $D^+\approx3.11$. The $k_\rho$ route gives the larger $D^+$ because its certificate measure was re-optimised.
\end{remark}


%%% --------------------------------------------------------------------
\section{$\ell$-adic nonvanishing}\label{sec:nonvanishing}

Let $\ell\equiv3\pmod4$ be a prime with $\ell\ge7$, $K=(\ell+1)/4$ (an integer $\ge2$), and let $K+1\le h\le\ell-1$. Let $\Z_{(\ell)}$ be the localization of $\Z$ at $\ell$.
Call the indices $a<K$ \emph{low} and $a\ge K$ \emph{high}.

\begin{lemma}[$\ell$-singular structure]\label{lem:struct}
Among the quantities $k,\;4k-1,\;b-a,\;4$ occurring in the entries of $M_0$ ($0\le a<b<h$, $1\le k\le h-1$) only $4k-1$ can be divisible by $\ell$, and this occurs only for $k=K$ (with $4K-1=\ell$). Consequently every high diagonal entry of $M_0$ contains the term $+\ell^{-2}$, every low--high entry contains $-1/(16K\ell(b-a))$, and all other entries lie in $\Z_{(\ell)}$.
\end{lemma}
\begin{proof}
$k,b-a<\ell$. If $\ell\mid4k-1$ with $4k-1\le4h-5<4\ell$ then $4k-1\in\{\ell,2\ell,3\ell\}$; $2\ell$ is even and $3\ell\equiv1\pmod4$, so $4k-1=\ell$, i.e.\ $k=K$.
The sum $\sum_{a<k\le b}$ contains $k=K$ iff $a<K\le b$ (low--high); the diagonal sum $\sum_{k\le a}$ contains it iff $a\ge K$.
\end{proof}

\begin{lemma}\label{lem:1}
Let $D=\mathrm{diag}(I_K,\ell I_{h-K})$. Then $D(XI+M_0)D\in M_h(\Z_{(\ell)}[X])$ and, modulo $\ell$,
\[
 D(XI+M_0)D\equiv\begin{pmatrix}XI_K+A&\hat B\\\hat B^\top&I\end{pmatrix},\qquad
 A=M_0|_{\rm low}\bmod\ell,\quad \hat B_{ab}\equiv-\frac1{4(b-a)}\ (a<K\le b),
\]
and hence $\ell^{2(h-K)}P_h(X)\equiv\det(XI_K+A-\hat B\hat B^\top)\pmod\ell$. The right-hand side is monic of degree $K$; in particular $v_\ell(\den P_h)=2(h-K)$.
\end{lemma}
\begin{proof}
By Lemma~\ref{lem:struct}: the high diagonal entries times $\ell^2$ are $\equiv1$ (the term $\ell^{-2}$), the low--high entries times $\ell$ are $\equiv-1/(16K(b-a))\equiv-1/(4(b-a))$ since $16K=4\ell+4$, the high--high off-diagonal entries times $\ell^2$ and the term $X\ell^2$ are $\equiv0$, and the low block is unchanged. Take the Schur complement of the identity block; $\det D=\ell^{h-K}$.
\end{proof}

\begin{lemma}[key identity]\label{lem:2}
Let $s=\ell-h\in[1,\ell-K-1]$, $m=\ell-K$, and $C_s=\bigl[1/(a+j)\bigr]_{0\le a<K,\ 1\le j\le s}$ (entries in $\mathbb F_\ell$; note $a+j\le K-1+s<\ell$). Put $c=-H^{(2)}_m\in\mathbb F_\ell$ and $Y=16X+c$. Then, in $\mathbb F_\ell[X]$,
\[
 16^{K}\ell^{2(h-K)}P_h(X)\equiv\det\bigl(Y\,I_K+C_sC_s^\top\bigr).
\]
\end{lemma}
\begin{proof}
Modulo $\ell$, $4K=\ell+1\equiv1$, so $\frac4{4k-1}=\frac1{k-1/4}\equiv\frac1{k-K}\equiv\frac1{k+m}$ and thus
$-\frac1{k(4k-1)}\equiv\frac1k-\frac1{k+m}$ and $\frac{16}{(4k-1)^2}\equiv\frac1{(k+m)^2}$.
For $a<b<K$ (so $b+m<\ell$) summing over $a<k\le b$ and using $\frac1{(b-a)}\bigl(\frac1{a+j}-\frac1{b+j}\bigr)=\frac1{(a+j)(b+j)}$,
\[
 16A_{ab}\equiv\frac1{b-a}\sum_{k=a+1}^{b}\Bigl(\frac1k-\frac1{k+m}\Bigr)=\sum_{j=1}^{m}\frac1{(a+j)(b+j)},\qquad
 16A_{aa}\equiv\sum_{k\le a}\Bigl(\frac1{(k+m)^2}-\frac1{k^2}\Bigr)=\sum_{j=1}^m\frac1{(a+j)^2}-H^{(2)}_m.
\]
Thus $16A\equiv cI+C_mC_m^\top$. Next, $16(\hat B\hat B^\top)_{aa'}\equiv\sum_{b=K}^{h-1}\frac1{(b-a)(b-a')}$. Substituting $j=\ell-b$ (which runs over $s+1,\dots,m$ as $b$ runs over $h-1,\dots,K$) we have $a+j\equiv-(b-a)$, so this equals $\sum_{j=s+1}^m\frac1{(a+j)(a'+j)}$. Subtracting,
$16(A-\hat B\hat B^\top)\equiv cI+C_sC_s^\top$, and Lemma~\ref{lem:1} gives the claim, since $\det(XI+(cI+C_sC_s^\top)/16)=16^{-K}\det(YI+C_sC_s^\top)$.
\end{proof}
(This was checked by exact characteristic-polynomial reduction for $h=12,20$ and every $\ell\equiv3\pmod4$ in $(h,4h]$.)

Write $\Xi_s=C_s^\top C_s$, i.e.\ $(\Xi_s)_{ij}=\sum_{a=0}^{K-1}\frac1{(a+i)(a+j)}$. By Sylvester's determinant identity, for $s\le K$,
\begin{equation}\label{eq:sylv}
 \det(YI_K+C_sC_s^\top)=Y^{K-s}\det(YI_s+\Xi_s).
\end{equation}

\begin{theorem}[$\ell$-adic nonvanishing]\label{thm:nonvanish}
Let $\ell\equiv3\pmod4$ be prime, $\ell\ge7$, $K=(\ell+1)/4$, and let $r\in\Q$ with $\ell\nmid\den r$. Then at least one of
\[
 h_1=\ell-1,\qquad h_2=\ell-2,\qquad h_3=\ell-K=\tfrac{3\ell-1}4
\]
satisfies $v_\ell\bigl(P_{h}(r)\bigr)=-2(h-K)$; in particular $P_h(r)\ne0$.
\end{theorem}
\begin{proof}
All three $h_i$ satisfy $K+1\le h_i\le\ell-1$ (for $h_3$: $\frac{3\ell-1}4\ge\frac{\ell+5}4\iff\ell\ge3$), and have $s=1,2,K$ respectively, all $\le K$ since $K\ge2$.
By Lemmas~\ref{lem:1} and~\ref{lem:2}, $\ell^{2(h-K)}P_h(r)$ is an $\ell$-adic unit iff the polynomial \eqref{eq:sylv} does not vanish at $Y=y:=16r+c\in\mathbb F_\ell$. Suppose it vanishes at $y$ for all three values of $s$.
\begin{itemize}[leftmargin=*]
\item $s=K$: at $Y=0$ the value of \eqref{eq:sylv} is $\det(\Xi_K)=\det(C_K)^2$, where $C_K=[1/(a+j)]$ is a Cauchy matrix with $\det C_K=\prod_{a<a'}(a'-a)\prod_{j<j'}(j'-j)/\prod_{a,j}(a+j)$, a ratio of products of integers in $[1,2K-1]$, and $2K-1<\ell$. So $\det C_K\ne0$ in $\mathbb F_\ell$, the polynomial does not vanish at $0$, and hence $y\ne0$.
\item $s=1$: the value is $y^{K-1}(y+\Xi_{11})$ with $\Xi_{11}=\sum_{a<K}(a+1)^{-2}=H_K^{(2)}$. As $y\ne0$, $y=-\Xi_{11}\ne0$.
\item $s=2$: the value is $y^{K-2}\det\begin{pmatrix}0&\Xi_{12}\\\Xi_{12}&\Xi_{22}-\Xi_{11}\end{pmatrix}=-y^{K-2}\Xi_{12}^2$, and $\Xi_{12}=\sum_{a<K}\bigl(\frac1{a+1}-\frac1{a+2}\bigr)=\frac K{K+1}\ne0$ in $\mathbb F_\ell$ since $1\le K<K+1<\ell$. This is nonzero, a contradiction.
\end{itemize}
(If $\ell=7$ then $h_2=h_3$, $K=2$, and the argument is unchanged.) Hence for some $i$ the reduction is nonzero, so $v_\ell(P_{h_i}(r))=-2(h_i-K)$.
\end{proof}
(Exact characteristic polynomials for $\ell=7,11,19,23,31,43$ confirm that the three reductions have $\gcd=1$ and the stated valuations.)

\begin{lemma}[primes dividing the denominator]\label{lem:3}
Let $\ell\equiv3\pmod4$ be prime, $\ell\ge7$, $K+1\le h\le\ell-1$, and $r=p/q$ with $e=v_\ell(q)\ge1$; put $u=\ell^er\in\Z_{(\ell)}^\times$. If $e=1$, or $e\ge3$, or $e=2$ and $u\not\equiv-1\pmod\ell$, then $P_h(r)\ne0$ \emph{for every admissible $h$}.
\end{lemma}
\begin{proof}
Let $N=D(rI+M_0)D$. By Lemmas~\ref{lem:struct}, \ref{lem:1}, its blocks are $rI+M_{0,LL}$, $\ell M_{0,LH}\equiv\hat B$, and $\ell^2rI+J$ with $J\equiv I$, all entries other than the $r$'s being $\ell$-integral. Multiply the low rows by $\ell^e$ (and, if $e\ge3$, the high rows by $\ell^{e-2}$). Modulo $\ell$ the matrix becomes
$\begin{pmatrix}uI&0\\\hat B^\top&I\end{pmatrix}$ for $e=1$ (determinant $u^K$),
$\begin{pmatrix}uI&0\\\hat B^\top&(u+1)I\end{pmatrix}$ for $e=2$ (determinant $u^K(u+1)^{h-K}$),
$uI_h$ for $e\ge3$ (determinant $u^h$). These are units in the stated cases.
\end{proof}
(Checked on $1044$ exact cases; the exception $e=2$, $u\equiv-1$ does occur.)

\subsection{Proof of Theorem~\ref{thm:irr}}
Suppose $\lambda=a/b$ in lowest terms. Choose a prime $\ell\equiv3\pmod4$, $\ell\ge7$, $\ell\nmid b$, as large as we please (there are infinitely many primes $\equiv 3\pmod 4$). Theorem~\ref{thm:nonvanish} gives $h\ge(3\ell-1)/4$ with $P_h(\lambda)\ne0$. By Lemma~\ref{lem:integrality}, Theorem~\ref{thm:content} and Theorem~\ref{thm:decay},
\[
 1\le b^h\den(P_h)\,|P_h(\lambda)|\le\exp\bigl(h\log b-(D^+-\Lm-o(1))h^2\bigr),
\]
which is false for large $h$ because $D^+-\Lm>0.130$. \qed

%%% --------------------------------------------------------------------
\section{A determinant evaluation and the choice of $h$}\label{sec:krylov}

Theorem~\ref{thm:nonvanish} uses three values of $h$ whose ratio $h_{\max}/h_{\min}$ is $4/3$ (the factor $4/3$ in the exponent of Theorem~\ref{thm:measure}). In this section we show that already the two values
\[
 h\in\{\ell-K,\ \ell-K-1\}\qquad(\text{that is, } s=\ell-h\in\{K,K+1\})
\]
suffice; their ratio is $1+O(1/\ell)$. The proof rests on a statement about two consecutive Gram-type matrices over $\mathbb F_\ell$ (Lemma~CC below), which in turn reduces to the non-vanishing modulo~$\ell$ of an explicit Krylov determinant; the latter is evaluated in closed form (Theorem~A). By~\eqref{eq:sylv} the polynomial $\det(YI_K+C_sC_s^\top)$ has the factor $Y^{K-s}$ for $s<K$, so the value $y=0$ cannot be excluded for any $h>\ell-K$; this is why the smallest usable $s$ is $K$.

Throughout this section $\ell\equiv3\pmod 4$ is a prime, $\ell\ge7$, $K=(\ell+1)/4$, and for $s\ge0$ with $K+s<\ell$ we use, over $\mathbb F_\ell$,
\[
 C_s=\Bigl[\frac1{a+j}\Bigr]_{0\le a<K,\;1\le j\le s},\qquad N_s=C_sC_s^\top\in M_K(\mathbb F_\ell),\qquad w_j=\Bigl(\frac1{a+j}\Bigr)_{0\le a<K},
\]
so that $N_{s+1}=N_s+w_{s+1}w_{s+1}^\top$. All denominators are $\ell$-units because $a+j\le K-1+s+1<\ell$.

\subsection{Statement}

\begin{lemma}[Lemma CC]\label{lem:CC}
Let $s\ge K$ with $K+s<\ell$. Then $N_s$ and $N_{s+1}$ have no common eigenvalue in $\overline{\mathbb F}_\ell$. Equivalently, the polynomials $\det(YI_K+N_s)$ and $\det(YI_K+N_{s+1})$ have no common root in $\overline{\mathbb F}_\ell$. In particular this holds for $s=K$ (note $2K<\ell$).
\end{lemma}

The proof (\S\ref{ssec:CCproof}) is the combination of the Krylov reduction of \S\ref{ssec:kryred} with the evaluation of Theorem~A below.

\begin{theorem}[Choice of $h$]\label{thm:nonvanish2}
Let $\ell\equiv3\pmod4$ be prime, $\ell\ge7$, $K=(\ell+1)/4$, and let $r\in\Q$ with $\ell\nmid\den r$. Then at least one of
\[
 h_1=\ell-K=\tfrac{3\ell-1}4,\qquad h_2=\ell-K-1=\tfrac{3\ell-5}4
\]
satisfies $v_\ell\bigl(P_h(r)\bigr)=-2(h-K)$; in particular $P_h(r)\ne0$.
\end{theorem}

\subsection{The Krylov reduction}\label{ssec:kryred}

\begin{lemma}\label{lem:kr1}
Let $F$ be a field of characteristic $\ne2$, $N\in M_K(F)$ symmetric, $w\in F^K$. If $z\in F$ is a common root of $\det(z-N)$ and $\det(z-N-ww^\top)$, then there is $x\ne0$ with $Nx=zx$ and $w^\top x=0$.
\end{lemma}
\begin{proof}
The matrix determinant lemma gives $\det(z-N-ww^\top)=\det(z-N)-w^\top\operatorname{adj}(z-N)\,w$, so $w^\top\operatorname{adj}(z-N)\,w=0$. Since $\det(z-N)=0$, $\operatorname{rank}(z-N)\le K-1$. If $\operatorname{rank}(z-N)\le K-2$, then $\ker(z-N)$ has dimension $\ge2$ and so contains some $x\ne0$ with $w^\top x=0$. If $\operatorname{rank}(z-N)=K-1$, then $\operatorname{adj}(z-N)$ is symmetric of rank one with columns in $\ker(z-N)=\langle x\rangle$, so $\operatorname{adj}(z-N)=\gamma\,xx^\top$ with $\gamma\ne0$, and $\gamma(w^\top x)^2=0$ forces $w^\top x=0$.
\end{proof}
(This covers isotropic eigenvectors, $x^\top x=0$, which genuinely occur modulo $\ell$.)

\begin{lemma}\label{lem:kr2}
Let $N\in M_K(F)$ be symmetric and $w\in F^K$. If some $x\ne0$ satisfies $Nx=zx$ and $w^\top x=0$, then $\det[w,Nw,\dots,N^{K-1}w]=0$.
\end{lemma}
\begin{proof}
$x^\top N^kw=(N^kx)^\top w=z^kx^\top w=0$ for every $k$, so $x$ lies in the left kernel of the Krylov matrix.
\end{proof}

\begin{corollary}\label{cor:kr}
Let $s\ge1$, $K+s<\ell$, and $u(K,s):=\det[w_{s+1},N_sw_{s+1},\dots,N_s^{K-1}w_{s+1}]\in\mathbb F_\ell$. If $u(K,s)\ne0$, then $N_s$ and $N_{s+1}$ have no common eigenvalue in $\overline{\mathbb F}_\ell$.
\end{corollary}
\begin{proof}
Apply Lemmas~\ref{lem:kr1} and~\ref{lem:kr2} over $F=\overline{\mathbb F}_\ell$ with $N=N_s$ and $w=w_{s+1}$.
\end{proof}

\subsection{Theorem A}

For real $L,\rho\notin\{0,-1,-2,\dots\}$ and $n\ge1$ let $c_y=(1/(a+y))_{a=0}^{n-1}\in\R^n$ and
\[
 N(L,\rho):=\sum_{k\ge0}\bigl[c_{L+k}c_{L+k}^\top-c_{\rho+k}c_{\rho+k}^\top\bigr]\in M_n(\R),
\]
which converges absolutely (each term is $O(k^{-3})$). We write $\theta=L-1$, $s=\rho-L$ (so $\rho=\theta+s+1$), and
\[
 u_n(\theta,s):=\det\bigl[w,Nw,\dots,N^{n-1}w\bigr],\qquad N=N(\theta+1,\theta+s+1),\ w=c_{\theta+s+1}.
\]

\begin{lemma}\label{lem:N}
(a) $N(L,\rho)$ is real symmetric, meromorphic in $(L,\rho)$, with poles only where $L+k+a=0$ or $\rho+k+a=0$ for some $0\le a<n$, $k\ge0$; near any point only finitely many terms are singular.
(b) If $s=\rho-L\in\Z_{\ge0}$, then $N(L,\rho)=\sum_{y=L}^{\rho-1}c_yc_y^\top$ (a sum of $s$ terms); if $L-\rho=c\in\Z_{\ge0}$, then $N(L,\rho)=-\sum_{y=\rho}^{L-1}c_yc_y^\top$.
(c) $N(L,\rho+1)=N(L,\rho)+c_\rho c_\rho^\top$ and $N(L,\rho)=c_Lc_L^\top+N(L+1,\rho)$.
(d) Off-diagonal entries are rational functions of $(L,\rho)$, and $N(L,\rho)=N^{\rm rat}(L,\rho)+(\psi'(L)-\psi'(\rho))\,I$ with $N^{\rm rat}$ rational.
(e) $u_n$ is a rational function of $(\theta,s)$.
(f) For integer $s\ge0$ and $\theta=0$: $N(1,s+1)=N_s$ (the matrix of~\S\ref{ssec:kryred}, over $\Q$) and $w=c_{s+1}=w_{s+1}$; thus $u_n(0,s)=u(n,s)$ with $K=n$.
\end{lemma}
\begin{proof}
(a)--(c) are immediate from the definition (telescoping). For (d), $\sum_{k\ge0}\frac1{(a+y+k)(b+y+k)}=\frac1{b-a}\sum_{j=0}^{b-a-1}\frac1{a+y+j}$ for $a<b$, and $\sum_{k\ge0}(a+y+k)^{-2}=\psi'(y)-\sum_{j<a}(y+j)^{-2}$, i.e.\ $\psi'(y)$ minus a rational function. For (e), the Krylov determinant is invariant under $N\mapsto N+\tau I$ (the change of basis from $N^kw$ to $(N+\tau)^kw$ is unitriangular), so $\psi'(L)-\psi'(\rho)$ can be dropped, and $u_n$ is a polynomial in the entries of $N^{\rm rat}$ and $w$. (f) is (b).
\end{proof}

\begin{theorem}[Theorem A]\label{thm:A}
Put $T(m)=m(m+1)/2$ and
\begin{align*}
 A_n(s)&=\prod_{c=1}^{n-1}\bigl[(s+c)(s+1-c)\bigr]^{T(n-c)}, & B_n(\theta)&=\prod_{i=1}^{n-1}(\theta+i)^{2i(n-i)},\\
 C_n(t)&=\prod_{c=1}^{n}(t+c)^{\beta_c},\quad \beta_c=2c(n+1-c)-n, & \kappa_n^+&=\prod_{m=1}^{n-1}m^{2(n-m)(n-2m)}.
\end{align*}
Then there is a constant $\kappa_n\in\R$ with $|\kappa_n|=\kappa_n^+$ such that, as rational functions of $(\theta,s)$,
\[
 u_n(\theta,s)=\kappa_n\,\frac{A_n(s)}{B_n(\theta)\,C_n(\theta+s)} .
\]
\end{theorem}


\begin{remark}\label{rem:signA}
\emph{Only $|u_n|$ is needed.} Numerically $\kappa_n=(-1)^{\lfloor n/2\rfloor}\kappa_n^+$; the sign is not proved here and is not used. The formula was also verified exactly (rational arithmetic) for every $n\le12$, at random rational $(\theta,s)$ including non-integer and negative $s$, and with the Cauchy Gram matrices $N_s$ themselves (\S\ref{sec:formal}). For $n=3$, $\theta=0$ it gives $u(3,s)=-s^3(s-1)/\bigl(64(s+2)^4(s+3)^3\bigr)$.
\end{remark}

\subsection{Tools}

\begin{lemma}[resultant identity]\label{lem:res}
Let $N$ be real symmetric $n\times n$ with eigenvalues $\lambda_1,\dots,\lambda_n$, let $w\in\R^n$, and let $N'=N+ww^\top$ have eigenvalues $\mu_1,\dots,\mu_n$. Then
\[
 \det[w,Nw,\dots,N^{n-1}w]^2=\pm\prod_{i,k}(\lambda_i-\mu_k).
\]
\end{lemma}
\begin{proof}
Write $N=U\operatorname{diag}(\lambda)U^\top$ with $U$ orthogonal and $\tilde w=U^\top w$. Then $\det[w,Nw,\dots]=\det U\cdot\prod_k\tilde w_k\cdot\prod_{i<j}(\lambda_j-\lambda_i)$. On the other hand $\chi_{N'}(z)=\chi_N(z)\bigl(1-\sum_k\tilde w_k^2/(z-\lambda_k)\bigr)$, so $\chi_{N'}(\lambda_k)=-\tilde w_k^2\prod_{j\ne k}(\lambda_k-\lambda_j)$, and $\prod_{i,k}(\lambda_i-\mu_k)=\pm\prod_k\chi_{N'}(\lambda_k)=\pm\prod_k\tilde w_k^2\prod_{i<j}(\lambda_i-\lambda_j)^2$. (This is the classical identity for the resultant of the characteristic polynomials of $N$ and $N'$.)
\end{proof}
We call this identity (R). It exhibits $|u_n|^2$ as a product over pairs $(i,k)$ of the differences $|\lambda_i-\mu_k|$.

\begin{lemma}[cluster lemma]\label{lem:cluster}
Let $M(\varepsilon)=\varepsilon^{-2}\sigma J+\varepsilon^{-1}Y+Z(\varepsilon)$ be real symmetric $n\times n$ matrices, $\varepsilon\to0$, where $\sigma=\pm1$, $J=I_r\oplus0_{n-r}$, $Y$ is real symmetric with both diagonal blocks (the $r\times r$ and the $(n-r)\times(n-r)$ block) equal to zero, and $Z(\varepsilon)$ is bounded. Then $M(\varepsilon)$ has $r$ eigenvalues $\sigma\varepsilon^{-2}+O(1)$ and $n-r$ eigenvalues $O(1)$.
\end{lemma}
\begin{proof}
Write $\varepsilon^2M=\sigma J+\varepsilon Y+O(\varepsilon^2)$. Let $S=\begin{pmatrix}0&X\\-X^\top&0\end{pmatrix}$ with $X=-\sigma Y_{12}$ ($Y_{12}$ the off-diagonal block of $Y$), and $Q=\exp(\varepsilon S)$, which is orthogonal. Since $[J,S]=\begin{pmatrix}0&X\\X^\top&0\end{pmatrix}$,
\[
 Q^\top(\varepsilon^2M)Q=\sigma J+\varepsilon\bigl(Y+\sigma[J,S]\bigr)+O(\varepsilon^2)=\sigma J+O(\varepsilon^2),
\]
because the diagonal blocks of $Y$ vanish and the off-diagonal block of $Y+\sigma[J,S]$ is $Y_{12}+\sigma X=0$. Weyl's inequality $|\lambda_i(X_1+X_2)-\lambda_i(X_1)|\le\|X_2\|$ shows that $\varepsilon^2M$ has $r$ eigenvalues $\sigma+O(\varepsilon^2)$ and $n-r$ eigenvalues $O(\varepsilon^2)$.
\end{proof}

\subsection{Proof of Theorem~A}\label{ssec:proofA}

Throughout, $\theta\in\R\setminus\Z$ is fixed (generic) unless said otherwise, so that the points $s=-\theta-c$ (poles) and the integers $s=j$ (zeros) below are distinct. We write $\lambda_i$, $\mu_k$ for the eigenvalues of $N$ and $N'=N+ww^\top$ and use (R) throughout.

\medskip\noindent\textbf{Step 1: zeros in $s$.}
\emph{At $s=j$, $0\le j\le n-2$.} By Lemma~\ref{lem:N}(b), $N$ is a sum of $j$ rank-one terms, so $0$ is an eigenvalue of $N$ of multiplicity $\ge n-j$ and of $N'$ of multiplicity $\ge n-j-1$. Near $s=j$ the matrix $N$ is analytic in $s$ (for $\theta\notin\Z$ no term is singular), so by Weyl's inequality these eigenvalues are $O(|s-j|)$. This gives at least $(n-j)(n-j-1)$ factors of size $O(|s-j|)$ in (R), all other factors being bounded; hence $u_n$ vanishes at $s=j$ to order $\ge T(n-1-j)$.
\emph{At $s=-c$, $1\le c\le n-1$.} By Lemma~\ref{lem:N}(b), $N=-\sum_{y=\rho}^{L-1}c_yc_y^\top$ has rank $\le c$, and $w=c_\rho$ is one of these terms, so $N'$ has rank $\le c-1$. The same argument gives $(n-c)(n-c+1)$ small factors and order $\ge T(n-c)$.
Since $A_n(s)$ vanishes exactly at these points to exactly these orders ($s+1-c=0$ at $s=j=c-1$, with exponent $T(n-c)=T(n-1-j)$), $u_n/A_n$ has no pole in $s$ at the zeros of $A_n$.

\medskip\noindent\textbf{Step 2: poles in $s$.} For fixed $\theta$ the entries of $N$ and $w$ are singular only at $\rho=1-c$, i.e. $s=-\theta-c$, $c\ge1$. Near $\rho=1-c+\varepsilon$, the singular terms are $c_{\rho+k}c_{\rho+k}^\top$ with $k=c-1-a$ for $a$ in the set $\{0,\dots,\min(c,n)-1\}$ (call its size $c^*=\min(c,n)$); entry $a$ of $c_{\rho+k}$ has a simple pole, so each singular term contributes $-\varepsilon^{-2}e_ae_a^\top$ plus an $\varepsilon^{-1}$ term plus a bounded part. Thus
\[
 N=-\varepsilon^{-2}J_{c^*}+\varepsilon^{-1}Y+O(1).
\]
The $\varepsilon^{-1}$-coefficient $Y$ has vanishing diagonal blocks: for $p\ne q$ in the singular set, the contributions $1/(q-p)$ (from the term singular at $p$) and $1/(p-q)$ (from the term singular at $q$) cancel; the diagonal entries have no $\varepsilon^{-1}$ term since $(\varepsilon+m)^{-2}$ is regular for $m\ne0$; and entries with both indices outside the singular set are regular.
\emph{Case $c>n$.} Then $J_{c^*}=I$, $w$ is regular, and $N$, $N'$ are $-\varepsilon^{-2}I+\varepsilon^{-1}Y+O(1)$ with $Y=0$ (only the first block exists); all eigenvalues of both are $-\varepsilon^{-2}+O(1)$, all differences are $O(1)$, so $u_n$ is bounded: no pole.
\emph{Case $c\le n$.} Here $w=c_\rho$ has the pole in coordinate $a=c-1$, and $ww^\top$ is exactly the term $k=0$ of the $\rho$-series, so $N'=N+ww^\top$ has the same form with $J_{c-1}$ (and the same vanishing of the diagonal blocks of its $\varepsilon^{-1}$ coefficient). By Lemma~\ref{lem:cluster} (with $\sigma=-1$), $N$ has $c$ ``big'' eigenvalues $-\varepsilon^{-2}+O(1)$ and $n-c$ ``small'' ones $O(1)$, and $N'$ has $c-1$ big and $n-c+1$ small ones. In (R), big--big and small--small differences are $O(1)$, and each of the
\[
 c(n-c+1)+(c-1)(n-c)=2c(n+1-c)-n=\beta_c
\]
big--small differences has size $\varepsilon^{-2}$. Hence $u_n=O(\varepsilon^{-\beta_c})$.

\medskip\noindent\textbf{Step 3: poles in $\theta$ at $\theta=-i$, $1\le i\le n-1$, $s$ generic.} Put $L=1-i+\varepsilon$. The singular terms are $c_{L+k}c_{L+k}^\top$, $k=i-1-a$, $a\in\{0,\dots,i-1\}$, and exactly as in Step 2, but with sign $+$, $N=+\varepsilon^{-2}J_i+\varepsilon^{-1}Y+O(1)$ with $Y$ having vanishing diagonal blocks. Now $w=c_\rho$ is regular, so $N'$ has the same singular part. Both have $i$ big and $n-i$ small eigenvalues, giving $2i(n-i)$ big--small differences, so $u_n=O(\varepsilon^{-2i(n-i)})$.

\medskip\noindent\textbf{Step 4: identification in $s$.}
As $\rho\to\infty$ with $L$ fixed, $N=R_\infty(L)+O(1/\rho)$ with $R_\infty(L)=\sum_{k\ge0}c_{L+k}c_{L+k}^\top$, analytic in $1/\rho$, and $\rho w=\mathbf 1+O(1/\rho)$ with $\mathbf 1$ the all-ones vector. Hence $\rho^nu_n\to\det[\mathbf 1,R_\infty\mathbf 1,\dots,R_\infty^{n-1}\mathbf 1]$, finite, and $u_n=O(s^{-n})$.
Consider $F(s):=u_nC_n(\theta+s)/A_n(s)$ as a rational function of $s$. By Steps 1--2, $F$ has no poles (the poles of $u_n$ are at $s=-\theta-c$, $c\le n$, of order $\le\beta_c$, where $A_n\ne0$; the zeros of $A_n$ are integers, where $u_n$ vanishes to at least the order of $A_n$ and $C_n(\theta+s)\ne0,\infty$), so $F$ is a polynomial in $s$. Both $A_n$ and $C_n$ are monic and
\[
 \deg C_n-\deg A_n=\Bigl[\tfrac{n(n+1)(n+2)}3-n^2\Bigr]-\tfrac{(n-1)n(n+1)}3=n,
\]
so $F(s)=O(s^{-n}s^{n})=O(1)$, and $F$ is a constant $K_n(\theta)$:
\begin{equation}\label{eq:Ks}
 u_n(\theta,s)=K_n(\theta)\,\frac{A_n(s)}{C_n(\theta+s)}.
\end{equation}
Moreover $K_n(\theta)=\lim_{s\to\infty}s^nu_n=\det[\mathbf 1,R_\infty\mathbf 1,\dots]$ with $L=\theta+1$.

\medskip\noindent\textbf{Step 5: identification in $\theta$.}
By Lemma~\ref{lem:N}(d), $K_n$ is a rational function of $\theta$ (the scalar $\psi'(L)$ in $R_\infty$ drops out of the Krylov determinant), and its entries $\sum_{k\ge0}c_{L+k}c_{L+k}^\top-\psi'(L)I$ have poles only where $L+j=0$ for some $0\le j\le n-2$, that is $\theta=-i$, $1\le i\le n-1$. Comparing \eqref{eq:Ks} with Step~3 (at $\theta=-i$ and generic $s$, $A_n(s)/C_n(\theta+s)$ is finite and nonzero), the pole of $K_n$ at $-i$ has order $\le2i(n-i)$. Hence $K_nB_n$ is a polynomial in $\theta$.
Now let $L\to\infty$ with $\rho$ fixed and generic. Then $N(L,\rho)\to-\sum_{k\ge0}c_{\rho+k}c_{\rho+k}^\top$ (the $L$-series is $O(1/L)$) and $w=c_\rho$ is fixed, so $u_n$ stays bounded. In \eqref{eq:Ks}, $s=\rho-L$ and $C_n(\theta+s)=C_n(\rho-1)$ is a fixed nonzero number, while $A_n(\rho-L)$ is of exact order $L^{\deg A_n}$. Hence $K_n(\theta)=O(\theta^{-\deg A_n})$. Since $\deg B_n=\sum_i2i(n-i)=\frac{(n-1)n(n+1)}3=\deg A_n$, the polynomial $K_nB_n$ is bounded, hence a constant $\kappa_n$. This proves the shape of Theorem~A:
\begin{equation}\label{eq:shapeA}
 u_n=\kappa_n\frac{A_n(s)}{B_n(\theta)\,C_n(\theta+s)}.
\end{equation}

\medskip\noindent\textbf{Step 6: the constant, by recursion in $n$.}

\begin{lemma}[recursion]\label{lem:rec}
Fix non-integer $s$. As $\theta=-1+\varepsilon\to-1$,
\[
 \varepsilon^{2(n-1)}\,|u_n(-1+\varepsilon,s)|\ \longrightarrow\ \frac1{|s|}\,\bigl|u_{n-1}(1,\,s-1)\bigr|\qquad(n\ge2),
\]
where on the right $u_{n-1}$ is the $(n-1)$-dimensional function evaluated at $\theta'=1$, $s'=s-1$.
\end{lemma}
\begin{proof}
Here $L=\varepsilon$ and $\rho=s+\varepsilon$. We have $c_L=\varepsilon^{-1}p$ with $p=e_0+O(\varepsilon)$, and by Lemma~\ref{lem:N}(c) $N=\varepsilon^{-2}pp^\top+A(\varepsilon)$ with $A(\varepsilon)=N(1+\varepsilon,\rho)$ analytic at $\varepsilon=0$. Put $N'=N+ww^\top=\varepsilon^{-2}pp^\top+A'(\varepsilon)$, $A'=A+ww^\top$; $w=c_\rho$ has first coordinate $w_0=1/\rho\to1/s$.
\emph{Big eigenvalues.} The largest eigenvalues are $\lambda_{\rm big}=\varepsilon^{-2}|p|^2+\hat p^\top A\hat p+O(\varepsilon^2)$ and $\mu_{\rm big}=\varepsilon^{-2}|p|^2+\hat p^\top A'\hat p+O(\varepsilon^2)$ with $\hat p=p/|p|\to e_0$; so $\mu_{\rm big}-\lambda_{\rm big}\to w_0^2=1/s^2$.
\emph{Small eigenvalues.} $\varepsilon^2\chi_N(z)=\varepsilon^2\det(z-A)-p^\top\operatorname{adj}(z-A)p\to-\det\bigl(z-A(0)[1{:}]\bigr)$, where $A(0)[1{:}]$ is the compression of $A(0)=N(1,s)$ to the coordinates $1,\dots,n-1$; likewise for $N'$ with $A'(0)[1{:}]=A(0)[1{:}]+\tilde w\tilde w^\top$, $\tilde w=(w_a)_{a\ge1}$. Hence the other $n-1$ eigenvalues of $N$ (resp.\ $N'$) converge to those of $A(0)[1{:}]$ (resp.\ $A'(0)[1{:}]$). Re-indexing $a'=a-1$, the entries of $A(0)[1{:}]$ are $\sum_k\bigl[\frac1{(a'+2+k)(b'+2+k)}-\frac1{(a'+s+1+k)(b'+s+1+k)}\bigr]=N(2,s+1)_{a'b'}$ and $\tilde w=c_{s+1}$, that is, the $(n-1)$-dimensional problem with $L'=2$, $\rho'=s+1$, i.e.\ $\theta'=1$, $s'=s-1$.
Now apply (R). The big--big difference tends to $1/s^2$; each of the $2(n-1)$ big--small differences is $\pm\varepsilon^{-2}(1+o(1))$; the small--small differences tend to the corresponding product for the $(n-1)$-dimensional problem, which is $\pm u_{n-1}(1,s-1)^2$ by (R). Thus $|u_n|^2\sim s^{-2}\varepsilon^{-4(n-1)}u_{n-1}(1,s-1)^2$.
\end{proof}

\begin{lemma}[exponent bookkeeping]\label{lem:expo}
$C_n(s-1)\,A_{n-1}(s-1)=s\,C_{n-1}(s)\,A_n(s)$ as polynomials in $s$.
\end{lemma}
\begin{proof}
Compare the exponent of $(s+d)$ on both sides ($T(m)=0$ for $m\le0$):
for $1\le d\le n-3$: $2(d+1)(n-d)-n+T(n-2-d)=2d(n-d)-n+1+T(n-d)$;
$d=n-2$: both sides equal $3n-4$; $d=n-1$: both equal $n$; $d=0$: $n+T(n-2)=1+T(n-1)$; $d<0$ (the factors $s+1-c$): both equal $T(n-1+d)$.
\end{proof}
(The identity of Lemma~\ref{lem:expo} was also checked by exact computation for $2\le n\le8$.)

\begin{proof}[Proof of Theorem~A, conclusion]
Evaluate \eqref{eq:shapeA} at $\theta=-1+\varepsilon$. Since $B_n(-1+\varepsilon)=\varepsilon^{2(n-1)}\prod_{i=2}^{n-1}(i-1+\varepsilon)^{2i(n-i)}$,
\[
 \varepsilon^{2(n-1)}|u_n|\to\frac{|\kappa_n|\,|A_n(s)|}{\prod_{i=2}^{n-1}(i-1)^{2i(n-i)}\,|C_n(s-1)|}.
\]
By Lemma~\ref{lem:rec} and \eqref{eq:shapeA} for $n-1$ at $(\theta',s')=(1,s-1)$, the limit is also $\dfrac{|\kappa_{n-1}|\,|A_{n-1}(s-1)|}{|s|\,B_{n-1}(1)\,|C_{n-1}(s)|}$, with $B_{n-1}(1)=\prod_{i=1}^{n-2}(i+1)^{2i(n-1-i)}$. By Lemma~\ref{lem:expo}, $A_{n-1}(s-1)/(s\,C_{n-1}(s))=A_n(s)/C_n(s-1)$, so all $s$-dependence cancels and
\[
 |\kappa_n|=|\kappa_{n-1}|\cdot\frac{\prod_{i=2}^{n-1}(i-1)^{2i(n-i)}}{\prod_{i=1}^{n-2}(i+1)^{2i(n-1-i)}}=|\kappa_{n-1}|\prod_{m=1}^{n-1}m^{2(2n-1-3m)},
\]
where the exponent of $m$ is $2[(m+1)(n-m-1)-(m-1)(n-m)]=2(2n-1-3m)$. Since $u_1=1/(\theta+s+1)$, $\kappa_1=1$, and summing the exponents of $m$ over $N=m+1,\dots,n$ gives $2\sum_{N=m+1}^{n}(2N-1-3m)=2(n-m)(n-2m)$. Hence $|\kappa_n|=\prod_{m=1}^{n-1}m^{2(n-m)(n-2m)}=\kappa_n^+$. \qed
\end{proof}

\subsection{Proofs of Lemma CC and Theorem~\ref{thm:nonvanish2}}\label{ssec:CCproof}

\begin{proof}[Proof of Lemma~CC]
Let $n=K$ and $s\ge K$, $K+s<\ell$. By Lemma~\ref{lem:N}(f) and Theorem~A at $\theta=0$,
\[
 u(K,s)=\kappa_K\,\frac{A_K(s)}{B_K(0)\,C_K(s)}\quad\text{in }\Q,
\]
with $\kappa_K$ of absolute value $\prod_{m<K}m^{2(K-m)(K-2m)}$ (exponents may be negative). The bases occurring are: the integers $m\in[1,K-1]$ (from $\kappa_K$ and $B_K(0)=\prod i^{2i(K-i)}$); $s+c$ with $1\le c\le K$, so $s+c\in[s+1,s+K]$ (from $A_K$ and $C_K$); and $s+1-c\in[s-K+2,s]\subset[2,s]$ (from $A_K$). All are integers in $[1,K+s]$, and $K+s<\ell$, so every base is an $\ell$-unit and $u(K,s)$ is a nonzero element of $\Z_{(\ell)}$ that is a unit. The entries of $N_s$ and $w_{s+1}$ are $\ell$-integral, so the reduction modulo~$\ell$ of $u(K,s)$ is the Krylov determinant over $\mathbb F_\ell$ of Corollary~\ref{cor:kr}, which is therefore nonzero. Corollary~\ref{cor:kr} gives Lemma~CC. \qed
\end{proof}
Note that only $|u|$ enters: we use only that $u(K,s)$ is a nonzero rational number all of whose prime factors are $<\ell$.

\begin{proof}[Proof of Theorem~\ref{thm:nonvanish2}]
For $h_1=\ell-K$ and $h_2=\ell-K-1$ we have $s=\ell-h=K$ and $K+1$ respectively. Both satisfy $K+1\le h_i\le\ell-1$ ($h_2\ge K+1\iff\ell\ge5$) and $1\le s\le\ell-K-1$ ($K+1\le\ell-K-1\iff\ell\ge5$), so Lemmas~\ref{lem:1} and~\ref{lem:2} apply: $\ell^{2(h-K)}P_h(r)$ is an $\ell$-adic unit iff $D_s(Y):=\det(YI_K+N_s)$ does not vanish at $Y=y:=16r+c\in\mathbb F_\ell$. Suppose $D_K(y)=D_{K+1}(y)=0$. As in the proof of Theorem~\ref{thm:nonvanish}, $D_K(0)=\det(C_K)^2\ne0$ (Cauchy matrix with entries of size $<\ell$), so $y\ne0$; thus $-y\ne0$ is a common eigenvalue of $N_K$ and $N_{K+1}$, contradicting Lemma~CC with $s=K$ ($2K<\ell$). Hence for some $i$ the reduction is nonzero and $v_\ell(P_{h_i}(r))=-2(h_i-K)$. \qed
\end{proof}

\begin{remark}[numerical checks]
Independently of Theorem~A, the nonvanishing of $u(K,K)$ modulo $\ell$ (and hence Lemma~CC for $s=K$) was verified by exact computation of the Krylov determinant for all $\ell\equiv3\pmod4$ with $7\le\ell<4500$ (307 primes); for $39$ primes $\ell<400$ the characteristic polynomials of $N_K$ and $N_{K+1}$ over $\mathbb F_\ell$ were checked to be coprime, and Lemma~\ref{lem:2} was checked for $\ell=7,11,19,23$ and $h=\ell-K,\ell-K-1$ and all $X\in\mathbb F_\ell$. These are consistency checks, not used in the proof.
\end{remark}

%%% --------------------------------------------------------------------
\section{Node removal and the root-free window}\label{sec:nodeline}

\subsection{Why a small value is not enough}
Let $Q_h=\den(P_h)P_h$ and $|\lambda-p/q|=|t|$. The mean-value route, $|Q_h(p/q)-Q_h(\lambda)|\le|t|\max|Q_h'|$, with $\max|Q_h'|$ of size $e^{\Lm h^2+O(h)}$, only yields a bound of the shape $\exp(-C\log^2q)$, which is not a finite irrationality measure. What is needed is smallness of $|P_h|$ on a whole window $|x-\lambda|\le e^{-ch}$ of \emph{exponential radius}, together with $Q_h(p/q)\ne0$. The second requirement is supplied by Theorem~\ref{thm:nonvanish} (valid at every rational, including $p/q$). For the first, note that
\begin{equation}\label{eq:expand}
 P_h(\lambda+t)=\det(T_h+tI)=\sum_{S\subseteq\{0,\dots,h-1\}}t^{\,h-|S|}\det T_h[S],
\end{equation}
by multilinearity in the columns, where $T_h[S]$ is the principal submatrix on $S$. So the window requires a bound for \emph{every} principal minor, not only for the full determinant, and uniform in the number $j=h-|S|$ of removed nodes. The proof of Theorem~\ref{thm:decay} extends to principal minors; the only new feature is a cost for each removed node.

\subsection{Statement}
Retain the notation of \S\ref{ssec:cert}: the certificate measure $\sigma$ with $n_\sigma=300$ point masses, $\rho=1/1024$, $g=U^\sigma_\rho-U^\tau+(1-\kappa)\pi|\Im s|$, and the certified data $m'\ge1.9718000062$, $S_{\rm cert}\ge-0.1107505186$, $g_{\min}\ge-1.4640996549$ and $D_{\rm cert}\ge3.1113790$. On the real line the field term vanishes, so $g=U^\sigma_\rho-U^\tau$ there, and $g_{\min}$ bounds it below on $[0,1]$. Recall $c=2(m'-g_{\min})\le c_{\rm cert}:=6.8717994$.

\begin{theorem}[node removal $N_\rho$; computer-assisted]\label{thm:N}
There are $\eta(h)=O((\log h)/h)$ and a constant $c_\kappa$, both independent of $S$, such that for every $h\ge2^{30}$ and every $S\subseteq\{0,\dots,h-1\}$, with $k=|S|$, $j=h-k$ and $\beta=j/h$,
\[
 |\det T_h[S]|\ \le\ 16^{-k}\,(2h+c_\kappa)^k\,\exp\bigl(-h^2[D_{\rm cert}-\beta c]+h^2\eta(h)\bigr).
\]
Explicitly, $\eta(h)=\frac{|\log\rho|}h+\frac{2\log(1+8h/3)}h+\varepsilon_2(h)+\frac{\log^+C_E}h+\frac4h\sum_\ell(m'-\gamma_\ell)^+$.
\end{theorem}
For $S=\{0,\dots,h-1\}$ ($\beta=0$) this is Theorem~\ref{thm:decay}. The proof is that of Lemmas~\ref{lem:A} and~\ref{lem:B} line by line; we record what changes. (The coefficient of the last term of $\eta$ is $4/h$ because $H_S$ carries the factor $2/h$ in front of $\sum g$, so the error $\frac2h\sum_\ell(m'-\gamma_\ell)^+$ of Lemma~\ref{lem:B} is doubled.)

\subsection{Proof of Theorem~\ref{thm:N}}\label{ssec:proofN}
\textbf{Inputs.} By Proposition~\ref{prop:contour} the identity $L_F[a][b]=\int v_av_b\,dm$ is entrywise, so it holds for principal submatrices: $L_F[S]=\int\phi_S\phi_S^{\top}dm$ with $\phi_S=(v_a)_{a\in S}$, and $T_h[S]=L_F[S]/16$. As $m$ has finite total variation and $|v_a|\le8/3$ on its support, Andr\'eief's identity gives $\det L_F[S]=\frac1{k!}\int\det[v_a(z_i)]_{i\le k,a\in S}^2\prod dm(z_i)$ (for $k=0$ both sides are $1$), and therefore
\[\begin{gathered}
 |\det L_F[S]|\le\frac1{k!}\int|C_S(z)|^2\prod_{i\le k}d|m|(z_i),\\
 |C_S|^2=\frac{\prod_{i\ne i'}|z_i-z_{i'}|\prod_{a\ne b\in S}|a-b|}{\prod_{i,\,a\in S}|z_i-a|^2},
\end{gathered}\]
by Cauchy's determinant. Split $d|m|=w_\kappa\,d\Omega_\kappa$ as in the proof of Theorem~\ref{thm:decay}; then $\Omega_\kappa(\Sigma_h)\le2h+c_\kappa$ and $c_\kappa$ is independent of $h$. Every point of the support of $m$ has $\Re z\le-\frac38$ while the nodes lie in $[0,h-1]$, so free points and nodes never coincide. It therefore suffices to prove the pointwise bound
\begin{equation}\label{eq:pointN}
 |C_S(z)|^2\prod_iw_\kappa(z_i)\le\exp\bigl(-h^2[D_{\rm cert}-\beta c-\eta(h)]\bigr),
\end{equation}
since integrating it against $\prod\Omega_\kappa$ and using $1/k!\le1$ gives the theorem.

\textbf{Scaling.} Put $s_i=z_i/h$, $t_a=a/h$. Counting powers of $h$: $-k(k-1)$ from the free pairs, $-k(k-1)$ from the node pairs inside $S$ and $+2k^2$ from the cross terms, in total $2k\log h\ge0$ in the exponent of the reciprocal, which is dropped. Hence $-\log(|C_S|^2\prod w_\kappa)\ge h^2H_S$ with
\[\begin{aligned}
 h^2H_S={}&\sum_{i\ne i'}-\log|s_i-s_{i'}|+\sum_{a\ne b\in S}-\log|t_a-t_b|\\
 &+2\!\!\sum_{i,\,b\in S}\!\log|s_i-t_b|+\sum_i-\log w_\kappa(z_i).
\end{aligned}\]

\textbf{Completing the node set.} Let $S^c$ be the $j$ removed nodes, and let all pairs be of distinct elements. Two exact identities, from adding and removing the pairs meeting $S^c$ (for the second: the ordered pairs meeting $S^c$ are those with first entry in $S^c$, plus those with second entry in $S^c$, minus those with both; the summand is symmetric), are
\[
 2\!\!\sum_{i,\,b\in S}\!\log|s_i-t_b|=2\!\sum_{i,\,\text{all }b}\!\log|s_i-t_b|+2\!\!\sum_{i,\,a\in S^c}\!(-\log|s_i-t_a|),
\]
\[
 \sum_{b\ne b'\in S}=\sum_{\text{all }b\ne b'}-2\sum_{a\in S^c}\sum_{b\ne a}+\sum_{a\ne a'\in S^c}\qquad\text{(applied to }-\log|t_b-t_{b'}|).
\]
Let $X$ be the set of the $h$ distinct points $s_1,\dots,s_k$ and $t_a$ ($a\in S^c$); it has $k+j=h$ elements. The free--free, free--removed and removed--removed terms assemble to $\sum_{x\ne y\in X}-\log|x-y|$, and
\[\begin{aligned}
 h^2H_S={}&\underbrace{\sum_{x\ne y\in X}-\log|x-y|}_{[B]}+\underbrace{2\!\sum_{i,\,\text{all }b}\!\log|s_i-t_b|}_{[R1]}\\
 &+\underbrace{\sum_{\text{all }b\ne b'}\!-\log|t_b-t_{b'}|}_{[R2]}\\
 &+\underbrace{2\!\sum_{a\in S^c}\sum_{b\ne a}\!\log|t_a-t_b|}_{[R3]}+\underbrace{\sum_i-\log w_\kappa(z_i)}_{[R4]}.
\end{aligned}\]
For $S=\{0,\dots,h-1\}$ the term $[R3]$ is absent and this is the decomposition in the proof of Lemma~\ref{lem:A}.

\textbf{The term $[B]$ (replaces circle smoothing).} The pairs of removed nodes occur in $[B]$ and once more, with the opposite sign, in the second identity above; so \emph{the interaction of the removed nodes among themselves is counted exactly once and no energy of the removed set is left over}. Apply Lemma~\ref{lem:K0} to every ordered pair of distinct points of $X$, free or removed, and then Corollary~\ref{cor:Kp} with this $X$:
\begin{multline*}
 [B]=\sum_{x\ne y\in X}-\log|x-y|\ \ge\ \sum_{x\ne y\in X}k_\rho(x,y)\\
 \ge\ h^2\Bigl[\tfrac2h\textstyle\sum_{x\in X}U^\sigma_\rho(x)-I_\rho(\sigma)\Bigr]+h\log\rho,
\end{multline*}
that is, $\frac{[B]}{h^2}\ge\frac2h\sum_{i\le k}U^\sigma_\rho(s_i)+\frac2h\sum_{a\in S^c}U^\sigma_\rho(t_a)-I_\rho(\sigma)-\frac{|\log\rho|}h$. There is no Lipschitz term, no smoothing and no continuum energy: only the finite-sum Lemma~\ref{lem:K}.

\textbf{The terms $[R1]$, $[R2]$, $[R4]$.} As in Lemma~\ref{lem:A}: $[R1]/h^2\ge-\frac2h\sum_iU^\tau(s_i)-\frac{2k}{h^2}\log(1+\frac{8h}3)$ (monotonicity of $u\mapsto\log|s-u|$ on $[0,1]$ for $\Re s\le0$, and $|s_i|\ge\frac3{8h}$), and $[R4]/h^2\ge\frac2h\sum_i(1-\kappa)\pi|\Im s_i|-\frac k{h^2}\log^+C_E$. The term $[R2]=h(h-1)\log h-2\sum_{n<h}\log n!$ involves all $h$ nodes and is independent of $S$. Using $\log n!\le(n+\frac12)\log n-n+1$ for $n\ge1$ gives $[R2]/h^2\ge\frac32-\varepsilon_2(h)$ with $\varepsilon_2(h)=O((\log h)/h)$. Nodes carry no weight, so no field term arises at the removed nodes.

\textbf{The term $[R3]$: the removed-node term.} For a node $a$ with $t=a/h$,
\[
 \sum_{b\ne a}\log|t_a-t_b|=\log\bigl(a!\,(h-1-a)!\bigr)-(h-1)\log h .
\]
By $n!\ge(n/e)^n$, valid for all $n\ge0$, applied to $a!$ and to $(h-1-a)!=(h-a)!/(h-a)$, we get $\log(a!(h-1-a)!)\ge h[t\log t+(1-t)\log(1-t)]+h\log h-h-\log h$, that is, $\sum_{b\ne a}\log|t_a-t_b|\ge-h\,U^\tau(t_a)$ \emph{with no error term}. Hence $[R3]/h^2\ge-\frac2h\sum_{a\in S^c}U^\tau(t_a)$.

\textbf{Collecting.} On the real line $U^\sigma_\rho-U^\tau=g$, so
\begin{equation}\label{eq:HS}
 H_S\ \ge\ \frac2h\sum_{i\le k}g(s_i)+\frac2h\sum_{a\in S^c}g(t_a)-I_\rho(\sigma)+\frac32-\eta_1(h),
\end{equation}
with $\eta_1(h)=\frac{|\log\rho|}h+\frac{2\log(1+8h/3)}h+\varepsilon_2(h)+\frac{\log^+C_E}h$, which uses only $k\le h$ and so is uniform in $S$.

\emph{Removed nodes.} Since $t_a\in[0,1)$, (L) gives $\frac2h\sum_{a\in S^c}g(t_a)\ge2\beta\,g_{\min}$. As $g_{\min}<0$, a removed node costs.

\emph{Free points (Lemma~\ref{lem:B}).} The free points have mass $k/h=1-\beta$, and Lemma~\ref{lem:B} gives
\[
 \frac1h\sum_{i\le k}g(s_i)\ \ge\ (1-\beta)m'+S_{\rm cert}-\frac2h\sum_\ell(m'-\gamma_\ell)^+ .
\]
Inserting both into \eqref{eq:HS},
\[
 \begin{aligned}H_S&\ge2(1-\beta)m'+2S_{\rm cert}+2\beta g_{\min}-I_\rho(\sigma)+\tfrac32-\eta_1-\tfrac4h\textstyle\sum_\ell(m'-\gamma_\ell)^+\\
 &=D_{\rm cert}-\beta\cdot2(m'-g_{\min})-\eta(h),\end{aligned}
\]
where (since $D_{\rm cert}=2(m'+S_{\rm cert})-I_\rho(\sigma)+\frac32$, with $S_{\rm cert}$ the sum of Lemma~\ref{lem:B}) the expression in the statement is exactly $D_{\rm cert}-\beta c$, and $\eta=\eta_1+\frac4h\sum_\ell(m'-\gamma_\ell)^+$ is independent of $S$. This is \eqref{eq:pointN}. \qed

\subsection{The certificate for $g_{\min}$}\label{ssec:cert2}
The number $g_{\min}$ is part of the certificate of \S\ref{ssec:cert}, produced by the same script \texttt{kcert.py} on the same \texttt{sigma300.json}. The interval $[-\frac1{64},1]\supset[0,1]$ is divided into $4423$ exact cells and each cell becomes one Arb enclosure \cite{Arb} (full-cell enclosures, so no Lipschitz bound is needed). On each cell the lower bound for $U^\sigma_\rho$ is the one described in \S\ref{ssec:cert}, and $U^\tau(t)=1-t\log|t|-(1-t)\log|1-t|$ has $(U^\tau)'=\log|1-t|-\log|t|$, which is positive for $t<\frac12$ and negative for $t>\frac12$ (also across $0$), so its supremum on a cell $[a,b]$ is $U^\tau(b)$ if $b\le\frac12$, $U^\tau(a)$ if $a\ge\frac12$, and $U^\tau(\frac12)=1+\log2$ otherwise. The certified infimum is
\[
 g_{\min}\ \ge\ -1.4640996549,
\]
with the minimiser near $t\approx0.7499$ in floating point. Hence $c=2(m'-g_{\min})\le2(1.9718000062+1.4640996549)=6.8717993\ldots\le c_{\rm cert}:=6.8717994$.

\subsection{The window corollary}
\begin{corollary}\label{cor:window}
For $h\ge2^{30}$ and every real $t$ with $|t|\le e^{-ch}$,
\[
 |P_h(\lambda+t)|\ \le\ \exp\bigl(-D_{\rm cert}h^2+h^2\eta(h)+h\log\bigl(2(2h+c_\kappa)\bigr)\bigr)=e^{-(D_{\rm cert}-o(1))h^2},
\]
where the $o(1)$ does not depend on $t$. In particular this holds for $|t|\le e^{-c_{\rm cert}h}$.
\end{corollary}
\begin{proof}
By \eqref{eq:expand}, grouping the subsets $S$ by $j=h-|S|$ (there are $\binom hj$ of them) and applying Theorem~\ref{thm:N}, with $\beta=j/h$ so that $e^{h^2\beta c}=e^{jch}$,
\[
 |P_h(\lambda+t)|\le(2h+c_\kappa)^he^{h^2\eta}e^{-D_{\rm cert}h^2}\sum_{j=0}^h\binom hj16^{-(h-j)}\bigl(|t|e^{ch}\bigr)^j .
\]
(Here $(2h+c_\kappa)^{k}\le(2h+c_\kappa)^h$ for $2h+c_\kappa\ge1$.) If $|t|e^{ch}\le1$ every summand is at most $\binom hj$, and the sum is at most $2^h$. \end{proof}

\begin{remark}[earlier versions]
An earlier version of this paper used a Gram-eigenvalue bound, giving $\mu(\lambda)\le126$. The node-removal bound replaced it, first with the circle-smoothing energy step (Remark~\ref{rem:circle}, $\mu(\lambda)\le79.35$) and now with the $k_\rho$ energy step.
\end{remark}

%%% --------------------------------------------------------------------
\section{Proof of the main theorem}\label{sec:measure}

\begin{theorem}\label{thm:measure}
Let $\eta=D_{\rm cert}-\Lm\ge0.1304266$ and $c_{\rm cert}=6.8717994$. Then for every $\nu>c_{\rm cert}\,(1/\eta+3/4)$ there are only finitely many rationals $p/q$ with $|\lambda-p/q|<q^{-\nu}$. In particular, since $c_{\rm cert}(1/\eta+3/4)\le57.841$ (it is $57.8410$ for the unrounded values) and $57.841<57.85$, Theorem~\ref{thm:main} holds.
\end{theorem}

\begin{proof}
Let $\nu>c_{\rm cert}(\frac1\eta+\frac34)$ and suppose $|\lambda-p/q|=:|t|<q^{-\nu}$ with $q$ large; put $L=\log q$. Choose $\delta>0$ so small that still $\nu>c_{\rm cert}\,(1+\delta)\bigl(\frac1\eta+\frac34\bigr)$, and set
\[
 k=\frac{1+\delta}\eta,\qquad X=\tfrac43kL,\qquad Y=X+(1+\delta)L.
\]
\emph{Choice of $\ell$.} By the prime number theorem in arithmetic progressions mod~$4$ (de la Vall\'ee Poussin; \cite[Ch.~20]{Davenport}, \cite[Cor.~11.19]{MV}), $\vartheta(x;4,3)=x/2+o(x)$, so the primes $\ell\equiv3\pmod4$ in $(X,Y]$ have product $\exp\bigl(\tfrac12(Y-X)(1+o(1))\bigr)>q^{1/2}$ for large $q$. The primes $\ell$ with $v_\ell(q)=2$ have product at most $q^{1/2}$, since their squares divide $q$. Hence there is a prime $\ell\equiv3\pmod4$ in $(X,Y]$ (so $\ell\ge7$) with $v_\ell(q)\ne2$.

\emph{Choice of $h$.} Put $K=(\ell+1)/4$. If $\ell\nmid q$, Theorem~\ref{thm:nonvanish2} gives $h\in\{\ell-K,\ell-K-1\}=\{\frac{3\ell-1}4,\frac{3\ell-5}4\}$ with $P_h(p/q)\ne0$. If $\ell\mid q$, then $v_\ell(q)=1$ or $\ge3$ and Lemma~\ref{lem:3} gives $P_h(p/q)\ne0$ for every admissible $h$, in particular for $h=\ell-K$. In every case
\[
 \tfrac{3\ell-5}4\le h\le\tfrac{3\ell-1}4,\qquad\text{so}\qquad kL-\tfrac54< h<\tfrac34Y
\]
(because $\ell>X=\frac43kL$ and $\ell\le Y$), $h\ge2^{30}$ for large $q$, and $Q_h(p/q)\ne0$, whence by Lemma~\ref{lem:integrality}
\[
 |Q_h(p/q)|\ge q^{-h}=e^{-hL}.
\]
\emph{Window.} Here $\frac34Y=kL+\frac34(1+\delta)L=(1+\delta)\bigl(\frac1\eta+\frac34\bigr)L$, so $h<\frac34Y$ and the choice of $\delta$ give $c_{\rm cert}\,h<\nu L$. Hence $|t|<e^{-\nu L}\le e^{-c_{\rm cert}h}\le e^{-ch}$, and Corollary~\ref{cor:window} and Theorem~\ref{thm:content} give
\[
 |Q_h(p/q)|\le\den(P_h)\,|P_h(\lambda+t)|\le\exp\bigl((\Lm-D_{\rm cert}+o(1))h^2\bigr)=\exp\bigl(-(\eta-o(1))h^2\bigr).
\]
\emph{Contradiction.} Combining, $(\eta-o(1))h\le L$. But $h>kL-\frac54$ gives $\eta h\ge(1+\delta)L-\frac54\eta$, a contradiction for large $q$. \end{proof}

\begin{remark}[sensitivity]
The exponent is $c(\frac1\eta+\frac34)$ with $\eta=D_{\rm cert}-\Lm$ and $c=2(m'-g_{\min})$. The term $c/\eta=52.687$ is what remains if the interval for $\ell$ has length $0$; the term $\frac34c$ comes from the length $(1+\delta)L$ of that interval, which is needed to find a prime $\ell\equiv3\pmod4$ with $v_\ell(q)\ne2$, multiplied by $h/\ell\approx\frac34$. Lemma~\ref{lem:3} is what lets one use $\ell\mid q$ when $v_\ell(q)\ne2$; without it one must avoid every $\ell\mid q$, the interval for $\ell$ must have length $(2+\delta)L$, the term $\frac34$ becomes $\frac32$, and the bound is $c(\frac1\eta+\frac32)=62.995$ (see the table of \S\ref{ssec:cert}).
\end{remark}

\begin{remark}[a simpler choice of $h$ gives $77.13$]\label{rem:77}
If in the proof of Theorem~\ref{thm:measure} one uses Theorem~\ref{thm:nonvanish} instead of Theorem~\ref{thm:nonvanish2}, that is $h\in\{\ell-1,\ell-2,\frac{3\ell-1}4\}$ (and $h=\ell-1$ when $\ell\mid q$), only $\frac{3\ell-1}4\le h\le\ell-1$ is available. With $X=\frac43kL$, $Y=X+(1+\delta)L$ as above one gets $kL-\frac14\le h\le Y$, and the window condition $c_{\rm cert}\,Y<\nu L$ requires $\nu>c_{\rm cert}\bigl(\frac{4(1+\delta)}{3\eta}+1+\delta\bigr)$. The rest of the proof is unchanged and gives: for every $\nu>c_{\rm cert}\,(4/(3\eta)+1)$ there are only finitely many $p/q$ with $|\lambda-p/q|<q^{-\nu}$. Since $4/(3\eta)+1\le11.2229$ ($11.22286\ldots$ unrounded) and $c_{\rm cert}\cdot11.2229\le77.1217<77.13$ (unrounded $77.1212$), this gives $\mu(\lambda)\le77.13$. The factor $\frac43$ comes from the spread of the three values of $h$, and without Lemma~\ref{lem:3} the bound is $c(4/(3\eta)+2)<84$ ($83.99$). The Lean formalization currently covers this version (\S\ref{sec:formal}), and it does not use Theorem~A.
\end{remark}

%%% --------------------------------------------------------------------
\section{The mirror constant}\label{sec:mirror}

Put
\[
 \lambda'\;:=\;\sum_{k\ge1}\Bigl(\frac1{(4k-3)^2}-\frac1{16k^2}\Bigr)\;=\;\Bigl(\frac{\pi^2}{16}+\frac G2\Bigr)-\frac{\pi^2}{96}\;=\;\frac G2+\frac{5\pi^2}{96}\;=\;0.972\ldots,
\]
using $\sum_{n\equiv1(4)}n^{-2}=\frac{\pi^2}{16}+\frac G2$ (the companion of the sums in \S\ref{sec:matrix}); thus $G+5\pi^2/48=2\lambda'$ and $\lambda+\lambda'=G$.
The proof of Theorem~\ref{thm:mirror} is the proof of Theorems~\ref{thm:irr} and~\ref{thm:main} for the \emph{mirror frame}
\[
 \bar F(z)=\psi(z+1)-\psi(z+\tfrac14)=\sum_{j\ge0}\Bigl[\frac1{z+j+\frac14}-\frac1{z+j+1}\Bigr],
\]
with poles $+1$ at $-(j+\frac14)$ and $-1$ at $-(j+1)$. The shift $\frac34\to\frac14$ exchanges the roles of the residue classes $1$ and $3$ modulo~$4$.
We list \emph{every} place where the argument changes; everything not listed is used verbatim. The constants $\Lm=313/105$, $D_{\rm cert}\ge3.1113790$, $c_{\rm cert}=6.8717994$, $\eta\ge0.1304266$ are the same, so the numerical value of the exponent is the same, $57.841\le57.85$ (and $77.1213$ for the simpler choice of $h$).

\subsection{The matrices}
From $\bar F(k)-\bar F(k-1)=\frac1k-\frac4{4k-3}=-\frac3{k(4k-3)}$ and the analogous identity for $\bar F'$ one gets, exactly as for \eqref{eq:Loff}, \eqref{eq:Ldiag},
\[
 \bar L[a][b]=-\frac3{b-a}\sum_{a<k\le b}\frac1{k(4k-3)}\ (a<b),\qquad
 \bar L[a][a]=\bar F'(0)+\sum_{k\le a}\Bigl(\frac{16}{(4k-3)^2}-\frac1{k^2}\Bigr),
\]
and $\bar F'(0)=\psi'(1)-\psi'(\tfrac14)=\frac{\pi^2}6-(\pi^2+8G)=-16\lambda'$. Let $N_0=N_0^{(h)}$ be the symmetric rational matrix with
\[
 N_0[a][b]=\frac3{16(b-a)}\sum_{a<k\le b}\frac1{k(4k-3)}\ (a<b),\qquad
 N_0[a][a]=\frac1{16}\sum_{k\le a}\Bigl(\frac1{k^2}-\frac{16}{(4k-3)^2}\Bigr),
\]
and put $\bar T_h=\frac1{16}\bar L=-(\lambda'I+N_0)$, $\bar P_h(X)=\det(XI+N_0)\in\Q[X]$ (monic of degree $h$). Then $\bar P_h(\lambda'+t)=(-1)^h\det(\bar T_h-tI)$, and Lemma~\ref{lem:integrality} holds verbatim for $\bar Q_h=\den(\bar P_h)\bar P_h$.

\subsection{Content}
For $p\ge5$ and $q=p^e$ the least $k\ge1$ with $q\mid k(4k-3)$ is $k_q=(j_qq+3)/4$, where now $j_q=1$ if $q\equiv1\pmod4$ and $j_q=3$ if $q\equiv3\pmod 4$ (indeed $4k-3=jq$ forces $jq\equiv1\pmod 4$; $k$ and $4k-3$ are coprime for $p\ge5$, and $k_q<q$). So the classes are \emph{swapped} relative to the construction before Lemma~\ref{lem:dual}; the dual certificate (the potentials $u,v$ and Lemma~\ref{lem:dual}) is verbatim. Since $|k_q-j_qq/4|\le\frac34$, the per-level bound is $\delta(q)\le h\,g_{j_q}(q/h)+\frac32$, and the error costs $O(h)$ in total. The two integrals $\int g_1=21/5$ and $\int g_3=37/21$ are now attached to the classes $1$ and $3$ mod~$4$ respectively; by the prime number theorem in progressions the two classes have equal weight, so $\Lm=\frac12(\frac{21}5+\frac{37}{21})=\frac{313}{105}$ again. The primes $2,3$ contribute $O(h\log h)$ as before. Hence $\log\den\bar P_h\le(\frac{313}{105}+o(1))h^2$.

\subsection{Nonvanishing}
Let $\ell\equiv1\pmod4$ be prime, $\ell\ge13$, $K=(\ell+3)/4$ (so $4K-3=\ell$), and $K+1\le h\le\ell-1$; low and high indices are as in \S\ref{sec:nonvanishing}.

\begin{lemma}[$\ell$-singular structure, mirror]\label{lem:mstruct}
Among $k,\ 4k-3,\ b-a,\ 16$ and the numerator $3$ occurring in the entries of $N_0$ only $4k-3$ can be divisible by $\ell$, and this occurs only for $k=K$. Hence every high diagonal entry of $N_0$ contains the term $-\ell^{-2}$, every low--high entry contains $3/(16K\ell(b-a))$, and all other entries lie in $\Z_{(\ell)}$.
\end{lemma}
\begin{proof}
If $\ell\mid4k-3$ with $4k-3\le4h-7<4\ell$, then $4k-3\in\{\ell,2\ell,3\ell\}$; $2\ell$ is even and $3\ell\equiv3\not\equiv1\pmod4$. The numerator $3$ is a unit since $\ell\ne3$. The rest is as in Lemma~\ref{lem:struct}.
\end{proof}

\begin{lemma}\label{lem:m1}
With $D=\mathrm{diag}(I_K,\ell I_{h-K})$, modulo $\ell$
\[
 D(XI+N_0)D\equiv\begin{pmatrix}XI_K+A&\hat B\\\hat B^\top&-I\end{pmatrix},\qquad \hat B_{ab}\equiv\frac3{16K(b-a)}\equiv\frac1{4(b-a)},
\]
because $16K=4\ell+12\equiv12$. The Schur complement of the block $-I$ gives
$\ell^{2(h-K)}\bar P_h(X)\equiv(-1)^{h-K}\det(XI_K+A+\hat B\hat B^\top)$, which is monic of degree $K$ up to sign; so $v_\ell(\den\bar P_h)=2(h-K)$. (The sign in front of $\hat B\hat B^\top$ is the \emph{opposite} of that in Lemma~\ref{lem:1}, because the high block is $-I$ instead of $+I$.)
\end{lemma}

\begin{lemma}[key identity, mirror]\label{lem:m2}
Let $s=\ell-h$, $m=\ell-K$, let $C_s$ be as in Lemma~\ref{lem:2} (note $a+j\le K-1+s<\ell$), and put $Y=-(16X+H^{(2)}_m)$. Then in $\mathbb F_\ell[X]$
\[
 16^K\ell^{2(h-K)}\bar P_h(X)\equiv(-1)^h\det\bigl(YI_K+C_sC_s^\top\bigr).
\]
\end{lemma}
\begin{proof}
Modulo $\ell$, $4K\equiv3$, so $\frac4{4k-3}=\frac1{k-3/4}\equiv\frac1{k-K}\equiv\frac1{k+m}$, hence $\frac3{k(4k-3)}\equiv\frac1{k+m}-\frac1k$ and $\frac{16}{(4k-3)^2}\equiv\frac1{(k+m)^2}$. For $a<b<K$ (so $b+m<\ell$)
\[
 16A_{ab}\equiv-\frac1{b-a}\sum_{k=a+1}^b\Bigl(\frac1k-\frac1{k+m}\Bigr)=-\sum_{j=1}^m\frac1{(a+j)(b+j)},\qquad
 16A_{aa}\equiv\sum_{k\le a}\Bigl(\frac1{k^2}-\frac1{(k+m)^2}\Bigr)=H^{(2)}_m-\sum_{j=1}^m\frac1{(a+j)^2},
\]
that is $16A\equiv H^{(2)}_mI-C_mC_m^\top$. As in the proof of Lemma~\ref{lem:2} (substitute $j=\ell-b$, so $b-a\equiv-(a+j)$), $16\hat B\hat B^\top\equiv\sum_{j=s+1}^mw_jw_j^\top=C_mC_m^\top-C_sC_s^\top$. Therefore $16(A+\hat B\hat B^\top)\equiv H^{(2)}_mI-C_sC_s^\top$, and
$\det\bigl(XI+\frac1{16}(H^{(2)}_mI-C_sC_s^\top)\bigr)=16^{-K}(-1)^K\det(YI+C_sC_s^\top)$. Combine with Lemma~\ref{lem:m1}: the sign is $(-1)^{h-K}(-1)^K=(-1)^h$.
\end{proof}

\begin{theorem}[choice of $h$, mirror]\label{thm:mnonvanish2}
Let $\ell\equiv1\pmod4$ be prime, $\ell\ge13$, $K=(\ell+3)/4$, and $r\in\Q$ with $\ell\nmid\den r$. Then one of
$h\in\{\ell-K,\ \ell-K-1\}=\{\tfrac{3\ell-3}4,\tfrac{3\ell-7}4\}$ satisfies $v_\ell(\bar P_h(r))=-2(h-K)$; in particular $\bar P_h(r)\ne0$.
\end{theorem}
\begin{proof}
Here $s=\ell-h\in\{K,K+1\}$; both values satisfy $K+1\le h\le\ell-1$ and $1\le s\le\ell-K-1$ for $\ell\ge13$, so Lemmas~\ref{lem:m1}, \ref{lem:m2} apply: $\ell^{2(h-K)}\bar P_h(r)$ is an $\ell$-adic unit iff $D_s(Y)=\det(YI_K+N_s)$ does not vanish at $Y=y:=-(16r+H^{(2)}_m)\in\mathbb F_\ell$. Suppose $D_K(y)=D_{K+1}(y)=0$. As before $D_K(0)=\det(C_K)^2\ne0$ (a Cauchy determinant in integers $\le2K-1<\ell$), so $y\ne0$ and $-y\ne0$ is a common eigenvalue of $N_K$ and $N_{K+1}$. Lemma~CC holds for any $K$, any $s\ge K$ and any prime $\ell>K+s$: its proof (Krylov reduction and Theorem~A at $\theta=0$) uses only that $u(K,s)$ is a rational number all of whose prime factors are integers in $[1,K+s]$, and never mentions the residue class of $\ell$ (Theorem~A is an identity in $(n,s,\theta)$). Here $K+s=2K=(\ell+3)/2<\ell$ for $s=K$. This contradicts Lemma~CC.
\end{proof}
(The same argument gives the three-$h$ version, $h\in\{\ell-1,\ell-2,\ell-K\}$, with $\Xi_{11}=H^{(2)}_K$ and $\Xi_{12}=K/(K+1)\ne0$; this is the route to $77.13$.)

\begin{lemma}[primes dividing the denominator, mirror]\label{lem:m3}
Let $\ell\equiv1\pmod4$ be prime, $\ell\ge13$, $K+1\le h\le\ell-1$, $r=p/q$ with $e=v_\ell(q)\ge1$ and $u=\ell^er$. If $e=1$, or $e\ge3$, or $e=2$ and $u\not\equiv+1\pmod\ell$, then $\bar P_h(r)\ne0$ for every admissible $h$.
\end{lemma}
\begin{proof}
As for Lemma~\ref{lem:3}, with high block $\ell^2rI+J$ and now $J\equiv-I$ (Lemma~\ref{lem:mstruct}): the reduced determinants are $u^K$ for $e=1$, $u^K(u-1)^{h-K}$ for $e=2$, and $u^h$ for $e\ge3$.
\end{proof}
Thus the exceptional case is $e=2$, $u\equiv+1$ (instead of $u\equiv-1$); since it still requires $v_\ell(q)=2$, the selection ``$v_\ell(q)\ne2$'' in the proof of Theorem~\ref{thm:measure} is unchanged.
(Checked on $574$ exact cases for $\ell=13,17,29$, all admissible $h$ and $e\le4$; the exception $e=2$, $u\equiv1$ does occur. Lemma~\ref{lem:m2} was checked by exact characteristic polynomials for $h=12,20$ and all $\ell\equiv1\pmod4$ in $(h,4h]$.)

\subsection{Decay and node removal}
\emph{The certificate is the same.} The certificates for $D_{\rm cert}$ (\S\ref{ssec:cert}) and for $g_{\min}$ (\S\ref{ssec:cert2}) concern one fixed function $g$ on the scaled plane; they do not depend on $F$, on $h$, or on where the atoms lie modulo~$1$. What changes is only the contour identity and the $o(1)$ bookkeeping, as follows.

\emph{Reflection.} From $\psi(1+z)=\psi(-z)-\pi\cot\pi z$ and $\psi(\frac14+z)=\psi(\frac34-z)-\pi\cot(\pi z+\frac\pi4)$,
\[
 \bar F=\tilde F+\bar E,\qquad \tilde F(z)=\psi(-z)-\psi(\tfrac34-z),\quad
 \bar E(z)=\frac{-\pi\sin(\pi/4)}{\sin\pi z\,\sin(\pi z+\pi/4)},
\]
with $\tilde F$ analytic and $O(1/|z|)$ in $\Re z<0$, and $\bar E$ $1$-periodic with $|\bar E(x+iy)|\le Ce^{-2\pi|y|}$ for $|y|\ge1$.

\emph{The contour.} Take $N_0=\lfloor r_0h\rfloor$ as before and $x^*=-N_0+\frac38$, the midpoint of the gap $(-N_0,-N_0+\frac34)$ between consecutive poles of $\bar F$, so $x^*\equiv\frac38\pmod1$. On the part of $\Gamma_h=x^*+u(-1\pm3i)$ with $|\Im z|<1$ one has $\Re z\in(-N_0+\frac1{24},-N_0+\frac38]$, inside that gap. The poles to the right of $x^*$ are $-(j+\frac14)$, $0\le j<N_0$ (residue $+1$), and $-(j+1)$, $0\le j<N_0-1$ (residue $-1$); so the set of atoms is
\[
 \bar{\mathcal A}=\{j+\tfrac14:\ 0\le j<N_0\}\ \cup\ \{j+1:\ 0\le j<N_0-1\},\qquad|\bar{\mathcal A}|=2N_0-1,
\]
with $c_\mu=+1$ on the first part and $-1$ on the second, and Proposition~\ref{prop:contour} holds verbatim:
$\bar L=-\sum_{\mu\in\bar{\mathcal A}}c_\mu u_\mu u_\mu^\top+\frac1{2\pi i}\int_{\Gamma_h}\bar E\,vv^\top dz$
(checked entrywise to $2.4\cdot10^{-17}$ at 25 digits for $N_0=1,2,7$, and independently for $(N_0,h)=(4,3),(3,4)$).

\emph{Changes in Andr\'eief's bound, Lemmas~\ref{lem:A}, \ref{lem:B} and Theorem~\ref{thm:N}.} All are in $o(1)$ terms:
\begin{itemize}[leftmargin=*]
\item the support of the measure has $\Re z\le-\frac14$ (the nearest atom is $-\frac14$), so $|v_a|\le4$ (instead of $\frac83$), and the corresponding loss is $\log(1+4h)/h\to0$;
\item $C_E$ and $c_\kappa$ are independent of $h$, by periodicity and $x^*\equiv\frac38\pmod 1$ (the contour crosses the real axis at distance $\frac38$ from the poles at $0$ and $\frac34$ modulo~$1$);
\item in Lemma~\ref{lem:B}, $x^*/h=-r_0+\delta_h$ with $\delta_h=(r_0h-N_0+\frac38)/h\in[\frac3{8h},\frac{11}{8h})$, so $|\delta_h|\le\omega=2^{-30}$ as soon as $h\ge\frac{11}8\cdot2^{30}$; the atoms lie in $-(\Z_{\ge0}+\frac14)/h\,\cup\,-\Z_{\ge1}/h\subset[-r_0,0]$ (the largest is $N_0-\frac34\le r_0h$), two progressions of spacing $1/h$, so $N_\ell\le2h(b_\ell-a_\ell)+2$ exactly as before.
\end{itemize}
The function $g$, the boxes (R), the segment cells (S), the line cells (L) and $I_\rho(\sigma)$ are identical. Hence Theorem~\ref{thm:decay} holds for $\bar P_h$ with $D^+=D_{\rm cert}\ge3.1113790$, Theorem~\ref{thm:N} with $c\le6.8717994$, and Corollary~\ref{cor:window}: $|\bar P_h(\lambda'+t)|\le e^{-(D_{\rm cert}-o(1))h^2}$ for $|t|\le e^{-ch}$.

\begin{remark}
Exact determinants (Arb, rational entries) give $-\log|\det\bar T_h|/h^2=3.030,\,3.060,\,3.075,\,3.084,\,3.090,\,3.094$ at $h=32,48,64,80,96,112$, against $3.068,\,3.084,\,3.093,\,3.098,\,3.101,\,3.104$ for $T_h$. The difference of the logarithms is $\approx1.06\,h+4.7$, linear in $h$ (the near atom $-\frac14$ costs $O(h)$, not $O(h^2)$), so the $h^2$-constants coincide; a margin ``$\eta\approx0.10$'' read off from the finite-$h$ slopes is a finite-$h$ artefact. These computations are consistency checks and are not used in the proof.
\end{remark}

\subsection{The measure}
The proof of Theorem~\ref{thm:measure} is verbatim with: primes $\ell\equiv1\pmod4$ in $(X,Y]$ (using $\vartheta(x;4,1)=x/2+o(x)$); Theorem~\ref{thm:mnonvanish2} in place of Theorem~\ref{thm:nonvanish2}; Lemma~\ref{lem:m3} in place of Lemma~\ref{lem:3}; and $h\in\{\frac{3\ell-7}4,\frac{3\ell-3}4\}$, so that $kL-\frac74<h<\frac34Y$. With $\eta$, $c_{\rm cert}$ as above this yields: for every $\nu>c_{\rm cert}(1/\eta+\frac34)$ there are only finitely many $p/q$ with $|\lambda'-p/q|<q^{-\nu}$, and $c_{\rm cert}(1/\eta+\frac34)\le57.841<57.85$ ($57.8409$ unrounded). With the three values of $h$ one gets $c_{\rm cert}(4/(3\eta)+1)\le77.1217$. This proves Theorem~\ref{thm:mirror}: $\lambda'$ is irrational and $\mu(\lambda')\le57.85$; since $\mu(a\xi+b)=\mu(\xi)$ for rational $a\ne0$ and $b$, also $\mu(G+5\pi^2/48)\le57.85$ and $G+5\pi^2/48\notin\Q$.

\begin{remark}[the two measures do not combine]\label{rem:nocombine}
Since $\lambda+\lambda'=G$ one might hope to deduce something about $G$ from Theorems~\ref{thm:main} and~\ref{thm:mirror}. This is impossible by any black-box argument. Fix a Liouville number $\theta$. For Lebesgue-almost every real $x$ one has $\mu(x)=2$ and $\mu(\theta-x)=2$ (the exceptional sets are null and invariant under $x\mapsto\theta-x$), yet $x+(\theta-x)=\theta$ has $\mu=\infty$. So bounds on $\mu(\xi_1)$ and $\mu(\xi_2)$ say nothing about $\mu(\xi_1+\xi_2)$, and the pair of results gives no measure for $G$ (nor a new proof that $G$ is irrational, which is OpenAI's theorem \cite{OpenAI2026}). Theorems~\ref{thm:main} and~\ref{thm:mirror} concern two different linear forms, $G-5\pi^2/48$ and $G+5\pi^2/48$, in $1,G,\pi^2$.
\end{remark}

%%% ====================================================================
\part{The rectangular hybrid frame and $\mu(\lambda)\le46.45$}\label{part:hybrid}

\section{The rectangular hybrid frame}\label{sec:hframe}

The square frame of \S
ef{sec:matrix}--\S
ef{sec:mirror} uses $h$ nodes on both sides of the L\"owner matrix. In this part we use a \emph{rectangular} frame: the rows and the columns share a long common block of nodes $C$, and carry in addition two short, interleaved blocks $X'$ (rows) and $Y'$ (columns) of the same size $m$. The new feature is the following. The determinant of the square matrix of size $|C|+m$ is divisible, over $\Q(x)$, by $\det$ of the $m\times m$ block $R_{X'Y'}$, and the quotient is a monic polynomial $P_r$ of degree $r=|C|$ in the variable $x$; but the \emph{decay} of the determinant is that of a square frame with $(1+\alpha)r$ nodes ($\alpha=m/r$), while the $\ell$-adic window available for nonvanishing is shorter by the factor $1+\alpha$. Taking $\alpha=\frac18$ gives a much better balance between decay, content and the length of the interval of primes $\ell$.

Throughout this part $\alpha=\tfrac18$ unless said otherwise. All statements are labelled as in the rest of the paper; computer-assisted steps are marked as such, and the points at which the argument rests on external citations, on computation, or on material that is new in this text (and so was not covered by the audits of \S\ref{sec:formal}) are listed in \S\ref{ssec:htodo}, together with the limitations of this part.

\subsection{Notation}
\paragraph{The kernel.} For integers $a,b\ge0$ put
\[
 \Rk(a,b)=\frac1{b-a}\sum_{a<k\le b}\frac1{16k(4k-1)}\quad(a<b),\qquad \Rk(b,a)=\Rk(a,b),\qquad \Rk(a,a)=\Psi(a):=\sum_{k=1}^{a}\Bigl(\frac1{16k^2}-\frac1{(4k-1)^2}\Bigr).
\]
Thus $M_0=-\Rk|_{[0,h)^2}$ (the matrix of the Definition in \S\ref{sec:matrix}), and the entries are exactly (minus) the entries of $\frac1{16}L_F$: for $a\ne b$, $\frac1{16}\frac{F(b)-F(a)}{b-a}=-\Rk(a,b)$, and for $a=b$, $\frac1{16}F'(a)=\lambda-\Psi(a)$ (by \eqref{eq:Ldiag}).

\paragraph{Geometries.} Let $r,m\ge0$ and $e\in\{0,1\}$. The \emph{geometry} $\mathcal G(r,m,e)$ consists of the pairwise disjoint sets of nodes
\[
 C=\{0,1,\dots,r-1\}\cup\{r+2m\ \text{if }e=1\},\qquad X'=\{r+2i:0\le i<m\},\qquad Y'=\{r+2i+1:0\le i<m\}.
\]
Its \emph{scale} is $r=|C|-e$ (the length of the interval part of $C$), and $|C|=r+e$. We write $\mathcal G_1=\mathcal G(r,m,0)$ (``member~1'') and $\mathcal G_2=\mathcal G(r,m,1)$ (``member~2''). In member~2 one extra common node $b=r+2m$ sits directly above the interleaved block $X'\cup Y'$, which occupies $[r,r+2m)$.

\paragraph{The matrices.} For a geometry and a variable $X$ let $\Fu(X)=\Fu_{\mathcal G}(X)$ be the square matrix of size $|C|+m$ with rows indexed by $C\sqcup X'$, columns by $C\sqcup Y'$ (in this order), and entries
\[
 \Fu(X)[a][b]=\Rk(a,b)+X\cdot\mathbf 1[a=b\in C].
\]
(A diagonal entry occurs only at a common node $a=b\in C$; the nodes of $X'$ and $Y'$ are distinct from each other and from $C$.) For a real $t$ put
\[
 \mathcal T_t=-\Fu(-\lambda-t),\qquad\text{i.e.}\quad \mathcal T_t[a][b]=-\Rk(a,b)+(\lambda+t)\mathbf 1[a=b\in C]=\tfrac1{16}L_F[a][b]+t\,\mathbf 1[a=b\in C].
\]
When $m=0$ and $e=0$, $\mathcal T_t=T_h+tI$ with $h=r$, so the hybrid frame contains the frame of \S\ref{sec:matrix}--\S\ref{sec:mirror}. In general
\begin{equation}\label{eq:hsign}
 \det\Fu(-\lambda-t)=(-1)^{|C|+m}\det\mathcal T_t .
\end{equation}
The two sign conventions are both used: $\Fu$ is convenient for $p$-adic analysis (\S\ref{sec:hcontent}, \S\ref{sec:hnv}, \S\ref{sec:hl3}), $\mathcal T_t$ for analysis (\S\ref{sec:hdecay}).

\begin{lemma}[Schur complement; integrality]\label{lem:hschur}
Let $m\ge1$ and suppose $\det\Rk[X',Y']\ne0$ (the $m\times m$ block). Then $\det\Fu(X)=\det\Rk[X',Y']\cdot P(X)$ where $P(X)=\det(XI_{C}+M)$, $M=\Rk[C,C]-\Rk[C,Y']\Rk[X',Y']^{-1}\Rk[X',C]$, is monic of degree $|C|$. In particular $\det\Fu(X)$ has degree $|C|$ in $X$ with leading coefficient $\det\Rk[X',Y']\neq0$, and if $\den_F\in\Z_{>0}$ is the least integer with $\den_F\det\Fu(X)\in\Z[X]$, then $Q(X):=\den_F\det\Fu(X)\in\Z[X]$ has degree $|C|$, and for every rational $x_0=u/v$ in lowest terms with $\det\Fu(x_0)\ne0$,
\[
 |Q(x_0)|\ge v^{-|C|}.
\]
\end{lemma}
\begin{proof}
Schur complement on the block of rows $X'$ and columns $Y'$ (which is $X$-free). The rest is Lemma~\ref{lem:integrality}: $v^{|C|}Q(u/v)\in\Z\setminus\{0\}$.
\end{proof}
That $\det\Rk[X',Y']\ne0$ for the geometries used below is proved in \S\ref{sec:hl3} (Corollary~\ref{cor:hRXY}: it is an $\ell$-adic unit for the primes $\ell$ that we use, hence nonzero).

\subsection{The theorem}
The goal is Theorem~\ref{thm:hybrid}: $|\lambda-a/b|>b^{-46.45}$ for all large $b$, i.e.\ $\mu(\lambda)\le46.45$ (computed value $46.447$). This is the audited hybrid theorem of the title. It improves Theorem~\ref{thm:main} ($57.85$) for $\lambda=G/2-5\pi^2/96$, equivalently $G-5\pi^2/48=2\lambda$ (since $\mu(2\lambda)=\mu(\lambda)$). The hybrid for the mirror constant $\lambda'$ ($\ell\equiv1\pmod 4$) has \emph{not} been done; Theorem~\ref{thm:mirror} ($57.85$) is unchanged. \textbf{The Lean formalization does not cover Theorem~\ref{thm:hybrid}} (\S\ref{sec:formal}).

The theorem rests on five propositions about the geometries with $m=r/8+O(1)$, stated in \S\ref{ssec:hinputs}, and the deduction in \S\ref{sec:hdeduce}. The constants are:
\begin{center}
\begin{tabular}{lll}
\toprule
quantity & value & where\\
\midrule
content $\Lambda_F$ & $\le 3.7211401$ ($=\frac{313}{105}+\frac{554}{840}+\frac{118/21+2\log2}{64}-\frac{129}{4480}$) & \S\ref{sec:hcontent}\\
decay $D_F$ & $\ge 3.9062695$ & \S\ref{sec:hdecay} (computer-assisted)\\
window $c_{\rm hyb}$ & $\le 7.6540647$ & \S\ref{sec:hdecay} (computer-assisted)\\
margin $\eta_{\rm hyb}=D_F-\Lambda_F$ & $\ge 0.1851294$ & \\
window of $\ell$: $r/\ell\to\kappa$ & $\kappa=\frac{3/4}{1+\alpha}=\frac23$ & \S\ref{sec:hnv}\\
exponent $c_{\rm hyb}(1/\eta_{\rm hyb}+\kappa)$ & $\le46.4471<46.45$ & \S\ref{sec:hfinal}\\
\bottomrule
\end{tabular}
\end{center}

\subsection{The five inputs}\label{ssec:hinputs}
Call a geometry $\mathcal G(r,m,e)$ \emph{admissible} if $|m-r/8|\le1$. All $o(1)$ below are as $r\to\infty$, uniform over admissible geometries. For a geometry let $\den_F$ be as in Lemma~\ref{lem:hschur}, and $N_{\rm tot}=|C|+m=r+e+m$.

\begin{proposition}[content]\label{prop:hC}
$\log\den_F\le(\Lambda_F+o(1))\,r^2$ with $\Lambda_F=3.7211401$.
\end{proposition}
\begin{proposition}[decay and window; computer-assisted]\label{prop:hD}
Let $D_F=3.9062695$ and $c_{\rm hyb}=7.6540647$. For every real $t$ with $|t|\le e^{-c_{\rm hyb}r}$,
\[
 |\det\mathcal T_t|\le\exp\bigl(-(D_F-o(1))\,r^2\bigr).
\]
\end{proposition}
\begin{proposition}[nonvanishing]\label{prop:hNV}
Let $\ell\equiv3\pmod4$ be a prime, $\ell\ge7$, $K=(\ell+1)/4$, $0\le m\le K-1$ and $r_1=\ell-K-m$. For every rational $x_0$ with $\ell\nmid\den x_0$, $\det\Fu_{\mathcal G}(x_0)\ne0$ for $\mathcal G=\mathcal G(r_1,m,0)$ or $\mathcal G=\mathcal G(r_1,m,1)$; moreover $v_\ell(\det\Fu_{\mathcal G}(x_0))=-2(|C_{\rm high}|+m)$ for such a $\mathcal G$.
\end{proposition}
\begin{proposition}[$\ell\mid\den$]\label{prop:hL3}
In the setting of Proposition~\ref{prop:hNV}, let $x_0=u/v$ with $v_\ell(v)=e_0\ge1$, $e_0\neq2$. Then $\det\Fu_{\mathcal G}(x_0)\neq0$ for \emph{both} $\mathcal G=\mathcal G(r_1,m,0)$ and $\mathcal G(r_1,m,1)$ (for $m\ge1$).
\end{proposition}

\section{The deduction}\label{sec:hdeduce}

\begin{theorem}\label{thm:hmeasure}
Assume Propositions~\ref{prop:hC}--\ref{prop:hL3}. Let $\eta_{\rm hyb}=D_F-\Lambda_F$, $c=c_{\rm hyb}$ and $\kappa=\frac23$. For every $\nu>c\,(1/\eta_{\rm hyb}+\kappa)$ there are only finitely many rationals $p/q$ with $|\lambda-p/q|<q^{-\nu}$.
\end{theorem}
\begin{proof}
Let $\nu>c(1/\eta_{\rm hyb}+\kappa)$ and suppose $|\lambda-p/q|=:|t|<q^{-\nu}$ with $q$ large; $L=\log q$. Choose $\delta>0$ with $\nu>c(1+\delta)(1/\eta_{\rm hyb}+\kappa)$ and put
\[
 k=\frac{1+\delta}{\eta_{\rm hyb}},\qquad X=\tfrac32kL,\qquad Y=X+(1+\delta)L.
\]
\emph{Choice of $\ell$.} Exactly as in \S\ref{sec:measure}: by the prime number theorem in progressions modulo~$4$, the primes $\ell\equiv3\pmod4$ in $(X,Y]$ have product $>q^{1/2}$, while those with $v_\ell(q)=2$ have product $\le q^{1/2}$. So there is a prime $\ell\equiv3\pmod4$, $X<\ell\le Y$ (so $\ell\ge7$), with $v_\ell(q)\ne2$.

\emph{Choice of the geometry.} Put $K=(\ell+1)/4$, $m=\operatorname{round}\bigl((\ell-K)/9\bigr)$ and $r_1=\ell-K-m$. Then $\ell-K=9m+\delta''$ with an integer $\delta''$, $|\delta''|\le\frac92$ (as $|(\ell-K)/9-m|\le\frac12$), so $r_1=8m+\delta''$, whence $|m-r_1/8|\le\frac9{16}<1$: both geometries $\mathcal G(r_1,m,e)$, $e\in\{0,1\}$, which have the same scale $r_1$, are admissible. Also $0\le m\le K-1$ (for $\ell\ge7$: $m=1=K-1$ if $\ell=7$, and $m\le K-2$ for all $\ell\ge11$, since $(\ell-K)/9+\frac12\le K-2$ for $K\ge3$; $K=3$ gives $m=1$). Moreover $\bigl|r_1-\tfrac23\ell\bigr|\le1$, since $\ell-K=\frac{3\ell-1}4$ and $r_1=\frac89(\ell-K)+\frac{\delta''}9=\frac23\ell-\frac29+\frac{\delta''}9$ with $\frac{|\delta''|}9\le\frac12$.

Let $x_0=-p/q$. If $\ell\nmid q$, Proposition~\ref{prop:hNV} gives $e\in\{0,1\}$ with $\det\Fu_{\mathcal G}(x_0)\ne0$ for $\mathcal G=\mathcal G(r_1,m,e)$. If $\ell\mid q$ then $v_\ell(q)\in\{1\}\cup[3,\infty)$ and Proposition~\ref{prop:hL3} gives the same (here $m\ge1$ because $\ell$ is large). In every case let $\mathcal G$ be such a geometry, $r=r_1$ its scale, $|C|=r+e$. Then
\begin{equation}\label{eq:hr}
 kL-1<\tfrac23X-1\le r_1=r\le\tfrac23\ell+1\le\tfrac23Y+1=(1+\delta)\bigl(\tfrac1{\eta_{\rm hyb}}+\tfrac23\bigr)L+1,
\end{equation}
(using $r_1\ge\frac23\ell-1>\frac23X-1=kL-1$), and $r\ge2^{30}$ for large $q$.

\emph{Lower bound.} By Lemma~\ref{lem:hschur}, with $Q=\den_F\det\Fu_{\mathcal G}$: $|Q(-p/q)|\ge q^{-(r+e)}=e^{-(r+1)L}$.

\emph{Upper bound.} Write $-p/q=-\lambda-t$ with $t=p/q-\lambda$. By \eqref{eq:hsign}, $|\det\Fu_{\mathcal G}(-p/q)|=|\det\mathcal T_t|$. By \eqref{eq:hr} and the choice of $\delta$, $c\,r<\nu L$ for large $q$, so $|t|<e^{-\nu L}\le e^{-c r}$ and Proposition~\ref{prop:hD} gives $|\det\mathcal T_t|\le e^{-(D_F-o(1))r^2}$. With Proposition~\ref{prop:hC},
\[
 |Q(-p/q)|\le\den_F|\det\mathcal T_t|\le\exp\bigl(-(\eta_{\rm hyb}-o(1))r^2\bigr).
\]
\emph{Contradiction.} Comparing, $(\eta_{\rm hyb}-o(1))r^2\le(r+1)L$, i.e. $(\eta_{\rm hyb}-o(1))r\le L(1+1/r)$. But $r>kL-1$ gives $\eta_{\rm hyb}r\ge(1+\delta)L-\eta_{\rm hyb}$. Contradiction for $q$ large.
\end{proof}

\begin{remark}[where each constant enters]\label{rem:hsens}
The exponent is $c\,(1/\eta+\kappa)$: $c/\eta=41.35$ is the cost of a window of exponential radius around $\lambda$ for a determinant of size $\eta r^2$; $c\kappa=5.10$ is the cost of the length $(1+\delta)L$ of the interval of $\ell$ multiplied by $r/\ell\to\kappa=2/3$. Proposition~\ref{prop:hL3} is what allows $\ell\mid q$ (without it the interval has length $(2+\delta)L$ and $\kappa$ becomes $2\kappa$, giving $\approx51.5$). Compared with the square frame ($\kappa=\frac34$, $c=6.872$, $\eta=0.1304$) the hybrid has a larger window constant $c$ and a smaller $\kappa$, but a much larger $\eta$. The derivative of the exponent is $\partial/\partial D_F=-223$, so one more digit of $D_F$ matters.
\end{remark}

\section{The denominator of the hybrid determinant}\label{sec:hcontent}

This section proves Proposition~\ref{prop:hC}. The structure is that of \S\ref{sec:content}: a tropical bound for each prime power with an explicit dual certificate, summed with the prime number theorem in progressions. Two things are new. First, the dual certificate of Lemma~\ref{lem:dual} has to be replaced by asymmetric certificates (D1--D6 below), because rows and columns of the rectangular frame differ. Second, at a positive proportion of primes the tropical bound is not sharp, and a $p$-adic rank argument (Lemma~E, \S\ref{ssec:hLemmaE}) removes $m$ units from it; this is what turns the tropical content $3.7499$ into $\Lambda_F=3.7211$.

\subsection{Gauss valuations and the tropical lemma}
Fix a prime $p$. Let $\mathcal O_p\subset\Q(X)$ be the localization of $\Z_{(p)}[X]$ at the prime ideal $(p)$; it is a discrete valuation ring with uniformizer $p$ and residue field $\F_p(X)$. Its valuation $v$ is the \emph{Gauss valuation}: for $f\in\Q[X]$, $v(f)=\min_kv_p(\text{coefficient of }X^k)$, and $v(fg)=v(f)+v(g)$ (Gauss's lemma). For $f\in\Q[X]$ with least common denominator $\den f$ of the coefficients, $v_p(\den f)=\max(0,-v(f))$. In particular, for $f=\det\Fu(X)$,
\begin{equation}\label{eq:hden}
 v_p(\den_F)=\max\bigl(0,-v(\det\Fu(X))\bigr).
\end{equation}
\begin{lemma}[tropical expansion]\label{lem:htrop}
Let $B$ be a square matrix with entries in $\Q[X]$ (rows $I$, columns $J$), $W_{ab}=\max(0,-v(B_{ab}))$ and $\tau(W)=\max_\sigma\sum_aW_{a\sigma(a)}$, over bijections $\sigma:I\to J$. Then $v(\det B)\ge-\tau(W)$. If $u_a,v_b\ge0$ (or real numbers) satisfy $W_{ab}\le u_a+v_b$ for all $a,b$, then $\tau(W)\le\sum_au_a+\sum_bv_b$.
\end{lemma}
\begin{proof}
Each term of the Leibniz expansion has $v(\prod_aB_{a\sigma(a)})=\sum_av(B_{a\sigma(a)})\ge-\sum_aW_{a\sigma(a)}\ge-\tau(W)$, and $v$ of a sum is at least the minimum. For the second claim $\sum_aW_{a\sigma(a)}\le\sum_a(u_a+v_{\sigma(a)})$.
\end{proof}
For $\Fu(X)$ the entry $\Rk(a,b)+X\mathbf 1[a=b\in C]$ has $v=\min(0,v_p(\Rk(a,b)))$ (as $v(X)=0$), so \emph{for every entry}
\begin{equation}\label{eq:hW}
 W_{ab}=W_{ab}(p)=\max\bigl(0,-v_p(\Rk(a,b))\bigr).
\end{equation}
(Compare Lemma~\ref{lem:trop}, whose proof by principal minors is replaced here by the multiplicativity of the Gauss valuation; this is the form needed for rectangular matrices.)

\subsection{The level majorant}
Let $p$ be an odd prime and $\nu(k)=v_p(k(4k-1))$, $\mu(a)=\max_{1\le k\le a}\nu(k)$. Since $16$ is a $p$-adic unit, the computation leading to \eqref{eq:major} gives, for $a<b$ and for the diagonal,
\[
 W_{ab}\le v_p(b-a)+\max_{a<k\le b}\nu(k),\qquad W_{aa}\le2\mu(a)\quad(a\in C).
\]
For a prime power $q=p^e$ put $S_q=\{k\ge1:q\mid k(4k-1)\}=\{k\equiv0\text{ or }k_q\ (q)\}$, where $k_q=(j_qq+1)/4$, $j_q\in\{1,3\}$, $j_qq\equiv3\pmod 4$ (as in \S\ref{sec:content}; $4k_q\equiv1\pmod q$), and let $\beta(a)=\#(S_q\cap[1,a])$. The consecutive elements of $S_q$ are at distance $k_q$ and $q-k_q$ alternately, both $<q$, and the first block $[0,k_q)$ has length $<q$. Define
\[
 w^{(q)}(a,b)=\mathbf 1[\beta(a)\ne\beta(b)]+\mathbf 1[q\mid a-b]\ (a\ne b),\qquad w^{(q)}(c,c)=2\cdot\mathbf 1[c\ge k_q]\ (c\in C).
\]
Because $v_p(n)=\sum_{e\ge1}\mathbf 1[p^e\mid n]$ and $\max_k\nu(k)=\sum_e\mathbf 1[\exists k\in(a,b]:p^e\mid k(4k-1)]$, we get, exactly as in \S\ref{sec:content},
\begin{equation}\label{eq:hmaj}
 W(p)\le\sum_{e\ge1}w^{(p^e)}\ \text{entrywise},\qquad \tau(W(p))\le\sum_{e\ge1}\tau(w^{(p^e)}),
\end{equation}
and $w^{(q)}\in\{0,1,2\}$: if $q\mid a-b$, $a\ne b$, then $|a-b|\ge q$, so $(a,b]$ contains a multiple of $q$ and $\beta(a)\ne\beta(b)$; hence $w=2$ on congruent pairs and on high diagonal cells, $w=1$ on non-congruent pairs in different blocks, $w=0$ otherwise. Moreover $w^{(q)}\equiv0$ if $q>4N'$, where $N'$ bounds all indices (then $k_q>N'$ and no two indices are congruent).

\subsection{Dual certificates}\label{ssec:hdual}
Let $8m\le r$, $N=r+2m$, and consider the geometry $\mathcal G(r,m,0)$ (indices in $[0,N)$; $Z=[r,N)$ the zone, $X'$ the even and $Y'$ the odd positions of $Z$ relative to $r$). Fix an odd prime power $q$, $j=j_q$, $k=k_q$, and write $w=w^{(q)}$ on (rows $C\sqcup X'$) $\times$ (columns $C\sqcup Y'$). As in Lemma~\ref{lem:dual} let
\[
 \delta_A(q)=\min\bigl(2r,\ 2(r-n)+2\min(n,(r-q)_+)\bigr),\qquad n=\min(r,k_q),
\]
the value of the $C$-block. Let $\varphi:C\to\{0,1\}$ be the potential of Lemma~\ref{lem:dual} with $h=r$: $\varphi(c)=1$ for $c\ge k_q$ (``high'') and $\varphi(c)=\mathbf 1[c\le r-1-q]$ for $c<k_q$. By the proof of Lemma~\ref{lem:dual}, $w(c,c')\le\varphi(c)+\varphi(c')$ for $c,c'\in C$ (including the diagonal) and $2\sum_c\varphi(c)=\delta_A(q)$. For $q\le r$ the set $U=[r-q,k_q)$ (empty if $r-q\ge k_q$) consists of the ``low unpaired'' indices ($\varphi=0$); a zone node $x$ is \emph{low} if $x-q\in U$, and $L$ denotes the number of low zone nodes; for $q\ge2m$, $L=\min(|U|,2m)$ and $|U|=(k_q-r+q)_+$.

\begin{proposition}[dual certificates]\label{prop:hdual}
Let $q$ be an odd prime power. Then $\tau(w^{(q)})\le B(q):=\delta_A(q)+E(q)$, where
\begin{center}
\begin{tabular}{lll}
\toprule
range & $E(q)$ & certificate\\
\midrule
(R1) $q<2m$ & $m+M_q$, $\ M_q=m-q\|m/q\|$ & D1\\
(R2) $2m\le q\le r$ & $\min\bigl(2m+L,\ m+L+|U|\bigr)$ & D2, D3\\
(R3) $r<q<N$ & $2m+(N-q)$ & D4\\
(R4) $q\ge N$, $k_q\le r$ & $2m$ & D5\\
(R5) $q\ge N$, $k_q>r$ & $(N-k_q)_+$ & D6\\
\bottomrule
\end{tabular}
\end{center}
($\|\cdot\|$ is the distance to the nearest integer.) In (R1) $\delta_A=2r$.
\end{proposition}
\begin{proof}
In every case we exhibit potentials $u_a$ (rows) and $v_b$ (columns), $\ge0$, with $w(a,b)\le u_a+v_b$ everywhere, and apply Lemma~\ref{lem:htrop}. Two remarks used repeatedly: (i) a pair of distinct indices that are not congruent mod $q$ has $w\le1$; (ii) pairs of nodes of the zone are not congruent, since $|x-y|\le2m-1<q$ when $q\ge2m$ (and $q\neq2m$ is odd).

\emph{D1 ($q$ arbitrary; used for $q<2m$).} For a residue class $\varrho$ mod $q$ let $x_\varrho,y_\varrho$ be the numbers of nodes of $X'$, $Y'$ in $\varrho$. Rows in class $\varrho$ get $u=\tfrac12,\tfrac32,1$ and columns $v=\tfrac32,\tfrac12,1$ according as $x_\varrho>y_\varrho$, $x_\varrho<y_\varrho$, $x_\varrho=y_\varrho$. Pairs in the same class (congruent pairs and all diagonal cells) have $u+v=2\ge w$; pairs in different classes have $w\le1\le\frac12+\frac12\le u+v$. The value is $\sum_\varrho[(\#C_\varrho+x_\varrho)u+(\#C_\varrho+y_\varrho)v]=2r+2m-\frac12\sum_\varrho|x_\varrho-y_\varrho|$. The nodes of $Z$ are $2m$ consecutive integers alternating between $X'$ and $Y'$, and $q$ is odd, so the members of one class alternate in type: $|x_\varrho-y_\varrho|\in\{0,1\}$, with $1$ iff the class has an odd number of members. Writing $2m=Aq+s'$ ($0\le s'<q$), $s'$ classes have $A+1$ members and $q-s'$ have $A$; the number of odd classes is $s'$ if $A$ is even and $q-s'$ if $A$ is odd, which equals $2q\|m/q\|$ in both cases. So the value is $2r+2m-q\|m/q\|=2r+m+M_q$. For $q<2m\le r/4$ one has $k_q<q\le r-q$, whence $\delta_A=2r$.

\emph{D2 ($2m\le q\le r$).} On $C$ take $\varphi$ (both for rows and columns); a zone node $x$ (row or column) gets $1+\mathbf 1[x\text{ low}]$. A zone--$C$ pair with $\varphi(c)=1$ has $u+v\ge2\ge w$. If $\varphi(c)=0$, i.e.\ $c\in U$: either $z\not\equiv c$ and $w\le1\le u_z$, or $z\equiv c$, and then $0<z-c\le(N-1)-(r-q)<2q$ forces $z=c+q$, so $z$ is low and $u_z=2$. Zone--zone pairs have $w\le1\le1+1$. The value is $\delta_A+2m+L$.

\emph{D3 ($2m\le q\le r$).} Start from $\varphi$ on $C$. (1) For each \emph{high} zone node $z\in X'$ put $u=\frac12$ on all rows and $v=\frac32$ on all columns of its residue class; for each high $y\in Y'$ put $u=\frac32$ on rows and $v=\frac12$ on columns of its class. These classes contain no other zone node (zone nodes are pairwise incongruent) and no index of $U$ (an index $c\in U$ congruent to a zone node $x$ forces $x=c+q$, i.e.\ $x$ low); so every $C$-member of such a class has $\varphi=1$ and $u+v=2$ as before. (2) For $c\in U$ put $u_c=v_c=\frac12$. (3) A low zone node gets potential $\frac32$, a high one $\frac12$ (as a row, resp.\ column; these agree with the class rule in (1)). Check: every potential is $\ge\frac12$, which settles all non-congruent pairs ($w\le1$). Congruent pairs: inside a class modified in (1), $\frac12+\frac32=2$; low $x$ with $x-q\in U$: $\frac32+\frac12=2$; low $x$ with $x-2q,x-3q,\dots\ (<r-q)$, which have $\varphi=1$ and unmodified class: $\frac32+1$; $C$--$C$ congruent pairs in unmodified classes: $1+1$ (indices of $U$ have no congruent partner in $C$ since $c+q\ge r$ and $c-q<0$); diagonal cells: $\frac12+\frac32$ in a modified class, $1+1$ in an unmodified one, and the diagonal cells of $U$ have $w=0$. The value is $\delta_A+|U|+m+L$ (the potentials of $U$ add $|U|$, the zone adds $\frac32L+\frac12(2m-L)=m+L$).

\emph{D4 ($r<q<N$).} $\varphi'=\mathbf 1[c\ge k_q]$ on $C$ and $1+\mathbf 1[x\ge q]$ on the zone. Since $2q>N$ the only congruences inside $[0,N)$ are $x\equiv x-q$ with $x\ge q$, and then $x-q<2m<r$ lies in $C$ and $u_x=2$. $C$--$C$: no congruences ($q>r$), so $w\le1$, and $w=0$ inside the first block; if one of $c,c'$ is high, $\varphi'$ gives $\ge1$ (or $2$ for both high, and for high diagonal cells). The value is $2(r-k_q)_++2m+(N-q)$, and $\delta_A=2(r-n)=2(r-k_q)_+$ here.

\emph{D5 ($q\ge N$, $k_q\le r$).} $\varphi'$ on $C$, potential $1$ on the zone. There are no congruences, so $w\le1$ off the diagonal, and the diagonal cells are covered by $\varphi'$. Value $2(r-k_q)+2m=\delta_A+2m$.

\emph{D6 ($q\ge N$, $k_q>r$).} Potential $0$ on $C$ and $\mathbf 1[x\ge k_q]$ on zone nodes. The only element of $S_q$ below $N$ is $k_q$ (if $k_q<N$). Then $w(z,c)=\mathbf 1[z\ge k_q]$ for $c\in C$; for two zone nodes $w=1$ only if $k_q\in(\min,\max]$, and then the larger node has potential $1$; $C\times C$ has $w=0$ (one block, no congruences, no high diagonal cell). Value $(N-k_q)_+$, and $\delta_A=0$.
\end{proof}

\begin{remark}
These certificates are also sharp (the corresponding primal assignments exist, and the exact tropical value computed by the Hungarian algorithm equals $B(q)$ for all odd $q$ and all tested $(r,m)$, with $\tau_p$ computed from the exact entries equal to $\sum_e B(p^e)$ for every prime $p\ge17$). Sharpness is not used.
\end{remark}

\subsection{Asymptotics of the level bounds}
Put $\theta=q/r$ and $a=m/r$. For $j\in\{1,3\}$ let $g_j$ be the function of \S\ref{sec:content} (so that $\delta_A(q)=r\,g_j(\theta)+O(1)$ and $\int g_1=\frac{21}5$, $\int g_3=\frac{37}{21}$), and let $\varepsilon_j(\theta;a)$ be the piecewise function
\begin{center}
\begin{tabular}{ll}
\toprule
$\theta$ & $\varepsilon_j(\theta;a)$\\
\midrule
$(0,2a)$ & $2a-\theta\|a/\theta\|$\\
$(2a,\ 4/(4+j))$ & $a$\\
$\bigl(\tfrac4{4+j},\ \tfrac{4(1+2a)}{4+j}\bigr)$, $\ell:=\tfrac{(4+j)\theta}4-1$ & $\min(a+2\ell,\ 2a+\ell)$\\
$\bigl(\tfrac{4(1+2a)}{4+j},\ 1\bigr)$ & $4a$\\
$(1,\ 1+2a)$ & $4a-(\theta-1)$\\
$(1+2a,\ 4/j)$ & $2a$\\
$(4/j,\ 4(1+2a)/j)$ & $1+2a-j\theta/4$\\
$>4(1+2a)/j$ & $0$\\
\bottomrule
\end{tabular}
\end{center}
(for $a\le\frac18$; for $j=1$ the fourth interval is empty when $a=\frac18$; the formula is the translation of Proposition~\ref{prop:hdual}: in (R2) $|U|=\ell r+O(1)$, $L=\min(\ell,2a)r+O(1)$, so that $E=\min(2m+L,m+L+|U|)$ is $\min(a+2\ell,2a+\ell)r$ for $\ell\le2a$ and $4ar$ beyond; in (R3) $E=(2a+(1+2a-\theta))r$; in (R4) $E=2m=2ar$; in (R5) $N-k_q=(1+2a-j\theta/4)r+O(1)$.) One checks directly that $\varepsilon_j$ is continuous in $\theta$ (and jointly continuous in $(\theta,a)$), bounded by $4a$, supported on $[0,4(1+2a)/j]$, and Lipschitz in $a$ with constant $\le8$ (each piece has $|\partial_a|\le4$). For every odd prime power $q$ and all $r,m$ with $8m\le r$,
\begin{equation}\label{eq:hB}
 B(q)\le r\,f_j(q/r;a)+C_0,\qquad f_j=g_j+\varepsilon_j,\qquad j=j_q,\ a=m/r,\qquad C_0=0.
\end{equation}
\begin{proof}[Proof of \eqref{eq:hB}]
Put $\theta=q/r$ and $n'=\min(r,jq/4)$. The function $\delta(n)=\min\bigl(2r,2(r-n)+2\min(n,(r-q)_+)\bigr)$ is non-increasing in $n$ (constant for $n\le(r-q)_+$, then of slope $-2$), and $\delta(n')=rg_j(\theta)$: for $\theta\le\frac4{4+j}$ one has $n'\le r-q$ and $\delta=2r$; for $\frac4{4+j}\le\theta\le1$, $\delta=2(r-\frac{jq}4)+2(r-q)$; for $1\le\theta\le\frac4j$, $\delta=2(r-\frac{jq}4)$; beyond, $\delta=0$. Since $n_q=\min(r,k_q)\ge n'$ ($k_q=\frac{jq}4+\frac14$), $\delta_A(q)\le rg_j(\theta)$. It remains to compare $E(q)$ with $r\varepsilon_j(\theta;a)$; as $a\le\frac18$, $2a<\frac4{4+j}$, $\frac{4(1+2a)}{4+j}\le1$ and $1+2a<\frac4j$, so the intervals of the table are in the stated order.
(R1) $\theta<2a$: $E=2m-q\|m/q\|=r(2a-\theta\|a/\theta\|)=r\varepsilon_j$, and $\delta_A=2r=rg_j$.
(R2) $2a\le\theta\le1$: with $\ell=\frac{(4+j)\theta}4-1$ one has $k_q-r+q=r\ell+\frac14$, an integer, so $|U|=(r\ell+\frac14)_+$ and $L=\min(|U|,2m)$. The function $\mathcal E(x)=\min\bigl(2m+\min(x,2m),\,m+\min(x,2m)+x\bigr)$ is non-decreasing and $2$-Lipschitz, $E=\mathcal E(|U|)$, and $\mathcal E(r\ell_+)=r\varepsilon_j(\theta;a)$ (it is $m$, $\min(2m+r\ell,m+2r\ell)$, $4m$ for $\ell\le0$, $0\le r\ell\le2m$, $r\ell\ge2m$). If $\ell<0$, then $r\ell+\frac14<\frac14$ is a non-positive integer, $|U|=0$ and $E=m=r\varepsilon_j$. If $\ell\ge0$, then $|U|=r\ell+\frac14$ and $E\le r\varepsilon_j+\frac12$; but then $k_q>r-q$, $k_q<r$, and $\delta_A(q)=2(r-k_q)+2(r-q)=rg_j(\theta)-\frac12$. In both cases $\delta_A+E\le r(g_j+\varepsilon_j)$.
(R3) $r<q<N$: $E=2m+(r+2m-q)=r(4a-(\theta-1))=r\varepsilon_j$.
(R4) $q\ge N$, $k_q\le r$: then $1+2a\le\theta<\frac4j$ and $E=2m=r\varepsilon_j$.
(R5) $q\ge N$, $k_q>r$: $E=(r+2m-k_q)_+$. If $\theta<\frac4j$ then $E<2m=r\varepsilon_j$; if $\theta\ge\frac4j$ then $E\le r(1+2a-\frac{j\theta}4)_+=r\varepsilon_j$, as $k_q>\frac{jq}4$.
In every range $B(q)=\delta_A(q)+E(q)\le rf_j(\theta;a)$.
\end{proof}
(This agrees with the exact evaluation: $\max_q[B(q)-rf_j(q/r;a)]=0$ over all odd $q$ at $r=10^4,4\cdot10^4$, $m=r/8$, and at $21$ further pairs $(r,m)$, $97\le r\le10^4$, $m\in\{\lfloor r/8\rfloor-5,\lfloor r/8\rfloor-1,\lfloor r/8\rfloor\}$; script \texttt{check\_hB.py}.)

\begin{lemma}[summation]\label{lem:hS}
Let $f:[0,\infty)\to[0,\infty)$ be bounded, Riemann integrable and supported in $[0,K_0]$, and $a\in\{1,3\}$. Then $\displaystyle\lim_{r\to\infty}\frac1r\sum_{p\equiv a\,(4)}f(p/r)\log p=\frac12\int_0^\infty f$.
\end{lemma}
\begin{proof}
Given $\epsilon>0$ choose a partition $0=\theta_0<\dots<\theta_n=K_0$ with $\sum_ic_i(\theta_i-\theta_{i-1})\le\int f+\epsilon$, $c_i=\sup_{[\theta_{i-1},\theta_i]}f$. Every prime $p\le K_0r$ lies in some $(\theta_{i-1}r,\theta_ir]$, so $\sum_pf(p/r)\log p\le\sum_ic_i[\vartheta(\theta_ir;4,a)-\vartheta(\theta_{i-1}r;4,a)]$. By the prime number theorem in progressions, $\vartheta(x;4,a)=x/2+o(x)$, so each bracket is $\frac12(\theta_i-\theta_{i-1})r+o(r)$ ($n$ fixed). Hence $\limsup\le\frac12(\int f+\epsilon)$. The lower Darboux sum gives the same for $\liminf$.
\end{proof}
(Both an upper and a lower bound for $\vartheta(x;4,a)$ are needed here, since $f_j$ is not monotone; in \S\ref{sec:content} monotonicity of $g_j$ made an upper bound enough.)

\begin{proposition}[tropical content law]\label{prop:htropical}
For $\mathcal G(r,m,e)$ admissible, the quantity $T(r):=\sum_{p\ge5}\sum_e\tau(w^{(p^e)})\log p$ (over all prime powers $p^e\le4N'$) satisfies
\[
 T(r)\le\bigl(\Lambda_{\rm trop}+o(1)\bigr)r^2,\qquad \Lambda_{\rm trop}=\frac{313}{105}+\frac{554}{105}\alpha+\Bigl(\frac{118}{21}+2\log2\Bigr)\alpha^2\Big|_{\alpha=1/8}=3.7499347\ldots
\]
\end{proposition}
\begin{proof}
\emph{Extension to admissible geometries.} Let $m_0=\lfloor r/8\rfloor$, so $m_0\le m\le m_0+2$ and $8m_0\le r$. The matrix for $\mathcal G(r,m,e)$ is obtained from that of $\mathcal G(r,m_0,0)$ by adding at most $2+1$ rows and the same number of columns (the added nodes are $r+2i,r+2i+1$ for $m_0\le i<m$ and $b$). Each $w^{(q)}\le2$, so extending a certificate by the potential $2$ on every new row and column keeps $w\le u+v$ and adds $\le4$ per new row-column pair. Hence $\tau(w^{(q)}(\mathcal G(r,m,e)))\le B_{m_0}(q)+12$.

\emph{Prime powers $p^e$, $e\ge2$, and $p\le3$.} For $q=p^e$ with $e\ge2$ use the trivial potentials $u=v=1$: $\tau(w^{(q)})\le2(N_{\rm tot})$, and $\sum_{p^e\le4N',e\ge2}\log p=\psi(4N')-\vartheta(4N')=O(\sqrt r)$ (Chebyshev), so these levels contribute $O(r^{3/2})$. For $p=3$ all levels $3^e\le4N'$ are bounded by the trivial potentials, total $O(r\log r)$. For $p=2$: every entry of $\Rk$ has $-v_2\le4+2\log_2N'$ (from $16$, $k$, $b-a$ and the ultrametric inequality), so $\tau_2\le N_{\rm tot}(4+2\log_2N')$, a contribution $O(r\log r)$.

\emph{Level $1$.} By \eqref{eq:hB} (applied with $m_0$), for $p\ge5$, $p\le4N'$:
$\tau(w^{(p)})\log p\le(r f_{j_p}(p/r;a_0)+C_0+12)\log p$, $a_0=m_0/r\in(\frac18-\frac1r,\frac18]$, and $f_j$ is $8$-Lipschitz in $a$, so $rf_j(p/r;a_0)\le rf_j(p/r;\frac18)+8$. The $O(1)$ errors sum to $O(\vartheta(4N'))=O(r)$. By Lemma~\ref{lem:hS},
$\sum_{p\equiv a_j(4)}rf_j(p/r;\tfrac18)\log p=\frac{r^2}2\int f_j+o(r^2)$. Hence $T(r)\le\frac{r^2}2\bigl(\int f_1+\int f_3\bigr)+o(r^2)$.

\emph{The integral.} Put $I=\int_{1/2}^\infty\|t\|t^{-3}dt$. On each cell $[n-\frac12,n+\frac12]$ the antiderivative of $|t-n|t^{-3}$ is elementary; summing, $\sum_n[\frac1n-\frac1{2n-1}-\frac1{2n+1}]+\sum_n[\frac1{4(n-\frac12)^2}-\frac1{4(n+\frac12)^2}]=(1-2\log2)+1$, so $I=2-2\log2$. The substitution $t=a/\theta$ gives $\int_0^{2a}\theta\|a/\theta\|d\theta=a^2I$. The seven pieces of $\int\varepsilon_j$ are
$4a^2-a^2I$; $a(\frac4{4+j}-2a)$; $\frac{22a^2}{4+j}$ (the middle piece: $\frac4{4+j}\int_0^{2a}\min(a+2\ell,2a+\ell)d\ell=\frac4{4+j}\cdot\frac{11a^2}2$); $\frac{4a(j-8a)}{4+j}$; $6a^2$; $2a(\frac4j-1-2a)$; $\frac{8a^2}j$. Adding,
\[
 \int\varepsilon_1=\tfrac{38}5a+(10-I)a^2,\qquad\int\varepsilon_3=\tfrac{62}{21}a+(\tfrac{110}{21}-I)a^2 .
\]
Therefore $\frac12\sum_j\int f_j=\frac12(\frac{21}5+\frac{37}{21})+\frac12\sum_j\int\varepsilon_j=\frac{313}{105}+\frac{554}{105}a+(\frac{160}{21}-I)a^2$, and $\frac{160}{21}-I=\frac{118}{21}+2\log2$. At $a=\frac18$: $2.9809524+0.6595238+0.1094585=3.7499347$.
\end{proof}
(Negative control: at $\alpha=0$ every $\varepsilon_j$ vanishes and the law is $\frac{313}{105}$, the content constant of Theorem~\ref{thm:content}. The same bound computed from the exact entries, $\tau_p$ from exact $p$-adic valuations, agrees with $\sum_eB(p^e)$ for every prime $p\ge17$ and $(r,m)\in\{(64,8),(96,12),(128,16),(160,20)\}$, and the finite-$r$ values of the full tropical bound are $3.7027,3.6927,3.7216,3.7403$, increasing towards $3.7499$.)

\subsection{Lemma E: a rank saving at clean primes}\label{ssec:hLemmaE}
The tropical bound is not sharp: at primes $p\in(\frac58r,\frac45r)$, $p\equiv3\pmod4$, and $p\in(\frac57r,r)$, $p\equiv1\pmod4$, the true denominator is smaller by $m$ units (up to $O(1)$). We prove this in the Gauss valuation of $\mathcal O_p$. Note that this is a statement about the polynomial $\det\Fu(X)$ with $X$ a \emph{formal variable}: it is false for $\det\Fu(x_0)$ at specialisations $x_0\in\Q$ with $16x_0\equiv H^{(2)}_{K'-1}\pmod p$ (below), and it is never applied to a specialised value in this paper.

The case $e=1$ (member~2, with the extra common node $b$) is treated below together with $e=0$. This treatment is new in this text and is not covered by the audits listed in \S\ref{sec:formal}; it was re-derived line by line by a further independent check, and confirmed by two independent exact computations (Remark~\ref{rem:hmember2}).
\paragraph{Set-up.} Let $\mathcal G=\mathcal G(r,m,e)$ be admissible, $p\ge11$ prime, as specified below. Put $Z_e=[r,r+2m+e)$ for the zone nodes: $X''=X'\cup\{b\}$, $Y''=Y'\cup\{b\}$ if $e=1$ and $X''=X',Y''=Y'$ if $e=0$. Take the \emph{interval} $C_0=[0,r)$, $A=\Rk[C_0,C_0]+XI$, and write
\begin{equation}\label{eq:hschur2}
 \Fu(X)=\begin{pmatrix}A&\Rk[C_0,Y'']\\ \Rk[X'',C_0]&\Rk[X'',Y'']+eX\,E_{bb}\end{pmatrix},\qquad \det\Fu(X)=\det A\cdot\det T,
\end{equation}
$T=\Rk[X'',Y'']+eX\,E_{bb}-\Rk[X'',C_0]A^{-1}\Rk[C_0,Y'']$ (an $m''\times m''$ matrix, $m''=m+e$). Here $E_{bb}$ is the matrix unit at $(b,b)$. ($\det A=X^r+\dots\ne0$.) By Lemma~\ref{lem:htrop},
\begin{equation}\label{eq:hredE}
 v(\det\Fu(X))\ \ge\ -\tau(W_A)+v(\det T),\qquad v(\det T)\ge\sum_{z}\min_y v(T_{zy})\ \text{(Leibniz)}.
\end{equation}

\paragraph{Singular indices and levels.} Let $F(n)=-\sum_{k\le n}\frac1{4k(4k-1)}$, so $\Rk(a,b)=\frac{F(a)-F(b)}{4(b-a)}$. The indices $k$ with $p\mid k(4k-1)$ are $k\equiv0$ and $k\equiv4^{-1}\pmod p$. Since $1/(4k-1)$ and $1/k$ have valuation $\ge-1$ for $k<p^2/4$, $v(F(n))\ge-1$ and $v(\Psi(n))\ge-2$ on the indices considered. For $n$ put $\mathsf L(n)=\operatorname{lead}_{-1}F(n)\in\F_p$ and $\pi(n)=\operatorname{lead}_{-2}\Psi(n)$ (residues of $pF$, $p^2\Psi$). They are step functions, jumping at singular $k$: if $4k-1=j'p$ the jump of $\mathsf L$ is $-1/j'$ and of $\pi$ is $-1/j'^2$ (as $4k\equiv1$); at $k=p$ the jumps are $+\frac14$ and $+\frac1{16}$. Consequently, for $a\ne b$ among the indices:
\begin{itemize}[leftmargin=*]
\item[(V1)] if $p\nmid b-a$: $v(\Rk(a,b))\ge-1$ and $\operatorname{lead}_{-1}\Rk(a,b)=\frac{\mathsf L(a)-\mathsf L(b)}{4(b-a)}$;
\item[(V2)] if $b-a=p$ (and no other multiples of $p$ occur): $v(\Rk(a,b))\ge-2$ and $\operatorname{lead}_{-2}\Rk(a,b)=\frac{\mathsf L(a)-\mathsf L(b)}4$.
\end{itemize}
Two further elementary facts: the \emph{block-scaling lemma} (if $D=\operatorname{diag}(p^{d_i})$ and $\tilde A=DAD\in M(\mathcal O_p)$ has $\tilde A\bmod p$ invertible over $\F_p(X)$, then $A^{-1}=D\tilde A^{-1}D$ with $\tilde A^{-1}\in M(\mathcal O_p)$, because $\det\tilde A$ is a unit) and the \emph{partner decomposition}: for a zone row $z$ and a zone column $y$, $z\ne y$, whose partners $c_z=z-p$, $c_y=y-p$ lie in $C_0$ and are distinct, put
$\ell_z(c)=\Rk(z,c)-\Rk(c_z,c)-X\mathbf 1[c=c_z]$ ($c\in C_0$), so that $\Rk[z,C_0]=A[c_z,C_0]+\ell_z$ and likewise $\Rk[C_0,y]=A[C_0,c_y]+\ell_y^{\top}$. Since $A[c_z,\cdot]A^{-1}=e_{c_z}^\top$,
\begin{equation}\label{eq:hpartner}
 T_{zy}=\Delta^2\Rk(z,y)-\ell_zA^{-1}\ell_y^{\top},\qquad \Delta^2\Rk(z,y):=\Rk(z,y)-\Rk(z,c_y)-\Rk(c_z,y)+\Rk(c_z,c_y).
\end{equation}
(The variable $X$ cancels in $\Delta^2\Rk$ because $c_z\ne c_y$.) With $\delta_z:=F(z)-F(c_z)$ one computes, for $c\in C_0\setminus\{c_z\}$ with $p\nmid c-z$,
\begin{equation}\label{eq:hdefect}
 \ell_z(c)=\frac{p\,\Rk(c_z,c)}{c-z}+\frac{\delta_z}{4(c-z)},\qquad \ell_z(c_z)=-\frac{\delta_z}{4p}-\Psi(c_z)-X.
\end{equation}
For the diagonal entry $(b,b)$ (when $e=1$), with $c=c_b\in C_0$ and the same decomposition $\Rk[b,C_0]=A[c,C_0]+\ell_b$, $\Rk[C_0,b]=A[C_0,c]+\ell_b^\top$ (by symmetry of $\Rk$): $\Rk[b,C_0]A^{-1}\Rk[C_0,b]=A(c,c)+2\ell_b(c)+\ell_bA^{-1}\ell_b^\top$ with $A(c,c)=\Psi(c)+X$ and $\ell_b(c)=\Rk(b,c)-\Psi(c)-X$, so
$T_{bb}=\Psi(b)+X-\Rk(b,c)\cdot2+\bigl(\Psi(c)+X\bigr)-\ell_bA^{-1}\ell_b^\top$.

\begin{theorem}[Lemma E]\label{thm:hE}
Let $\mathcal G(r,m,e)$ be admissible, $r$ large, $p\ge11$ a prime, $p\ne139$, and let $v$ be the Gauss valuation at $p$.
\begin{enumerate}[label=(\alph*),leftmargin=*]
\item If $p\equiv3\pmod 4$, $2p>r+2m$ and $5p+1<4r$, then $v(\det\Fu(X))\ge-2r-2e$.
\item If $p\equiv1\pmod 4$, $p<r$, $7p+1>4r+8m$ (for $e=0$ it suffices that $7p+1\ge4r+8m$), and $K_p:=\frac{3p+1}4$, then $v(\det\Fu(X))\ge-2(2r-p-K_p)-3m-4e$.
\end{enumerate}
\end{theorem}
(For $m=r/8$, $e=0$ the conditions are exactly $5r<8p$, $5p+1<4r$, resp.\ $5r\le7p+1$, $p<r$.)

\begin{proof}
Throughout, $V1$, $V2$ are used with all indices $<2p$ (this holds: all indices are $\le r+2m<2p$ in both cases), so that $b-a$ is a $p$-unit unless $|b-a|=p$.

\textbf{Case (a), $p\equiv3\pmod 4$.} Put $k_1=\frac{p+1}4$ (so $4k_1-1=p$), $k_2=\frac{5p+1}4=p+k_1$ ($4k_2-1=5p$). The singular indices below $r+2m+1$ are $k_1,p,k_2$: indeed $2p>r+2m$ and $\frac{9p+1}4>2p$. Also $k_2<r$ by $5p+1<4r$, and $p<r$. So $[0,r)$ splits into blocks $B_0=[0,k_1)$, $B_1=[k_1,p)$, $B_2=[p,k_2)$, $B_3=[k_2,r)$ and the zone $[r,r+2m]$ lies in the block above $k_2$, with levels and diagonal leads
\begin{center}
\begin{tabular}{c|cccc}
block & $B_0$ & $B_1$ & $B_2$ & $B_3\ (\text{and zone})$\\ \hline
$\mathsf L$ & $0$ & $-1$ & $-\frac34$ & $-\frac{19}{20}$\\
$\pi$ & $0$ & $-1$ & $-\frac{15}{16}$ & $-\frac{391}{400}$
\end{tabular}
\end{center}
(jumps $-\frac1j$, $+\frac14$, $-\frac1{5}$ for $\mathsf L$ and $-\frac1{j^2},+\frac1{16},-\frac1{25}$ for $\pi$). The partner map $a\mapsto a+p$ within $C_0$ sends $B_0$ onto $B_2$ and $[k_1,r-p)$ onto $B_3$ (here $r-p\ge k_1$ by $5p+1\le4r$); the set $\mathcal U=[r-p,p)\subset B_1$ is unpaired (for $a\in\mathcal U$: $a+p\ge r$ and $a-p<0$). Each zone node $z\in[r,r+2m]$ has partner $c_z=z-p\in[r-p,r-p+2m]\subset\mathcal U$ (as $r+2m-p<p$ by $2p>r+2m$).

\emph{Leading form of $A$.} By V1, V2: $p^2A\in M(\mathcal O_p)$, and $p^2A\bmod p$ is the matrix $M_0$ with entries $\pi(a)$ on the diagonal, $\frac{\mathsf L(a)-\mathsf L(a+p)}4$ at $(a,a+p)$ and $(a+p,a)$, and $0$ otherwise. After a permutation $M_0$ is block diagonal: $1\times1$ blocks $(-1)$ on $\mathcal U$; $2\times2$ blocks $\begin{psmallmatrix}0&3/16\\3/16&-15/16\end{psmallmatrix}$ (determinant $-\frac9{256}$) for $a\in B_0$; and $\begin{psmallmatrix}-1&-1/80\\-1/80&-391/400\end{psmallmatrix}$ (determinant $\frac{6255}{6400}$, $6255=3^2\cdot5\cdot139$) for $a\in[k_1,r-p)$. Hence $M_0$ is invertible for $p\notin\{3,5,139\}$, and $N:=(p^2A)^{-1}\in M(\mathcal O_p)$ satisfies $A^{-1}=p^2N$, $N\equiv M_0^{-1}\pmod p$. Let $P$ be the set of paired indices; block inversion of $p^2A=\begin{psmallmatrix}\tilde A_{\mathcal U\mathcal U}&\tilde A_{\mathcal UP}\\ \tilde A_{P\mathcal U}&\tilde A_{PP}\end{psmallmatrix}$ gives (the entries between distinct unpaired indices, and between $\mathcal U$ and $P$, are $p^2A(c,c')$ with $v(A(c,c'))\ge-1$, i.e. lie in $p\mathcal O_p$): $N_{\mathcal U\mathcal U}=S^{-1}$ with $S=-I+pS_1$, $S_1(c,c')\equiv pA(c,c')=\operatorname{lead}_{-1}A(c,c')=\frac{\mathsf L(c)-\mathsf L(c')}{4(c'-c)}=0\pmod p$ for $c\ne c'$ in $\mathcal U$ (same level), and $N_{\mathcal UP}\in p\mathcal O_p$. Therefore
\begin{equation}\label{eq:hNform}
 N(c,c)\equiv-1\ (p),\qquad N(c,c')\equiv0\ (p^2)\ \text{for }c\ne c'\in\mathcal U,\qquad N(c,c')\in p\mathcal O_p\ (c\in\mathcal U,c'\ne c).
\end{equation}

\emph{Entries of $T$.} Let $z\in X''$, $y\in Y''$, $z\ne y$, $d=y-z$ (a nonzero integer, $|d|\le2m<p$, a unit). Then $c_z\ne c_y$, both in $\mathcal U$. The first term of \eqref{eq:hpartner}: $z,y$ lie in one block, hence $v(\Rk(z,y))\ge0$ ($\mathsf L(z)=\mathsf L(y)$, $d$ a unit); likewise $v(\Rk(c_z,c_y))\ge0$. With $\sigma:=\mathsf L(z)-\mathsf L(c_z)$ (the same for all zone nodes), $\operatorname{lead}_{-1}\Rk(z,c_y)=\frac{\sigma}{4(d-p)}\equiv\frac\sigma{4d}$ and $\operatorname{lead}_{-1}\Rk(c_z,y)=-\frac{\sigma}{4(d+p)}\equiv-\frac\sigma{4d}$; these cancel mod $p$, so $v(\Delta^2\Rk)\ge0$.
For the second term, by \eqref{eq:hdefect}: for $c\ne c_z$ we have $c-z\equiv c-c_z\not\equiv0$, $p\ell_z(c)\equiv\frac\sigma{4(c-c_z)}\pmod p$ (since $p^2\Rk(c_z,c)\in p\mathcal O_p$), and $p^2\ell_z(c_z)\equiv-\frac\sigma4-\pi(c_z)=:-\beta\pmod p$, the same for all $z$ (the signs: $\Rk(z,c_z)=-\delta_z/(4p)$). Put $\lambda_z=p^2\ell_z\in\mathcal O_p^{C_0}$: $\lambda_z(c_z)\equiv-\beta$, $\lambda_z(c)=pu_z(c)$ with $u_z(c)\equiv\frac\sigma{4(c-c_z)}$ ($c\neq c_z$). Then $\ell_zA^{-1}\ell_y^\top=p^{-2}\lambda_zN\lambda_y^\top$ and, by \eqref{eq:hNform}, modulo $p^2$:
the $(c_z,c_y)$ term is $\equiv0$ because $N(c_z,c_y)\equiv0\ (p^2)$; the terms with $c\ne c_z$ and $c'\neq c_y$ are $O(p^2)$; the terms $(c_z,c')$, $c'\ne c_z,c_y$ and $(c,c_y)$ are $O(p^2)$ by $N(c_z,c')\in p\mathcal O_p$; the remaining ones are $(c_z,c_z)$ and $(c_y,c_y)$:
\begin{multline*}
 \lambda_zN\lambda_y^\top\equiv(-\beta)(-1)\,p\,u_y(c_z)+p\,u_z(c_y)(-1)(-\beta)\\=p\beta\Bigl[\frac\sigma{4(c_z-c_y)}+\frac\sigma{4(c_y-c_z)}\Bigr]\equiv0\pmod{p^2}.
\end{multline*}
So $v(\ell_zA^{-1}\ell_y^\top)\ge0$ and $v(T_{zy})\ge0$ for all $z\neq y$.

\emph{The diagonal entry $T_{bb}$ ($e=1$).} Every term in the formula for $T_{bb}$ has $v\ge-2$: $\Psi(b),\Psi(c)$ by $v(\Psi)\ge-2$, $\Rk(b,c_b)$ by V2, and $\ell_bA^{-1}\ell_b^\top=p^{-2}\lambda_bN\lambda_b^\top$ with $\lambda_b\in\mathcal O_p^{C_0}$, $N\in M(\mathcal O_p)$. So $v(T_{bb})\ge-2$.

\emph{Conclusion.} For $e=0$: $v(\det T)\ge0$ by \eqref{eq:hredE}; for $e=1$ every permutation term of $\det T$ either contains $T_{bb}$ (then $v\ge-2$) or $T_{by}T_{zb}$ with all other factors $\ge0$: $v(\det T)\ge-2$. Finally $\tau(W_A)\le2r$ since all weights are $\le2$. This proves (a).

\textbf{Case (b), $p\equiv1\pmod 4$.} Now $4^{-1}\equiv K_p$, singular indices $K_p,p,\frac{7p+1}4=K_p+p$ ($4k-1=7p$) and $2p$. By the hypotheses $K_p<p<r$, $K_p+p>r+2m$ (since $7p+1>4r+8m$; in the case $e=0$, $=$ allowed, because the index $K_p+p=r+2m$ is then just above the zone $[r,r+2m)$) and $2p>r+2m$ (as $p>\frac57r$). Blocks $B_0=[0,K_p)$, $B_1=[K_p,p)$, $B_2=[p,r)$ and the zone, with $\mathsf L=0,-\frac13,-\frac1{12}$ and $\pi=0,-\frac19,-\frac19+\frac1{16}=-\frac7{144}$. Partners: $P_0=[0,r-p)\subset B_0$ ($r-p<K_p$) is paired with $B_2$; $B_1$ is unpaired ($K_p+p>r$); $\mathcal U=[r-p,K_p)$ is unpaired. Let $P=P_0\cup B_1\cup B_2$, $|P|=2r-p-K_p$. The partner of a zone node $z$ is $c_z=z-p\in[r-p,r-p+2m]\subset\mathcal U$, since $r-p+2m<K_p$ is the hypothesis $4r+8m<7p+1$.

\emph{Scaling and the scalar Schur complement.} Let $D=\operatorname{diag}(1\text{ on }\mathcal U,\ p\text{ on }P)$ and $\tilde A=DAD$. By V1, V2: $\tilde A\in M(\mathcal O_p)$ ($\mathcal U\times\mathcal U$ entries have $v\ge0$; $\mathcal U\times P$ entries have $v\ge-1$ since no pair in $\mathcal U\times P$ differs by $p$; $P\times P$ entries have $v\ge-2$). Modulo $p$: $\tilde A_{PP}\equiv M_{0P}$, block diagonal with $1\times1$ blocks $-\frac19$ on $B_1$ and $2\times2$ blocks $\begin{psmallmatrix}0&e\\e&\pi_2\end{psmallmatrix}$, $e=\frac1{48}$, $\pi_2=-\frac7{144}$, on the pairs $(a,a+p)$; $\tilde A_{\mathcal UP}\equiv V$ with $V_{u,c}=\frac1{12(c-u)}$ ($c\in B_1$), $\frac1{48(c-u)}$ ($c\in B_2$), $0$ ($c\in P_0$); and $\tilde A_{\mathcal U\mathcal U}\equiv\bar A_{\mathcal U\mathcal U}$ (the reduction of $\Rk+XI$). We claim
\begin{equation}\label{eq:hscalar}
 \bar S:=\bar A_{\mathcal U\mathcal U}-VM_{0P}^{-1}V^\top=s_0I_{\mathcal U},\qquad s_0=X-\tfrac1{16}H^{(2)}_{K_p-1}\in\F_p(X)^\times .
\end{equation}
\emph{Proof of \eqref{eq:hscalar}.} The inverse of $\begin{psmallmatrix}0&e\\e&\pi_2\end{psmallmatrix}$ has zero $(B_2,B_2)$ entry, and $V$ vanishes on $P_0$; so only the $1\times1$ blocks contribute, and $\bar S=\bar A_{\mathcal U\mathcal U}+9\,V_{B_1}V_{B_1}^\top=\bar A_{\mathcal U\mathcal U}+\frac1{16}WW^\top$, $W_{u,c}=\frac1{c-u}$, $K_p\le c\le p-1$. For $n<K_p$, using $4K_p\equiv1$ so $\frac4{4k-1}\equiv\frac1{k-K_p}$:
$4F(n)\equiv\varphi(n):=H_n+H_{K_p-1}-H_{K_p-1-n}$ and $16\Psi(n)\equiv H^{(2)}_n-H^{(2)}_{K_p-1}+H^{(2)}_{K_p-1-n}$. Hence $16\bar A_{uw}=\frac{\varphi(u)-\varphi(w)}{w-u}$ for $u\ne w$ in $\mathcal U$. Partial fractions: $\sum_{c=K_p}^{p-1}\frac1{(c-u)(c-w)}=\frac{[H_{p-1-u}-H_{K_p-1-u}]-[H_{p-1-w}-H_{K_p-1-w}]}{u-w}$. Adding, $16\bar S_{uw}=\frac{[H_{p-1-u}-H_u]-[H_{p-1-w}-H_w]}{u-w}\equiv0$ since $H_{p-1-n}\equiv H_n\pmod p$. On the diagonal, $16\bar S_{uu}=16X+H^{(2)}_u-H^{(2)}_{K_p-1}+H^{(2)}_{K_p-1-u}+(H^{(2)}_{p-1-u}-H^{(2)}_{K_p-1-u})\equiv16X-H^{(2)}_{K_p-1}$, using $H^{(2)}_{p-1-n}\equiv-H^{(2)}_n$ ($p\ge5$). $\square$

Hence $\tilde A\bmod p$ is invertible over $\F_p(X)$ (its determinant is $\det M_{0P}\,s_0^{|\mathcal U|}\neq0$: $-e^2$ and $-\frac19$ are nonzero for $p\ne3$), so $A^{-1}=D\tilde A^{-1}D$ and $(\tilde A^{-1})_{\mathcal U\mathcal U}\equiv s_0^{-1}I$. Moreover $v(\det A)\ge-2|P|$ (Lemma~\ref{lem:htrop} with $u=v=1$ on $P$ and $0$ on $\mathcal U$; $W\le2$ on $P\times P$, $\le1$ on $\mathcal U\times P$ and $0$ on $\mathcal U\times\mathcal U$).

\emph{Entries of $T$.} For $z\ne y$ in the zone: $v(\Delta^2\Rk)\ge-1$ (V1; all four terms). For the defect term write $\ell_zA^{-1}\ell_y^\top=p^{-4}\mu_z\tilde A^{-1}\mu_y^\top$ with $\mu_z=p^2\ell_zD$. By \eqref{eq:hdefect} and V1, $\ell_z(c)$ has $v\ge-1$ for $c\ne c_z$ (no $c\in C_0$ differs from $c_z\in\mathcal U$ by $p$), so $\mu_z\in\mathcal O_p^{C_0}$ and $\mu_z\equiv-\beta e_{c_z}\pmod p$ with $\beta=\sigma/4$ (here $\pi(c_z)=0$ on $B_0$). Hence $\mu_z\tilde A^{-1}\mu_y^\top\equiv\beta^2(\tilde A^{-1})(c_z,c_y)\equiv\beta^2s_0^{-1}\mathbf 1[c_z=c_y]=0\pmod p$ for $z\ne y$, i.e.\ $v\ge1$ and $v(\ell_zA^{-1}\ell_y^\top)\ge-3$. So $v(T_{zy})\ge-3$ for $z\ne y$. For $T_{bb}$ ($e=1$): the same expression with $c_z=c_y=c_b$ gives $v(\ell_bA^{-1}\ell_b^\top)\ge-4$, and the other terms have $v\ge-2$; so $v(T_{bb})\ge-4$.

\emph{Conclusion.} $v(\det T)\ge-3m$ for $e=0$; for $e=1$ every term of the Leibniz expansion has $v\ge-3m-4$ (a term containing $T_{bb}$ has $m$ further factors $\ge-3$; others have $m+1$ factors $\ge-3$). With $v(\det A)\ge-2|P|$ we obtain (b).
\end{proof}

\begin{remark}[formal variable]\label{rem:hformal}
The proof uses that $s_0=X-\frac1{16}H^{(2)}_{K_p-1}$ is a unit of $\mathcal O_p$, which holds because $X$ is an indeterminate: the residue of $s_0$ in $\F_p(X)$ is nonzero for every $p$. If $X$ is specialised to $x_0\in\Q$ with $16x_0\equiv H^{(2)}_{K_p-1}\pmod p$, then $\bar S\equiv0$ and the bound fails (we have observed $\min v(T_{zy})=-5,-6$ instead of $-3$ at $r=64$, $p=53,61$). Likewise (a) uses $p\nmid 6255$. The statement $v_p(\den_F)\le\dots$ is about the coefficients of $\det\Fu(X)\in\Q[X]$ and therefore about the Gauss valuation; in the proof of Theorem~\ref{thm:hmeasure} only $\den_F$ is used, never the value of a rank-saving bound at a specialisation.
\end{remark}

\begin{remark}[member~2: where the extra node enters]\label{rem:hmember2}
The node $b=r+2m$ enters the proof of Theorem~\ref{thm:hE} only through the following points, each of which uses exactly the hypothesis stated for $e=1$. (1) $b$ is a zone node with a partner: in (a), $b<2p$ (the hypothesis $2p>r+2m$) puts $b$ in the top block and $c_b=b-p\in\mathcal U$; in (b), the \emph{strict} inequality $7p+1>4r+8m$ gives $b<K_p+p$ (so $b$ lies in the block of the zone, with the same level $\mathsf L$) and $c_b<K_p$, $c_b\in\mathcal U$. Hence the off-diagonal entries of $T$ in row and column $b$ obey the same bounds ($\ge0$ in (a), $\ge-3$ in (b)) as those between ordinary zone nodes, by the same computation (it uses only that $z\ne y$ are zone nodes with distinct unpaired partners at a unit distance $<p$). (2) $b$ is the only common index of $X''$ and $Y''$, so $T_{bb}$ is the only diagonal entry of $T$; its bound ($\ge-2$ in (a), $\ge-4$ in (b)) uses $\lambda_b\in\mathcal O_p^{C_0}$, which holds because no $c\in C_0$ other than $c_b$ differs from $b$ by $p$ and $c_b$ is unpaired. (3) The Leibniz expansion of the $(m+1)\times(m+1)$ matrix $T$ then gives the extra $-2e$ resp.\ $-4e$. Lemma~E is applied (Corollary~\ref{cor:hsave}) only with the tropical bound $B_{m_0}(p)$ of member~1 as reference, and the extension paragraph in the proof of Proposition~\ref{prop:htropical} bounds the tropical value of member~2 by $B_{m_0}(q)+12$; so the extra node changes each per-prime bound by $O(1)$, and the total by $O(\vartheta(4N'))=O(r)=o(r^2)$, both in the tropical law and in the saving.
\emph{Exact checks.} (i) The script \texttt{member2\_content.py} gives $\tau_p-v_p(\den_F)=m$ (exact Hungarian $\tau_p$ for member~2) at every clean prime for $r=48,64,80,96$ ($m=r/8$). (ii) An independent script, \texttt{check\_hE\_member2.py} (own interpolation nodes $k+\frac13$, comparison with the bounds of Theorem~\ref{thm:hE} themselves), gives $v_p(\den_F)$ equal to the bound of Theorem~\ref{thm:hE} (slack $0$) in all $27$ cases $(r,m,e,p)$ with $r\in\{48,64,80\}$, $m\in\{\frac r8-1,\frac r8,\frac r8+1\}$, $e\in\{0,1\}$ and $p$ satisfying the hypotheses of (a) or (b); no violation.
\end{remark}

\begin{corollary}[content saving]\label{cor:hsave}
Let $\mathcal G(r,m,e)$ be admissible, $r$ large, and let $p\ne139$ be a prime with
\[
 p\equiv3\ (4),\ \ \tfrac58r+2<p<\tfrac45r-2,\qquad\text{or}\qquad p\equiv1\ (4),\ \ \tfrac57r+2<p<r-2.
\]
Then the hypotheses of Theorem~\ref{thm:hE} hold and $v_p(\den_F)\le B_{m_0}(p)-m_0+C_1$ with an absolute constant $C_1$ ($m_0=\lfloor r/8\rfloor$, $B_{m_0}$ as in Proposition~\ref{prop:hdual}). In these ranges $B_{m_0}(p)=2r+m_0$ (for $p\equiv3$) and $B_{m_0}(p)=2(2r-p-K_p)+4m_0$ (for $p\equiv1$).
\end{corollary}
\begin{proof}
The hypotheses: $2p>\frac54r+4>r+2m$ since $m\le\frac r8+1$; $5p+1<4r$; $7p+1>5r+15>4r+8m$ since $8m\le r+8$. For the values of $B$: for $p\equiv3$ and $\theta=p/r<\frac45$ one has $|U|=0$, $L=0$, $E=m_0$, $\delta_A=2r$; for $p\equiv1$ and $\theta>\frac57$, $|U|=\frac{7\theta}4r-r+O(1)>2m_0$, so $L=2m_0$, $E=\min(2m_0+L,m_0+L+|U|)=4m_0$ and $\delta_A=2(r-K_p)+2(r-p)$. Now Theorem~\ref{thm:hE} gives $v_p(\den_F)\le2r+2\le B-m_0+2$ in the first case and $\le2(2r-p-K_p)+3m+4\le B-m_0+3(m-m_0)+4\le B-m_0+10$ in the second (by \eqref{eq:hden}).
\end{proof}

\subsection{Proof of Proposition~\ref{prop:hC}}
By \eqref{eq:hden}, \eqref{eq:hmaj} and Proposition~\ref{prop:htropical},
$\log\den_F=\sum_pv_p(\den_F)\log p\le T(r)-\sum_{p\text{ clean}}(\ldots)$, where for every prime $p$ in the ranges of Corollary~\ref{cor:hsave} we replace the tropical bound $\sum_e\tau(w^{(p^e)})\le B_{m_0}(p)+12$ (note $p^2>4N'$ for these $p$, so only level $1$ occurs) by the bound of Corollary~\ref{cor:hsave}, which is smaller by $m_0-C_1'$. The primes involved (excluding $139$, which lies in the range only for $r<222$) have
\[
 \sum_{\substack{p\equiv3\\ \frac58r<p<\frac45r}}\log p+\sum_{\substack{p\equiv1\\ \frac57r<p<r}}\log p=\frac r2\Bigl[\bigl(\tfrac45-\tfrac58\bigr)+\bigl(1-\tfrac57\bigr)\Bigr]+o(r)=\frac{129}{560}r+o(r)
\]
(PNT in progressions, $\vartheta(x;4,a)=x/2+o(x)$). Therefore the total saving is $m_0\cdot\frac{129}{560}r+O(r)=\frac{129}{4480}r^2+o(r^2)$, and
\begin{gather*}
 \log\den_F\le\Bigl(\Lambda_{\rm trop}-\frac{129}{4480}+o(1)\Bigr)r^2,\\
 \Lambda_{\rm trop}-\frac{129}{4480}=3.7499347-0.0287946=3.7211401\ (\le 3.7211401).
\end{gather*}
This proves Proposition~\ref{prop:hC}. (More precisely $\Lambda_{\rm trop}-\frac{129}{4480}=3.72114002$.) \qed

\begin{remark}[range of validity; what is not claimed]
The saving $\frac{\alpha}2(\frac{51}{70}-\frac{15\alpha}7)$ for general $\alpha$ is valid only for $\frac18\le\alpha\le\frac16$; we use it only at $\alpha=\frac18$, where it is $\frac{129}{4480}$. Primes $p<\frac58r$ (clean primes with additional singular indices) and the non-clean primes in the band (where the saving is $m$ minus the number of zone nodes whose window $(z-p,z]$ meets an extra singular index) probably give a further saving, which is not proved and not used. Any finite-$r$ table of the saving must not count $p=139$ (for which $\det M_0\equiv0$), which is why the exact finite-$r$ values listed in our computations are given only as numerical evidence.
\end{remark}

\section{Decay and the window for the hybrid frame}\label{sec:hdecay}

This section proves Proposition~\ref{prop:hD}. The proof is that of Theorem~\ref{thm:decay} and of Theorem~\ref{thm:N} (node removal) with two changes: Andr\'eief's identity is used with two different families of functions (rows and columns), and the node set is no longer a segment of integers, so the exact factorial identities are replaced by elementary cell comparisons. Only \emph{common} nodes are removed. Throughout, $r$ is the scale of the geometry, $\alpha_m=m/r$, $L=2m/r$, $N_{\rm tot}=|C|+m$, and $\mathcal T_t$, $\mathcal G=\mathcal G(r,m,e)$ are as in \S\ref{sec:hframe}; the kernel $k_\rho$, Lemmas~\ref{lem:K0}--\ref{lem:K} and Corollary~\ref{cor:Kp} of \S\ref{ssec:krho} are used unchanged.

\subsection{Contour identity and two-family Andr\'eief}
Fix $r$ and put $N_0=\lfloor r_0r\rfloor$, $r_0=\frac7{50}$, $x^*=-N_0-\frac38$, and let $\Gamma$, $\mathcal A$, $m$ (the complex measure $m=-\sum_{\mu\in\mathcal A}c_\mu\delta_{-\mu}+\frac1{2\pi i}E(z)dz|_\Gamma$) be as in \S\ref{sec:decay} with $h$ replaced by $r$. The proof of Proposition~\ref{prop:contour} does not use the nodes except that they lie in $[0,\infty)$, hence to the right of $\Gamma$ and away from the poles of $F$; therefore for \emph{all} integers $a,b\ge0$
\begin{equation}\label{eq:hcontour}
 \frac{F(b)-F(a)}{b-a}=\int v_av_b\,dm\ (a\ne b),\qquad F'(a)=\int v_a^2\,dm,\qquad v_a(z)=\frac1{z-a}.
\end{equation}
For a set $S\subseteq C$ of common nodes put $\mathcal R_S=S\sqcup X'$, $\mathcal Q_S=S\sqcup Y'$, $n=|S|+m$, and $M_S=\frac1{16}L_F[\mathcal R_S,\mathcal Q_S]$. Since the diagonal of $\mathcal T_t$ is $\tfrac1{16}F'(c)+t$ at $c\in C$, multilinearity gives
\begin{equation}\label{eq:hexpand}
 \det\mathcal T_t=\sum_{S\subseteq C}t^{|C|-|S|}\det M_S .
\end{equation}
\begin{proposition}\label{prop:hAndreief}
\[\det L_F[\mathcal R_S,\mathcal Q_S]=\frac1{n!}\int\det[v_a(z_i)]_{i\le n,\,a\in\mathcal R_S}\ \det[v_b(z_i)]_{i\le n,\,b\in\mathcal Q_S}\prod_idm(z_i),\]
hence
\begin{gather*}
 |\det L_F[\mathcal R_S,\mathcal Q_S]|\le\frac1{n!}\int|\mathsf C_{\mathcal R_S}(z)|\,|\mathsf C_{\mathcal Q_S}(z)|\prod_id|m|(z_i),\\ \mathsf C_{\mathcal F}(z)=\frac{\prod_{i<i'}(z_i-z_{i'})\prod_{a<a'\in\mathcal F}(a'-a)}{\prod_{i,a\in\mathcal F}(z_i-a)}\ (\text{up to sign}).
\end{gather*}
\end{proposition}
\begin{proof}
Andr\'eief's identity $\det\bigl[\int f_ag_b\,d\mu\bigr]_{a,b}=\frac1{n!}\int\det[f_a(z_i)]\det[g_b(z_i)]\prod d\mu(z_i)$ holds for two families of $n$ functions, a finite complex measure $\mu$ and bounded $f_a,g_b$ (here $|v_a|\le\frac83$ on the support of $m$, since $\operatorname{Re}z\le-\frac38$); apply it with \eqref{eq:hcontour}. The Cauchy determinant evaluates $\det[v_a(z_i)]$.
\end{proof}

\subsection{The potentials and the certified numbers}\label{ssec:hcert}
Let $\alpha\in[0,\frac14]$ and $\nu_\alpha=\mathrm{Leb}|_{[0,1]}+\frac12\mathrm{Leb}|_{[1,1+2\alpha]}$ (mass $1+\alpha$). Put $U^{\nu_\alpha}(s)=-\int\log|s-u|\,d\nu_\alpha(u)$ and $I(\nu_\alpha)=\iint-\log|u-u'|\,d\nu_\alpha d\nu_\alpha$. Closed forms: for $s=x+iy$, $x<0$, $y\ge0$, with $\Phi(q)=a_q(\frac12\log(a_q^2+y^2)-1)-y\arctan(y/a_q)$, $a_q=q-x$: $U^{\nu_\alpha}(s)=-[\Phi(1)-\Phi(0)]-\frac12[\Phi(1+2\alpha)-\Phi(1)]$; on the real line, with $G(x)=x\log|x|-x$, $U^{\nu_\alpha}(t)=1-t\log|t|-(1-t)\log|1-t|-\frac12[G(1+2\alpha-t)-G(1-t)]$, so that $U^{\nu_\alpha\prime}(t)=\frac12\log|(1-t)(1+2\alpha-t)|-\log|t|$ is decreasing on $(0,1)$ (hence $U^{\nu_\alpha}$ is concave there), positive on $[-r_0,0)$, and negative on $(1,\infty)$; and $I(\nu_\alpha)=\frac32+I_{12}+\alpha^2(\frac32-\log2\alpha)$, with $I_{12}=-[\mathsf F(1+2\alpha)-\mathsf F(1)-\mathsf F(2\alpha)+\mathsf F(0)]$, $\mathsf F(x)=\frac{x^2}2\log|x|-\frac{3x^2}4$. At $\alpha=\frac18$: $I(\nu_{1/8})=1.70244575114$. (These were checked against quadrature to $15$ digits.)

For a probability measure $\sigma=\sum_km_k\delta_{c_k}$ and $\rho$, $U^\sigma_\rho$, $I_\rho(\sigma)$ are as in \S\ref{ssec:cert}. Fix $\kappa=10^{-3}$ and put
\[
 g_\alpha(s)=(1+\alpha)U^\sigma_\rho(s)-U^{\nu_\alpha}(s)+(1-\kappa)\pi|\operatorname{Im}s|,\qquad \partial_\alpha g_\alpha(s)=U^\sigma_\rho(s)+\log|s-(1+2\alpha)| .
\]
\paragraph{The certificate (computer-assisted).} Let $\sigma=\sigma_{\rm hyb}$ be the measure of the file \texttt{sig\_h500.json} (dyadic data in the format of \texttt{sigma300.json}: coordinates $X/2^{24}$, $Y/2^{24}$, masses $n_k/2^{32}$, $\sum m_k=1$ exactly, conjugation-symmetric, $\rho=1/2048$; the file is named after its design parameter $K=500$), $\alpha=\frac18$, $L_0=\frac14$, $u_{\max}=\frac32$, $\omega=2^{-30}$. The script \texttt{hcert.py} (Arb, $96$ bits) certifies, with $g=g_{1/8}$:
\begin{enumerate}[label=(\textbf{\arabic*}),leftmargin=*]
\item[(R)] $g\ge m'$ on all boxes $[-r_0-u_b-\omega,-r_0-u_a+\omega]\times[3u_a,3u_b]$ covering $u\in[0,u_{\max}]$ (rays $-r_0+u(-1\pm3i)$), and the tail bound $T(u_{\max})\ge m'$, where for $u\ge u_{\max}$: $|z|\le r_0+\omega+u\sqrt{10}$, $U^\sigma_\rho(z)\ge-\frac12\log((|z|+R_\sigma)^2+\rho^2)$, $U^{\nu}(z)\le-(1+\alpha)\log(|z|-1-2\alpha)$ (all of $\nu$ lies in $[0,1+2\alpha]$ and has mass $1+\alpha$) and the field term is $\ge(1-\kappa)3\pi u$; $T$ is increasing in $u$ (derivative $\ge9-1$).
\item[(S)] a partition of $[-r_0,0]$ into cells $[a_\ell,b_\ell]$ with $g\ge\gamma_\ell$ on them, and $S_{\rm cert}:=2\sum_\ell(b_\ell-a_\ell)\min(\gamma_\ell-m',0)$;
\item[(L)] a cover of $[0,1+2\alpha]$ by cells on which $g\ge g_{\min}$.
\end{enumerate}
On a box or cell the lower bound for $U^\sigma_\rho$ is $-\frac12\sum_km_k\log(d_{x,k}^2+d_{y,k}^2+\rho^2)$ ($d_{x,k},d_{y,k}$ the largest horizontal and vertical distances to $c_k$), $U^\nu$ is bounded above by ball arithmetic on the closed form on ray boxes, by its right endpoint value on segment cells ($U^\nu$ is increasing on $(-\infty,0]$), by its left endpoint value on cells in $[1,\infty)$ (decreasing) and on cells of $[0,1]$ by the endpoint values or the two tangent bounds (concavity). Put
\[
 D_F:=2\bigl((1+\alpha)m'+S_{\rm cert}\bigr)-(1+\alpha)^2I_\rho(\sigma)+I(\nu_\alpha),\qquad c_{\rm hyb}:=2(m'-g_{\min}).
\]
\begin{center}
\begin{tabular}{p{0.52\textwidth}p{0.3\textwidth}l}
\toprule
quantity & certified value & cells\\
\midrule
$I_\rho(\sigma)$ (exact finite sum) & $2.00548597469\ (\pm6\cdot10^{-26})$ & finite sum\\
$I(\nu_{1/8})$ & $1.70244575114$ & --\\
$m'$: $\inf g$ on the rays, $u\in[0,\frac32]$ & $\ge2.2185001686$ & $11{,}681$ boxes\\
tail $T(\frac32)$ & $\ge13.595$ & --\\
$S_{\rm cert}$ & $\ge-0.1248042185$ & $14{,}000$ cells\\
$g_{\min}$ on $[0,1.25]$ (attained near $t=0.8284$) & $\ge-1.6085322035$ & $4{,}512$ cells\\
\midrule
$D_F$ & $\ge3.9062695066$ & \\
$c_{\rm hyb}$ & $\le7.6540647441$ & \\
\bottomrule
\end{tabular}
\end{center}
Negative control: with the plain certificate $\sigma_{300}$ and $\alpha=0$ the same script returns $D\ge3.111379042$, $c\le6.871799888$, reproducing the numbers of \S\ref{ssec:cert} ($3.1113790$, $6.8717994$). The copy of \texttt{hcert.py} in \texttt{scripts-hybrid/} has been re-run (Python~3.13 with \texttt{python-flint} and \texttt{numpy}, the libraries supplied through \texttt{PYLIBS}): it reproduces $D_F\ge3.906269506610271$ and $c_{\rm hyb}\le7.654064744147046$ exactly, in $8$ seconds; the original was also re-run by an audit agent, with identical output. The script, the data \texttt{sig\_h500.json}, the exact run instructions and the expected output are in \texttt{scripts-hybrid/} (\S\ref{ssec:hrepro}). The certificate is valid for $r\ge2^{30}$ and does not depend on $r$.

\begin{remark}[$c_{\rm hyb}$ is below the true first-order slope]
The true profile $-r^{-2}\log|\sum_{|S^c|=j}\det M_S|$ has initial slope $\approx7.687$ in $\beta=j/r$ (exact computations at $r\le128$ extrapolate to $7.687,7.682,7.677,7.673,7.654$ for $j=1,2,3,4,8$), larger than the certified $c_{\rm hyb}=7.654$. There is no contradiction: Theorem~\ref{thm:hN} certifies the \emph{line} $D_F-\beta c_{\rm hyb}$, which lies below the true profile at every $\beta$ (the margins at $\beta=0,\frac1{16},\frac18,\frac14,\frac12,1$ are $0.008,0.026,0.081,0.29,1.06,3.87$; the true profile is strongly convex). The window argument uses only that line (Corollary~\ref{cor:hwindow}).
\end{remark}

\subsection{The pointwise inequality}
Let $\mathcal G=\mathcal G(r,m,e)$ be admissible, $S\subseteq C$, $S^c=C\setminus S$, $j=|S^c|$, $\beta=j/r$. Let $\mathcal R,\mathcal Q$ be the full row and column node sets ($|\mathcal R|=|\mathcal Q|=N:=r+m+e$), $n=|S|+m=N-j$ the number of free points, $t_a=a/r$. Free points $z_1,\dots,z_n$ in the support of $|m|$; $s_i=z_i/r$. Split $d|m|=w_\kappa d\Omega_\kappa$ as before ($w_\kappa=|E|^{1-\kappa}$ on $\Gamma$, $1$ at atoms), $\Omega_\kappa$ of total mass $\le2r+c_\kappa$.

\emph{Scaling.} Counting powers of $r$ in the two Cauchy determinants: free pairs $\frac{n(n-1)}2$ and node pairs $\frac{n(n-1)}2$ (since $|\mathcal R_S|=|\mathcal Q_S|=n$) against $n^2$ cross factors, so $|\mathsf C_{\mathcal R_S}\mathsf C_{\mathcal Q_S}|=r^{-2n}\,|\tilde{\mathsf C}_{\mathcal R_S}\tilde{\mathsf C}_{\mathcal Q_S}|$ with the same expression in the scaled variables; $r^{-2n}\le1$ is dropped. Hence $-\log\bigl(|\mathsf C_{\mathcal R_S}||\mathsf C_{\mathcal Q_S}|\prod w_\kappa(z_i)\bigr)\ge r^2H_S$ with
\begin{multline*}
 r^2H_S=\sum_{i\ne i'}\!-\log|s_i-s_{i'}|+\sum_{\mathcal F\in\{\mathcal R,\mathcal Q\}}\Bigl[\tfrac12\!\!\sum_{a\ne a'\in\mathcal F_S}\!\!-\log|t_a-t_{a'}|\\+\sum_{i}\sum_{a\in\mathcal F_S}\log|s_i-t_a|\Bigr]+\sum_i-\log w_\kappa(z_i).
\end{multline*}
(Ordered pairs; each family has weight $\frac12$ in the node--node and $1$ in the cross term, and the two together give weight $1$ to the free--free term.)

\emph{Completing the node set.} For each family $\mathcal F$ the two identities of the proof of Theorem~\ref{thm:N} (adding and removing the pairs meeting $S^c$) hold verbatim, because $S^c\subseteq C\subseteq\mathcal F$. Adding the two families with weights $\frac12$, the removed--removed pairs enter once with a plus sign (from the identity) and once with a minus sign (from the node energy), so they cancel \emph{exactly} against the removed--removed part of $[B]$ below --- this is why only common nodes can be removed (a node removed from one family only would leave a term of weight $\frac12$). Let $X$ be the set of the $n$ free points $s_i$ together with the $j$ removed nodes $t_a$ ($a\in S^c$); $|X|=N$, and the points are distinct ($\operatorname{Re}s_i<0\le t_a$). Then exactly
\[
 r^2H_S=[B]+[R1]+[R2]+[R3]+[R4],
\]
\begin{align*}
[B]&=\sum_{x\ne y\in X}-\log|x-y|, & [R1]&=\sum_{\mathcal F}\sum_i\sum_{a\in\mathcal F}\log|s_i-t_a|,\\
[R2]&=\tfrac12\sum_{\mathcal F}\sum_{a\ne a'\in\mathcal F}-\log|t_a-t_{a'}|, & [R3]&=\sum_{\mathcal F}\sum_{a\in S^c}\sum_{b\in\mathcal F,b\ne a}\log|t_a-t_b|,\\
[R4]&=\sum_i-\log w_\kappa(z_i).
\end{align*}

\emph{Error budget.} We use the following elementary bounds; $\ell_1:=\log(1+\tfrac{(1+L)8r}3)$ and $N=r+m+e\le\frac98r+3$.
\begin{itemize}[leftmargin=*]
\item $[B]$. Lemma~\ref{lem:K0} and Corollary~\ref{cor:Kp} for $X$ with the measure $N\sigma$ (weights $1$ on the $N$ points, $-Nm_k$ on the $c_k$; the proof of the Corollary is unchanged): $[B]\ge2N\sum_{x\in X}U^\sigma_\rho(x)-N^2I_\rho(\sigma)-N|\log\rho|$.
\item $[R1]$. For $\operatorname{Re}s\le-\frac38\cdot\frac1r$ the function $u\mapsto\log|s-u|$ is increasing on $[0,\infty)$. For a family $\mathcal F$ consisting of the nodes $t_a=a/r$ ($a\in C_0$, step $\frac1r$, weight $1$) and of the nodes of $X'$ (step $\frac2r$), a left Riemann sum of an increasing function $f$ loses at most $(f(\text{top})-f(\text{bottom}))$: $\sum_{a\in C_0}f(t_a)\ge r\int_0^1f-(f(1)-f(0))$ and $\sum_{i<m}f(1+\frac{2i}r)\ge\frac r2\int_1^{1+L}f-(f(1+L)-f(1))$. For $Y'$ (shifted by $\frac1r$) the sum is larger, and the extra node $b$ (at $t_b=1+L$, where $|s-t_b|\ge1$) contributes $\ge0$. Hence $\sum_{a\in\mathcal F}\log|s-t_a|\ge-rU^{\nu_L}(s)-\ell_1$ and $[R1]\ge-2r\sum_iU^{\nu_L}(s_i)-2n\ell_1$.
\item $[R2]$. Give each node $a$ of a family the cell $J_a=[t_a,t_a+d_a]$ of $\nu$-mass $\frac1r$ ($d_a=\frac1r$ in $C_0$, $\frac2r$ in $X'$/$Y'$; for $b$, $d_b=\frac2r$ in $\mathcal R$ and $\frac1r$ in $\mathcal Q$); for the families without $b$ the cells tile $\operatorname{supp}\nu_L$ exactly. For $a<a'$, $u\in J_a$, $u'\in J_{a'}$: $u'-u\ge t_{a'}-t_a-d_a$ and hence, if $a'$ is not the successor of $a$ in $\mathcal F$, $\log(t_{a'}-t_a)\le\log(u'-u)+\frac{d_a}{u'-u}\le\log(u'-u)+\frac1{\#\{\text{nodes between}\}+1}$; summing over pairs, the total error is $\le2N(1+\log N)$ per family; adjacent pairs cost $\le\frac32$ each, $\le\frac32N$ per family; the diagonal cells contribute $N(\log r+\frac32)$ to $r^2I(\nu_{\mathcal F})$ but nothing to the discrete sum. Altogether $[R2]\ge r^2\cdot\frac12\sum_{\mathcal F}I(\nu_{\mathcal F})-N(\log r+\tfrac32)-4N(1+\log N)-3N$, and $|I(\nu_{\mathcal F})-I(\nu_L)|\le\frac4r\cdot\sup|U|\le\frac{8}{r}$.
\item $[R3]$. For $a\in S^c$ ($t=t_a\in[0,1+L]$) and a family $\mathcal F$: for nodes $b<a$ the cell $J_b$ lies to the left of $t_a$ and $\log(t_a-t_b)\ge\frac r{d_b}\int_{J_b}\log(t_a-u)du$ exactly, because $\log(t_a-\cdot)$ is decreasing, so $\log(t_a-t_b)\ge\log(t_a-u)$ for $u\ge t_b$; for $b>a$ the error is $\frac{d_b}{t_b-t_a}$, summing to $\le2(1+\log N)$; the self cell contributes $\ge0$. Hence $\sum_{b\in\mathcal F,b\ne a}\log|t_a-t_b|\ge-rU^{\nu_L}(t_a)-2(1+\log N)-\frac12(\log r+1)-2$, and $[R3]\ge-2r\sum_{a\in S^c}U^{\nu_L}(t_a)-j\,(4(1+\log N)+\log r+5)$. (If $t_a=t_b$ is the extra node, $U^{\nu_L}$ is evaluated at $1+L$, where it is continuous.)
\item $[R4]$. $E$ is $1$-periodic, $x^*\equiv\frac58\pmod 1$, so $C_E=\sup_\Gamma|E|e^{2\pi|\operatorname{Im}z|}$ is independent of $r$, and $[R4]\ge(1-\kappa)\bigl(2\pi r\sum_i|\operatorname{Im}s_i|-n\log^+C_E\bigr)$.
\end{itemize}
Define $g^{(m)}:=g_{\alpha_m}$ for the actual $\alpha_m=m/r$ (note $\nu_L=\nu_{\alpha_m}$), put $\bar N=N/r$ and
\[
 G^{(m)}(x):=g^{(m)}(x)+\tfrac er\,U^\sigma_\rho(x).
\]
Dividing by $r^2$ and collecting (the coefficient of $\sum_{x\in X}U^\sigma_\rho(x)$ is $2\bar N=2(1+\alpha_m)+\frac{2e}r$, and the part $\frac{2e}r$ is put into $G^{(m)}$; the extra point $b$ is a free or removed point like the others):
\begin{equation}\label{eq:hHS}
 H_S\ge\frac2r\sum_{i\le n}G^{(m)}(s_i)+\frac2r\sum_{a\in S^c}G^{(m)}(t_a)-\bar N^2I_\rho(\sigma)+I(\nu_{\alpha_m})-\eta_1(r),
\end{equation}
\begin{multline*}
 \eta_1(r)=\frac{|\log\rho|\bar N}r+\frac{2\bar N\ell_1}r+\frac{\bar N(\log r+\frac{9}{2})+4\bar N(1+\log N)}r\\+\beta\frac{4(1+\log N)+\log r+5}r+\frac{\bar N\log^+C_E}r+\frac{16}r+\frac1r .
\end{multline*}
The term $\frac1r$ covers the finitely many pair terms involving $b$ that are not of the types treated above (in $\mathcal Q$ the cell of $b$ overlaps the top cell of $Y'$) and the self-energy of the mismatch $\nu_{\mathcal F}-\nu_L$ (at most three pieces of mass $\le\frac1r$ and length $\ge\frac1r$); together they are $O(\log r)$ in $r^2H_S$, hence $\le r$ for $r\ge2^{30}$. \emph{Numerical value.} For $r\ge2^{30}$: $\bar N\le1.1251$, $\log N\le\log r+0.118$, $\ell_1\le\log r+1.205$, $\beta\le1.0001$, $|\log\rho|=\log2048\le7.625$, and $\log^+C_E\le3.6$. For the last: $\Gamma$ consists of the rays $x^*+u(-1\pm3i)$, $u\ge0$, with $x^*\equiv\frac58\pmod1$, and $|\Im z|=3u$; for $u\le\frac18$ the real parts of $\pi z$ and of $\frac\pi4-\pi z$ lie in $[\frac\pi2,\frac{5\pi}8]$ and $[-\frac{3\pi}8,-\frac\pi4]$ modulo $\pi$, so $|E|\le\pi\sin\frac\pi4/(\sin\frac{3\pi}8\sin\frac\pi4)\le3.41$ and $|E|e^{6\pi u}\le3.41\,e^{3\pi/4}\le36$; for $u\ge\frac18$, $|\sin w|\ge\sinh|\Im w|$ gives $|E|e^{6\pi u}\le4\pi\sin\frac\pi4\,(1-e^{-3\pi/4})^{-2}\le10.9$. Altogether $\eta_1(r)\le(12.9\log r+52)/r$.

\begin{lemma}[certificates for the actual geometry]\label{lem:hB}
Let $r\ge2^{30}$, $|m-\frac r8|\le1$, $e\in\{0,1\}$. (i) For every free point $s$ on the rays, $G^{(m)}(s)\ge m'-\frac{12}r$. (ii) For the atoms, $\frac1r\sum_{\text{atoms}}\min\bigl(G^{(m)}(s)-m',0\bigr)\ge S_{\rm cert}-\frac{2}r\sum_\ell(m'-\gamma_\ell)^+-\frac3r$. (iii) For $a\in S^c\cap C_0$, $G^{(m)}(t_a)\ge g_{\min}-\frac{11}r$; and $G^{(m)}(t_b)\ge g_{\min}-\frac{2400}r$.
\end{lemma}
\begin{proof}
We use $\partial_\alpha g_\alpha(s)=U^\sigma_\rho(s)+\log|s-(1+2\alpha)|$, $|\alpha_m-\frac18|\le\frac1r$, $R_\sigma=\max_k|c_k|\le0.4767$ (from the data of $\sigma$), and
$-\log(|x|+R_\sigma+\rho)\le U^\sigma_\rho(x)\le\log\frac1\rho\le7.63$, since $|x-c_k|^2+\rho^2\le(|x|+R_\sigma+\rho)^2$.
(i) Compact part $u\le\frac32$: $|s|\le r_0+\omega+\frac32\sqrt{10}\le4.9$, so $U^\sigma_\rho(s)\ge-1.7$, and $1.39\le|s-(1+2\alpha)|\le6.2$; hence $|\partial_\alpha g|\le7.63+1.83\le10$, $|g^{(m)}(s)-g(s)|\le\frac{10}r$, and $\frac erU^\sigma_\rho(s)\ge-\frac{1.7}r$. Tail $u\ge\frac32$: by (R), $g(s)\ge T(u)\ge T(\frac32)+8(u-\frac32)\ge m'+11+8(u-\frac32)$, while $|\partial_\alpha g(s)|+|U^\sigma_\rho(s)|\le15.3+3\log(|s|+1.5)$ with $|s|\le0.15+3.17u$; for $r\ge2^{30}$ the loss $(15.3+3\log(3.2u+2))/r$ is far below $11+8(u-\frac32)$, so $G^{(m)}(s)\ge m'$.
(ii) The atoms lie in $\{-(\mathbb Z_{\ge0}+\frac34)/r\}\cup\{-\mathbb Z_{\ge1}/r\}\subset[-r_0,0]$, are distinct, at most $2r_0r+1$ in all, and at most $2r(b_\ell-a_\ell)+2$ of them lie in the closed cell $C_\ell$ of the segment cover (assign each atom to one cell). As in Lemma~\ref{lem:B}, $\frac1r\sum\min(g-m',0)\ge\sum_\ell\frac{2r(b_\ell-a_\ell)+2}r\min(\gamma_\ell-m',0)=S_{\rm cert}-\frac2r\sum_\ell(m'-\gamma_\ell)^+$ for $g=g_{1/8}$. On $[-r_0,0]$, $|s-(1+2\alpha)|\in[1.24,1.40]$, so $|\partial_\alpha g|\le10$, and $U^\sigma_\rho\ge-\log0.62>0$; so passing to $G^{(m)}$ costs at most $\frac1r(2r_0r+1)\frac{10}r\le\frac3r$.
(iii) For $t\in[0,1)$: $|t-(1+2\alpha)|\in[0.24,1.26]$ and $U^\sigma_\rho(t)\ge-0.4$, so $|\partial_\alpha g|\le7.63+1.43\le10$ and $G^{(m)}(t)\ge g(t)-\frac{10.4}r\ge g_{\min}-\frac{11}r$ by (L). For $t_b=1+2\alpha_m$ ($|t_b-\frac54|\le\frac2r$, the right end of $\operatorname{supp}\nu_{\alpha_m}$): put $\varphi(\alpha):=U^{\nu_\alpha}(1+2\alpha)=-\int_{2\alpha}^{1+2\alpha}\log v\,dv-\frac12\int_0^{2\alpha}\log v\,dv$; then $\varphi'(\alpha)=\log2\alpha-2\log(1+2\alpha)$, $|\varphi'|\le1.9$ for $|\alpha-\frac18|\le\frac1r$, and $U^\sigma_\rho$ is $\frac1{2\rho}=1024$-Lipschitz. So $|g^{(m)}(t_b)-g(\frac54)|\le1.126\cdot1024\cdot\frac2r+\frac{7.63}r+\frac{1.9}r\le\frac{2316}r$, $\frac erU^\sigma_\rho(t_b)\ge-\frac{0.6}r$, and $g(\frac54)\ge g_{\min}$ by (L).
\end{proof}
Combining \eqref{eq:hHS} with Lemma~\ref{lem:hB} ($n=N-j=r(1+\alpha_m-\beta)+e$ free points, some on the rays and the others at distinct atoms; $j$ removed nodes):
\[
 H_S\ge2(1+\alpha_m-\beta)m'+2S_{\rm cert}+2\beta g_{\min}-\bar N^2I_\rho+I(\nu_{\alpha_m})-\eta_2(r),
\]
\[
 \eta_2(r)=\eta_1(r)+\frac{24\bar N+6+22\beta}r+\frac4r\sum_\ell(m'-\gamma_\ell)^++\frac{4800}{r^2}
\]
(the term $\frac{24\bar N}r$: at most $n\le N$ ray points, each losing $\frac2r\cdot\frac{12}r$; $\frac6r$: twice the last term of (ii); $\frac{22\beta}r$: the removed nodes of $C_0$ in (iii); $\frac{4800}{r^2}$: the node $b$ if removed; we used $\frac erm'\ge0$). Finally we replace $\alpha_m$ by $\frac18$: $2(1+\alpha_m)m'\ge\frac94m'-\frac{4.44}r$; $|\bar N-\frac98|\le\frac2r$ gives $|\bar N^2-\frac{81}{64}|\le\frac{4.51}r$, and $I_\rho(\sigma)\le2.006$, a loss $\le\frac{9.05}r$; and $\frac d{d\alpha}I(\nu_\alpha)=2\varphi(\alpha)$ (as $\alpha$ grows, $\nu_\alpha$ gains mass at rate $1$ at the point $1+2\alpha$), with $0<\varphi\le0.7$ for $|\alpha-\frac18|\le\frac1r$ ($\varphi(\frac18)=0.6728$), so $|I(\nu_{\alpha_m})-I(\nu_{1/8})|\le\frac{1.4}r$. Hence
\[
 H_S\ge D_F-\beta c_{\rm hyb}-\eta(r),
\]
\begin{equation}\label{eq:heta}
 \eta(r)\le\frac{15\log r+K_0}{r},\qquad K_0=52+70+4\sum_\ell(m'-\gamma_\ell)^+\le25{,}100\le3\cdot10^4,
\end{equation}
where $52$ comes from $\eta_1$, $70\ge24\bar N+6+22\beta+4.44+9.05+1.4+4800/r$, and the dominant constant is $4\sum_\ell(m'-\gamma_\ell)^+=24{,}960.84$ for the $14{,}000$ segment cells (computed from the cell file \texttt{hcells\_h500.json} written by \texttt{hcert.py}). We do not try to optimise $\eta(r)$ or to give an explicit $r_0$ beyond $2^{30}$. All these constants affect only the $o(1)$ in Proposition~\ref{prop:hD}, not $D_F$ or $c_{\rm hyb}$.

\begin{theorem}[node removal, hybrid; computer-assisted]\label{thm:hN}
For admissible $\mathcal G(r,m,e)$ with $r\ge2^{30}$, every $S\subseteq C$ (with $j=|C\setminus S|$, $\beta=j/r$, $n=|S|+m$),
\[
 |\det M_S|\le16^{-n}(2r+c_\kappa)^n\exp\bigl(-r^2[D_F-\beta c_{\rm hyb}-\eta(r)]\bigr),
\]
with $\eta(r)$ as in \eqref{eq:heta}, independent of $S$.
\end{theorem}
\begin{proof}
Proposition~\ref{prop:hAndreief}, the splitting $d|m|=w_\kappa d\Omega_\kappa$, $\frac1{n!}\le1$, the pointwise bound just proved (it holds for all configurations of $n$ distinct free points on the support of $m$; coincident points give a vanishing Cauchy determinant), and $M_S=\frac1{16}L_F[\mathcal R_S,\mathcal Q_S]$.
\end{proof}
\begin{corollary}[window]\label{cor:hwindow}
For admissible $\mathcal G$, $r\ge2^{30}$, and real $t$ with $|t|\le e^{-c_{\rm hyb}r}$,
\[
 |\det\mathcal T_t|\le\exp\bigl(-D_Fr^2+\eta(r)r^2+(r+1)\log(2(2r+c_\kappa))\bigr)=e^{-(D_F-o(1))r^2}.
\]
This is Proposition~\ref{prop:hD}.
\end{corollary}
\begin{proof}
By \eqref{eq:hexpand} and Theorem~\ref{thm:hN}, with $e^{r^2\beta c}=e^{jcr}$ and $(2r+c_\kappa)^n\le(2r+c_\kappa)^{r+1+m}$,
\[
 |\det\mathcal T_t|\le(2r+c_\kappa)^{N}e^{r^2\eta}e^{-D_Fr^2}\sum_{j=0}^{|C|}\binom{|C|}j16^{-(N-j)}\bigl(|t|e^{c_{\rm hyb}r}\bigr)^j,
\]
and for $|t|e^{c_{\rm hyb}r}\le1$ the sum is $\le2^{|C|}$. (Here $N=r+m+e$, and $(2r+c_\kappa)^N=e^{O(r\log r)}$.)
\end{proof}

\subsection{Reproducing the certificate}\label{ssec:hrepro}
The directory \texttt{scripts-hybrid/} contains \texttt{hcert.py}, \texttt{sig\_h500.json}, \texttt{sigma300.json} (the plain control), the expected logs and a \texttt{README}. The run is
\begin{verbatim}
python hcert.py sig_h500.json 1/8 2.2185 -1.6090 h500 14000 512     # D_F, c_hyb
python hcert.py sigma300.json 0 1.97180 -1.46410 ctl0 14000          # negative control
\end{verbatim}
(arguments: measure file, $\alpha$, target for $m'$, target for $g_{\min}$, tag, number of segment cells, number of line cells beyond $1$; targets must lie below the float plateau minimum). Requirements: Python with \texttt{python-flint} (Arb) and \texttt{numpy}; set the environment variable \texttt{PYLIBS} if these live in a non-standard directory. Run time about $30$ s. The output was reproduced independently by an audit agent (identical numbers; $8$ s), and its closed forms were compared with quadrature.

\section{$\ell$-adic nonvanishing for the hybrid frame}\label{sec:hnv}

This section proves Proposition~\ref{prop:hNV}, modulo the nonvanishing of two Krylov determinants (Theorem~\ref{thm:hK}) which is proved in \S\ref{sec:hkrylov}. The structure is that of \S\ref{sec:nonvanishing}--\S\ref{sec:krylov}: reduce $\det\Fu$ modulo $\ell$ to a \emph{bordered Cauchy--Gram} determinant (Proposition~\ref{prop:hred}); the boundary member is controlled at $Y=0$ by a product of two Cauchy determinants; the other roots are controlled by a rank-one bordering argument that ends in a Krylov determinant.

Throughout $\ell\equiv3\pmod4$ is prime, $\ell\ge7$, $K=(\ell+1)/4\ge2$, $0\le m\le K-1$, $s=K-m\ge1$, $r_1=\ell-K-m$. Indices $0,\dots,K-1$ are \emph{low}, indices $K\le b\le\ell-1$ \emph{high}; for a high node $b$ let $j_b=\ell-b\in[1,\ell-K]$ and, for $j\ge1$, $w_j=\bigl(\frac1{a+j}\bigr)_{0\le a<K}\in\F_\ell^K$ (all denominators are units because $a+j\le K-1+\ell-K<\ell$). For a set $J$ of positive integers $W_J=[w_j]_{j\in J}$, a $K\times|J|$ matrix, written $C_J$ when $J$ has $K$ elements. A geometry is \emph{admissible for $\ell$} if

\textbf{(H)} $\{0,\dots,K-1\}\subseteq C$ and all nodes are $\le\ell-1$.

Both $\mathcal G(r_1,m,0)$ and $\mathcal G(r_1,m,1)$ satisfy (H): $r_1\ge K\iff m\le2K-1$, and the top node of member~2 is $r_1+2m=\ell-K+m\le\ell-1\iff m\le K-1$. In $j$-coordinates, member~1 has $\{j_b:b\notin C\cup X'\cup Y'\}=G_1:=\{1,\dots,s\}$ ($G$ = ``absent'' nodes), $J_{Y'}=s+\{1,3,\dots,2m-1\}$, $J_{X'}=s+\{2,4,\dots,2m\}$ (the node $b=r_1+2i+1$ has $j_b=s+(2m-2i-1)$ and $r_1+2i$ has $s+2(m-i)$); member~2 has $G_2=\{1,\dots,s-1\}$ (its extra node $b=r_1+2m$ is $j=s$) and the same $J_{X'},J_{Y'}$.

\subsection{The $\ell$-singular structure}
\begin{lemma}\label{lem:hstruct}
Assume (H). Then the only $\ell$-singular quantity among $k,\,4k-1,\,4,\,b-a$ occurring in the entries $\Rk(a,b)$ ($a,b\le\ell-1$) is $4K-1=\ell$ at $k=K$ (note $3\ell\equiv1\pmod4$ and $2\ell$ is even, so $4k-1\in\{\ell,2\ell,3\ell\}$ forces $k=K$). Consequently:
\begin{enumerate}[label=(\alph*),leftmargin=*]
\item $\Rk(a,b)\in\Z_{(\ell)}$ for $a,b$ low, and for $a\ne b$ both high;
\item for $a$ low and $b$ high: $\ell\,\Rk(a,b)\equiv-\tfrac14\,w_{j_b}(a)\pmod\ell$;
\item for $b$ high: $\ell^2\Rk(b,b)\equiv-1\pmod\ell$.
\end{enumerate}
\end{lemma}
\begin{proof}
The term $k=K$ occurs in $\sum_{a<k\le b}$ iff $a<K\le b$, and in $\Psi(b)$ iff $b\ge K$. For (b): $\ell\Rk(a,b)\equiv\frac1{16K(b-a)}\equiv\frac1{4(b-a)}$ since $16K=4\ell+4$, and $b-a=\ell-j_b-a\equiv-(a+j_b)$. For (c): $\Psi(b)\ni-\frac1{(4K-1)^2}=-\ell^{-2}$.
\end{proof}
\begin{lemma}[low block]\label{lem:hlow}
Let $A=-\Rk[\text{low},\text{low}]$ (the low block of $M_0$) and $H^{(2)}_n=\sum_{i\le n}i^{-2}$. Then in $\F_\ell$: $16A=W_{[1,\ell-K]}W_{[1,\ell-K]}^{\top}+H^{(2)}_{K-1}I_K$.
\end{lemma}
\begin{proof}
This is the computation of the proof of Lemma~\ref{lem:2} with $m=\ell-K$ (it only used $a<b<K$): $16A_{ab}\equiv\sum_{j=1}^{\ell-K}\frac1{(a+j)(b+j)}$ for $a<b$, and $16A_{aa}\equiv\sum_{j=1}^{\ell-K}\frac1{(a+j)^2}-H^{(2)}_{\ell-K}$; and $H^{(2)}_{\ell-K}\equiv-H^{(2)}_{K-1}$ because $H^{(2)}_{\ell-1}\equiv0$ and $1/(\ell-i)^2\equiv1/i^2$.
\end{proof}

\begin{proposition}[bordered reduction]\label{prop:hred}
Assume (H). Let $G$ be the set of $j_b$, $b\in[K,\ell-1]\setminus(C\cup X'\cup Y')$, $W_G$ the corresponding matrix ($K\times|G|$, $GG^{\top}:=W_GW_G^\top$ below), $\#C_{\rm high}=|C\setminus[0,K)|$ and, for an indeterminate $Y$,
\[
 B(Y)=\det\begin{pmatrix}Y\,I_K+W_GW_G^{\top}&W_{Y'}\\ W_{X'}^{\top}&0\end{pmatrix}\in\F_\ell[Y]\quad(\text{size }K+m).
\]
Then $D\,\Fu(X)\,D\in M(\Z_{(\ell)}[X])$ for $D=\operatorname{diag}$ ($1$ on low, $\ell$ on high rows $C_{\rm high}\cup X'$ and high columns $C_{\rm high}\cup Y'$), and
\[
 \ell^{2(\#C_{\rm high}+m)}\det\Fu(X)\equiv(-1)^{\#C_{\rm high}+m}(-16)^{-K}\,B\bigl(H^{(2)}_{K-1}-16X\bigr)\pmod\ell\quad\text{in }\F_\ell[X].
\]
\end{proposition}
\begin{proof}
By Lemma~\ref{lem:hstruct}, modulo $\ell$ the matrix $D\Fu(X)D$ (rows $L,C_{\rm high},X'$; columns $L,C_{\rm high},Y'$; $L$=low) is
\[
 \begin{pmatrix}XI-A&-\frac14W_{C}&-\frac14W_{Y'}\\-\frac14W_C^\top&-I&0\\-\frac14W_{X'}^\top&0&0\end{pmatrix}
\]
(the zero blocks: $\ell^2\Rk$ between distinct high nodes and $\ell^2X$ are $\equiv0$; a diagonal entry occurs only at $C\times C$). Taking the Schur complement of the block $-I$ (factor $(-1)^{\#C_{\rm high}}$; the blocks $C_{\rm high}\times Y'$ and $X'\times C_{\rm high}$ vanish so no other block changes) the upper left block becomes $XI-A+\frac1{16}W_CW_C^\top$, which by Lemma~\ref{lem:hlow} and the partition $[1,\ell-K]=J_C\sqcup G\sqcup J_{X'}\sqcup J_{Y'}$ equals $-\frac1{16}\bigl(YI+W_GW_G^\top+W_{X'}W_{X'}^\top+W_{Y'}W_{Y'}^\top\bigr)$, $Y=H^{(2)}_{K-1}-16X$. Adding $-\frac14W_{X'}\times(\text{bottom rows})$ to the top rows and $(\text{right columns})\times(-\frac14W_{Y'}^\top)$ to the left columns removes $W_{X'}W_{X'}^\top$ and $W_{Y'}W_{Y'}^\top$ (the bottom right block is $0$, so the off-diagonal blocks do not change). Finally multiply the top $K$ rows by $-16$, the $m$ bottom rows by $-4$ and the $m$ right columns by $\frac14$; the determinant changes by $(-16)^K(-1)^m$ and we obtain $B(Y)$. Also $\det D=\ell^{\#C_{\rm high}+m}$ for rows and the same for columns.
\end{proof}
Hence $\deg B\le K-m$, and, since $X\mapsto Y$ is a bijection of $\F_\ell$ \emph{independent of the geometry}, for a rational $x_0$ with $\ell\nmid\den x_0$ and $Y_0=H^{(2)}_{K-1}-16x_0\bmod\ell$: if $B(Y_0)\ne0$ then $v_\ell(\det\Fu(x_0))=-2(\#C_{\rm high}+m)$ exactly; in particular $\det\Fu(x_0)\ne0$. (Reduction commutes with evaluation because $x_0\in\Z_{(\ell)}$.)

\begin{corollary}[boundary value]\label{cor:hbdry}
For member~1 ($|G_1|+m=K$): $B_1(0)=\pm\det C_{G_1\cup Y'}\cdot\det C_{G_1\cup X'}\ne0$ in $\F_\ell$.
\end{corollary}
\begin{proof}
At $Y=0$ the matrix is $\begin{psmallmatrix}C_GC_G^\top&W_{Y'}\\W_{X'}^\top&0\end{psmallmatrix}=PQ$ with $P=\begin{psmallmatrix}C_G&W_{Y'}&0\\0&0&I_m\end{psmallmatrix}$, $Q=\begin{psmallmatrix}C_G^\top&0\\0&I_m\\W_{X'}^\top&0\end{psmallmatrix}$, both square of size $K+m=s+2m$; $\det P=\pm\det[C_G,W_{Y'}]$, $\det Q=\pm\det[C_G,W_{X'}]$. Each is a $K\times K$ Cauchy determinant $\prod_{a<a'}(a'-a)\prod_{j<j'}(j'-j)/\prod(a+j)$ whose factors are integers in $[1,2K+m-1]\subseteq[1,3K-2]$, and $3K-2<4K-1=\ell$. So the determinant is an $\ell$-unit.
\end{proof}

\subsection{The dual formulation and bordering}
For a member put $\Gamma_{ij}=\sum_{a<K}\frac1{(a+i)(a+j)}$, i.e.\ $\Gamma=W^\top W$, and let
\[
 \hat\Gamma=C_{G\cup X'}^\top C_{G\cup Y'}\ (=\Gamma[G\sqcup J_{X'},G\sqcup J_{Y'}]),\qquad E=\operatorname{diag}(I_G,0_m),\qquad D(z)=zE+\hat\Gamma .
\]
\begin{lemma}\label{lem:hdual}
For $z\ne0$, $\ker M(z)\cong\ker D(z)$ and the same for left kernels, where $M(z)=\begin{psmallmatrix}zI+C_GC_G^\top&W_{Y'}\\W_{X'}^\top&0\end{psmallmatrix}$ is the matrix of $B(z)$. Hence the nonzero roots of $B$ and of $\det D$ coincide.
\end{lemma}
\begin{proof}
Let $(x,c)\in\ker M(z)$: $(z+C_GC_G^\top)x+W_{Y'}c=0$ and $W_{X'}^\top x=0$. Put $\tilde u=-(C_G^\top x,\,c)/z$. Then $x=C_{G\cup Y'}\tilde u$ and the rows of $D(z)\tilde u$ are $C_{G\cup X'}^\top$ applied to $x$ (rows in $X'$ give $W_{X'}^\top x=0$; rows in $G$ give $C_G^\top x+z\tilde u_G=0$). Conversely, for $\tilde u\in\ker D(z)$ put $x=C_{G\cup Y'}\tilde u$, $c=-z\tilde u_{Y'}$; then $(z+C_GC_G^\top)x+W_{Y'}c=z(x-C_{G\cup Y'}\tilde u)=0$ using $C_G^\top x=-z\tilde u_G$, and $W_{X'}^\top x=0$. Both maps are injective as $z\ne0$. Left kernels: transpose (exchange $X'\leftrightarrow Y'$).
\end{proof}
Let member~1 have absent set $G_1=\{1..s\}$ and member~2 $G_2=G_1\setminus\{g\}$, $g=s$, with the same $X',Y'$. Then $\hat\Gamma_1$ is a $K\times K$ matrix $=C_{G_1\cup X'}^\top C_{G_1\cup Y'}$, invertible over $\F_\ell$ by the proof of Corollary~\ref{cor:hbdry}. Put
\[
 T=(\hat\Gamma_1^{-1})_{G_1G_1}\in M_s(\F_\ell),\qquad k_R=\det\bigl[e_g^\top T^k\bigr]_{0\le k<s},\qquad k_L=\det\bigl[T^ke_g\bigr]_{0\le k<s}.
\]
\begin{lemma}[no common root]\label{lem:hCC}
If $B_1$ and $B_2$ have a common root $z_0\ne0$ in $\overline{\F}_\ell$, then $k_R=0$ or $k_L=0$.
\end{lemma}
\begin{proof}
By Lemma~\ref{lem:hdual}, $\det D_1(z_0)=\det D_2(z_0)=0$. Now $D_2$ is $D_1$ with row $g$ and column $g$ deleted. Write $D_1=\begin{psmallmatrix}D_2&\gamma\\\varrho^\top&z+\Gamma_{gg}\end{psmallmatrix}$; then $\det D_1=(z+\Gamma_{gg})\det D_2-\varrho^\top\operatorname{adj}(D_2)\gamma$, so $\varrho^\top\operatorname{adj}(D_2)\gamma=0$ at $z_0$. If $\operatorname{rank}D_2\le n_2-2$ ($n_2$ its size) then $\ker D_2$ has dimension $\ge2$ and contains $x\ne0$ with $\varrho^\top x=0$. If $\operatorname{rank}D_2=n_2-1$, $\operatorname{adj}(D_2)=c\,xy^\top$ with $c\ne0$, $D_2x=0$, $y^\top D_2=0$, and $(\varrho^\top x)(y^\top\gamma)=0$.
\emph{Right case} ($\varrho^\top x=0$, $x\ne0$): $\tilde u=[x;0]\in\ker D_1(z_0)$. From $z_0E\tilde u+\hat\Gamma_1\tilde u=0$: $\tilde u=-z_0\hat\Gamma_1^{-1}E\tilde u$; restricting to $G_1$: $T\tilde u_G=-z_0^{-1}\tilde u_G$, and $\tilde u_G\ne0$ (otherwise $\tilde u=0$), with $\tilde u_g=0$. Thus $e_g^\top T^k\tilde u_G=(-z_0)^{-k}\tilde u_g=0$ for all $k$, i.e.\ the matrix $[e_g^\top T^k]_{k<s}$ is singular: $k_R=0$.
\emph{Left case} ($y^\top\gamma=0$): $[y^\top,0]$ is a left kernel vector of $D_1(z_0)$; likewise $y_G^\top T=-z_0^{-1}y_G^\top$ with $y_g=0$, so $[T^ke_g]_{k<s}$ is singular and $k_L=0$.
\end{proof}
(If $s=1$, i.e.\ $m=K-1$, then $G_2=\emptyset$, $T$ is $1\times1$ and $k_R=k_L=1$.)

\begin{theorem}[$\ell$-adic nonvanishing, assuming Theorem~\ref{thm:hK}]\label{thm:hNV}
Let $x_0\in\Q$, $\ell\nmid\den x_0$, $Y_0=H^{(2)}_{K-1}-16x_0\bmod\ell$. If $k_R\ne0\ne k_L$ in $\F_\ell$, then $B_1(Y_0)\neq0$ or $B_2(Y_0)\ne0$; hence one of $\mathcal G(r_1,m,0)$, $\mathcal G(r_1,m,1)$ has $v_\ell(\det\Fu(x_0))=-2(\#C_{\rm high}+m)$ and in particular $\det\Fu(x_0)\ne0$.
\end{theorem}
\begin{proof}
If $Y_0=0$: $B_1(0)\ne0$ (Corollary~\ref{cor:hbdry}). If $Y_0\ne0$ and $B_1(Y_0)=B_2(Y_0)=0$, Lemma~\ref{lem:hCC} contradicts $k_Rk_L\ne0$. Then apply the consequence of Proposition~\ref{prop:hred} (note $\#C_{\rm high}$ differs by one between the two members).
\end{proof}

\paragraph{The selection of $m$, and $\kappa$.} Theorem~\ref{thm:hNV} holds for every $m\in[0,K-1]$ with $r_1=\ell-K-m$, so we are free to take $m=\operatorname{round}((\ell-K)/9)$ (so that $m/r_1\to\frac18$). Then $r_1/\ell\to\frac{8}{9}\cdot\frac34=\frac23=\kappa=\frac{3/4}{1+\alpha}$ and the two members have scales equal to $r_1$. The condition $r+m\le\ell-K$ is exactly the existence of a forced common root: $\deg B_i$ drops (the top $K$ rows of $M(0)$ have rank $\le s+m$), so $Y^{\max(0,K-s-m)}\mid B$, and members with $r+m>\ell-K$ would have $B(0)=0$; this is why $\kappa$ cannot be raised. (Members with $r+m<\ell-K$ have no forced root but $B(0)$ is then a Gram-type minor which is not a product, so $Y_0=0$ is not covered by the argument.) The condition $\ell\nmid\den x_0$ is removed by Proposition~\ref{prop:hL3}.

The remaining input is the nonvanishing of $k_R,k_L$ modulo $\ell$ for \emph{every} prime $\ell\equiv3\pmod4$, $\ell\ge7$ and every $0\le m\le K-1$. For $m=0$ this is the nonvanishing of the Krylov determinant of Theorem~A; the general case is a closed form of the same kind, proved next.

\section{The bordered Krylov determinants}\label{sec:hkrylov}

In this section $K\ge2$, $0\le m\le K-2$ (the case $m=K-1$ is trivial: $k_R=k_L=1$), $s=K-m\ge2$, and
\begin{gather*}
 G=\{1,\dots,s\},\quad Y=\{s+1,s+3,\dots,s+2m-1\},\quad X=\{s+2,s+4,\dots,s+2m\},\\ Z=X\sqcup Y=\{s+1,\dots,s+2m\},\quad g=s .
\end{gather*}
For a real parameter $\theta$ let $y_a=a+\theta$ ($0\le a<K$), $c_j=(1/(y_a+j))_a\in\R^K$, $C_J=[c_j]_{j\in J}$ and
\begin{gather*}
 \hat\Gamma(\theta)=C_{G\cup X}^\top C_{G\cup Y},\qquad T(\theta)=(\hat\Gamma(\theta)^{-1})_{GG},\\ k_R(\theta)=\det[e_g^\top T^k]_{k<s},\quad k_L(\theta)=\det[T^ke_g]_{k<s}.
\end{gather*}
(For $\theta\ge0$ all entries $y_a+j$ are positive and the nodes are distinct, so the two $K\times K$ Cauchy matrices are invertible.) For $\theta=0$ these are the matrices of \S\ref{sec:hnv} computed over $\Q$; their reductions modulo $\ell=4K-1$ are the objects there, see Lemma~\ref{lem:hbridge}. Write $T_2(n)=n(n-1)$ for $n\ge1$, $T_2(n)=0$ for $n\le0$, $D=\frac{s(s-1)(4s+6m+1)}3$ and define, for $1\le c\le2K-1$,
\begin{equation}\label{eq:hec}
 e_c=\begin{cases}c\,\mathbf 1[c<s]+T_2(\min(c,s))+T_2(c-m)&1\le c\le K,\\ T_2(n)+T_2(o)+d'\,\mathbf 1[d'<s]&c=K+d',\ 1\le d'\le K-1,\end{cases}
\end{equation}
with $o=(s-d')_+$, $|F|=(d'-m)_+$, $n=s-|F|$.
\begin{theorem}\label{thm:hK}
$k_R(\theta)$ and $k_L(\theta)$ are polynomials in $\theta$ with rational coefficients, and
\[
 k(\theta)=\varkappa\prod_{c=1}^{2K-1}(\theta+c)^{e_c},\qquad k\in\{k_R,k_L\},
\]
where $\varkappa\in\Q^\times$ has only prime factors $\le2K-1$. Consequently $k_R(0),k_L(0)$ are nonzero rationals with numerator and denominator prime to $\ell=4K-1$, and for $\ell$ prime $k_R,k_L\not\equiv0\pmod\ell$ in the sense of \S\ref{sec:hnv}.
\end{theorem}
The proof occupies this section: reduction to a symmetric matrix of Hahn type (\S\ref{ssec:hH}), the order of vanishing at each $\theta=-c$ (\S\ref{ssec:hz1}, \S\ref{ssec:hz2}), the degree and leading coefficient (\S\ref{ssec:hlead}), and $\sum_ce_c=D$ (\S\ref{ssec:hsum}).

\subsection{Polynomiality and the symmetric form}\label{ssec:hH}
\begin{lemma}\label{lem:hpoly}
$k_R,k_L\in\Q[\theta]$.
\end{lemma}
\begin{proof}
The inverse of a Cauchy matrix $[1/(y_a+j_b)]$ has entries
\[\pm\frac{\prod_k(y_a+j_k)\prod_{a'}(y_{a'}+j_b)}{(y_a+j_b)\prod_{k\ne b}(j_k-j_b)\prod_{a'\ne a}(y_{a'}-y_a)};\]
the factor $y_a+j_b$ occurs in both products, so the entry is a polynomial in $\theta$ (the differences $y_{a'}-y_a=a'-a$ are constants). Hence $\hat\Gamma^{-1}=C_{G\cup Y}^{-1}C_{G\cup X}^{-\top}$, $T$, $k_R$, $k_L$ have polynomial entries.
\end{proof}
Put $\omega(x)=\prod_{i=1}^s(x-i)$, $\pi_a=\prod_{b\ne a}(b-a)=\pm a!(K-1-a)!$, $\Pi_X(z)=\prod_{x\in X}(z-x)$, $\Pi_Y$ likewise, $\Pi_Z=\Pi_X\Pi_Y$, and
\[
 \Omega_a=\frac{(y_a+1)_s(y_a+1)_{s+2m}}{\pi_a^2},\quad A=C_G^\top\operatorname{diag}(\Omega_a)C_G\in M_s(\Q[\theta]),\quad \Psi_j=\frac{(j+\theta)_K}{|\omega'(j)|\sqrt{\Pi_Z(j)}} .
\]
($\Pi_Z(j)>0$ for $j\le s$, as it is a product of $2m$ negative numbers.) Let $\Lambda=\operatorname{diag}\bigl((\Pi_X(j)/\Pi_Y(j))^{1/2}\bigr)_{j\in G}$, a \emph{constant} positive matrix (both products have sign $(-1)^m$), and $\mathsf H=\Psi A\Psi$, real symmetric, with entries polynomial in $\theta$ up to the constant square roots.
\begin{proposition}\label{prop:hH}
$T=\Lambda\,\mathsf H\,\Lambda^{-1}$. Consequently, with $\kappa_{\mathsf H}=\det[e_g^\top\mathsf H^k]_{k<s}=\det[\mathsf H^ke_g]_{k<s}$,
\[
 k_R=\Lambda_g^{\,s}(\det\Lambda)^{-1}\kappa_{\mathsf H},\qquad k_L=\Lambda_g^{-s}(\det\Lambda)\kappa_{\mathsf H},\qquad k_Rk_L=\kappa_{\mathsf H}^2,
\]
so $k_R/k_L$ is a nonzero constant independent of $\theta$.
\end{proposition}
\begin{proof}
\emph{Step 1.} Block inversion of $\hat\Gamma=\begin{psmallmatrix}C_G^\top C_G&C_G^\top C_Y\\C_X^\top C_G&C_X^\top C_Y\end{psmallmatrix}$ gives $T=S^{-1}$, $S=C_G^\top\Pi C_G$, $\Pi=I-C_Y(C_X^\top C_Y)^{-1}C_X^\top$, the projector with range $\ker C_X^\top$ and kernel $\operatorname{range}C_Y$. If the columns of $U$ span $\ker C_X^\top$ and those of $V$ span $\ker C_Y^\top$, then $\Pi=U(V^\top U)^{-1}V^\top$.
\emph{Step 2.} For $u\in\R^K$, $\sum_a\frac{u_a}{z+y_a}=\frac{p(z)}{\prod_b(z+y_b)}$ with $\deg p<K$ and $u_a=p(-y_a)/\pi_a$. Hence $u\in\ker C_X^\top$ iff $p=\Pi_Xq$, $\deg q<s$; similarly $\ker C_Y^\top$ with $\Pi_Y$. Choose $q,r$ in the Lagrange basis $\ell_i$ of the nodes $1,\dots,s$: then $C_G^\top U=D_X:=\operatorname{diag}\bigl(\Pi_X(j)/(j+\theta)_K\bigr)$, $V^\top C_G=D_Y$ (same with $\Pi_Y$), and $M:=V^\top U$ has entries $\sum_a\Pi_X(-y_a)\Pi_Y(-y_a)\ell_i(-y_a)\ell_{i'}(-y_a)/\pi_a^2$. Therefore $S=D_XM^{-1}D_Y$, $T=D_Y^{-1}MD_X^{-1}$.
\emph{Step 3.} $\ell_i(-y)=-\omega(-y)/((y+i)\omega'(i))$, $\Pi_Z(-y)=(y+s+1)_{2m}$, $\omega(-y)^2=(y+1)_s^2$; so $M=D_{\omega'}^{-1}AD_{\omega'}^{-1}$ and $T=\Lambda_1A\Lambda_2$ with $\Lambda_1=\operatorname{diag}\frac{(j+\theta)_K}{\Pi_Y(j)\omega'(j)}$, $\Lambda_2=\operatorname{diag}\frac{(j+\theta)_K}{\Pi_X(j)\omega'(j)}$. Since $\Lambda_1\Lambda_2=\Psi^2$ and $\Lambda_1/\Lambda_2=\Lambda^2$, $\Lambda_1=\Psi\Lambda$, $\Lambda_2=\Psi\Lambda^{-1}$, i.e.\ $T=\Lambda\Psi A\Psi\Lambda^{-1}$. The formulas for $k_R,k_L$: $e_g^\top T^k=\Lambda_ge_g^\top\mathsf H^k\Lambda^{-1}$ and $T^ke_g=\Lambda_g^{-1}\Lambda\mathsf H^ke_g$.
\end{proof}
\begin{lemma}[moment form]\label{lem:hmoment}
Let $\varrho_a=(y_a+s+1)_{2m}/\pi_a^2$. There is a constant invertible matrix $\mathcal E\in M_s(\Q)$ with $A=\mathcal E\,\mathbb M\,\mathcal E^\top$, $\mathbb M=\bigl[\sum_{a<K}\varrho_a\,y_a^{p+q}\bigr]_{p,q<s}$. Equivalently, for $f\in\R^G$ and $P_f(x)=\sum_jf_j\,\omega(x)/(x-j)$:
\begin{gather*}
 f^\top Af=\sum_{a<K}\varrho_a\,P_f(x_a)^2,\qquad f^\top\Psi^{-2}f=\sum_{j=1}^s\nu_jP_f(j)^2,\\ x_a=-y_a,\qquad \nu_j=\frac{\Pi_Z(j)}{\prod_b(j-x_b)^2},\qquad \varrho_a=\frac{\Pi_Z(x_a)}{\prod_{b\ne a}(x_a-x_b)^2}.
\end{gather*}
\end{lemma}
\begin{proof}
The polynomial $P_y(x)=\frac{1-\omega(x)/\omega(-y)}{x+y}=-\frac{Q(x,y)}{\omega(-y)}$, $Q=\frac{\omega(x)-\omega(-y)}{x+y}=\sum_{r<s}x^rq_r(y)$, interpolates $x\mapsto\frac1{y+x}$ at $x=1,\dots,s$; $q_r$ has degree exactly $s-1-r$ with leading coefficient $\pm1$, so $(q_r(y))_r=\mathcal N(y^p)_p$ with $\mathcal N$ constant and anti-triangular. Hence $C_G=-D_{1/\omega(-y)}\hat Y\mathcal N^\top V_s^\top$ ($\hat Y=[y_a^p]$, $V_s=[j^r]$), and $A=\mathcal E\hat Y^\top D_\varrho\hat Y\mathcal E^\top$, $\mathcal E=V_s\mathcal N$, since $\Omega_a/\omega(-y_a)^2=\varrho_a$. The quadratic-form identities follow from $\sum_jf_j/(y+j)=-P_f(-y)/\omega(-y)$ and $P_f(j)=f_j\omega'(j)$, $\Psi_j^{-2}=\omega'(j)^2\Pi_Z(j)/(j+\theta)_K^2$.
\end{proof}
So the eigenvalues of $\mathsf H$ are those of the pencil $\langle\cdot,\cdot\rangle_\varrho-\lambda\langle\cdot,\cdot\rangle_\nu$ on $\mathrm{Poly}_{<s}$, and $\mathsf H$ has entries polynomial in $\theta$ (up to constant roots).

\subsection{Analytic tools}
Fix $c\in\{1,\dots,2K-1\}$, $\theta=-c+\epsilon$ ($\epsilon$ real), and write $\mathsf H(\epsilon)$.
\begin{lemma}[Rellich]\label{lem:hRell}
Let $H(\epsilon)$ be a real symmetric $s\times s$ matrix, real-analytic in $\epsilon$ near $0$. There are analytic eigenvalue branches $\lambda_i(\epsilon)$ and an analytic orthogonal $U(\epsilon)$ of eigenvectors, and $\kappa(\epsilon):=\det[e_g,He_g,\dots,H^{s-1}e_g]=\pm\prod_iu_{ig}\prod_{i<i'}(\lambda_{i'}-\lambda_i)$ ($u_{ig}=U_{gi}$). Hence $\operatorname{ord}_0\kappa=\sum_i\operatorname{ord}u_{ig}+\sum_{i<i'}\operatorname{ord}(\lambda_i-\lambda_{i'})$. If for all small $\epsilon\ne0$ at least $d$ eigenvalues of $H(\epsilon)$ lie in a window of width $C\epsilon^p$, then some $d$ branches satisfy $\operatorname{ord}(\lambda_i-\lambda_{i'})\ge p$ pairwise.
\end{lemma}
\begin{proof}
Rellich's theorem; $H^ke_g=\sum_i\lambda_i^ku_{ig}\varphi_i$ so the Krylov matrix is $U\operatorname{diag}(u_{\cdot g})[\lambda_i^k]$ (Vandermonde), $\det U=\pm1$ constant. For the last claim, along a sequence $\epsilon_n\to0$ a fixed $d$-subset of the finitely many branches lies in the window; a nonzero analytic function of order $q<p$ is $\gg|\epsilon|^q$ near $0$. (Identically coinciding branches give $\kappa\equiv0$ and $\operatorname{ord}=\infty$; the inequality remains true.)
\end{proof}
\begin{lemma}[cluster lemma]\label{lem:hclu}
Let $\mathcal M(\epsilon)=\mathcal M_0+\epsilon\mathcal M_1+O(\epsilon^2)$ be real symmetric and analytic, suppose $\mathcal M_0$ has an eigenvalue $1$ with eigenspace $\mathcal E_1$ (dimension $d$) isolated from the rest of its spectrum, and $P_{\mathcal E_1}\mathcal M_1P_{\mathcal E_1}=0$. Then $\mathcal M(\epsilon)$ has $d$ eigenvalues $1+O(\epsilon^2)$.
\end{lemma}
\begin{proof}
Kato's first-order perturbation theory: the splittings of order $\epsilon$ are the eigenvalues of $P_{\mathcal E_1}\mathcal M_1P_{\mathcal E_1}$ (or: conjugate by $\exp(\epsilon S)$ with $S$ antisymmetric cancelling the $\mathcal E_1$--$\mathcal E_1^\perp$ blocks of $\mathcal M_1$, then Weyl's inequality).
\end{proof}
\begin{lemma}[scaled Schur lemma]\label{lem:hsch}
Let $H(\epsilon)=\Gamma\hat H(\epsilon)\Gamma$ with $\hat H$ real symmetric analytic, $\Gamma=\operatorname{diag}(\epsilon\text{ on }J_1,\,1\text{ on }J_2)$, and $\hat H_{22}(0)$ invertible. Then exactly $|J_1|$ eigenvalues of $H$ are $O(\epsilon^2)$, the other $|J_2|$ tend to the eigenvalues of $\hat H_{22}(0)$ (nonzero), and the unit eigenvectors of the latter have $J_1$-components $O(\epsilon)$; the eigenvectors of the former have $J_2$-components $O(\epsilon)$.
\end{lemma}
\begin{proof}
For $|\lambda|\le\delta$, $\det(H-\lambda)=\det(\hat H_{22}-\lambda+O(\epsilon^2))\det(\epsilon^2X(\lambda,\epsilon)-\lambda I)$ by the Schur complement with $X$ bounded; the second factor has exactly $|J_1|$ roots in $|\lambda|\le\delta$, all $O(\epsilon^2)$ (Rouch\'e). From $(\hat H_{22}-\lambda)u_2=-\epsilon\hat H_{21}u_1+O(\epsilon^2)u_2$ we get $u_2=O(\epsilon)u_1$ for the small eigenvalues, and symmetrically for the large ones.
\end{proof}

\subsection{Zeros at $\theta=-c$ for $1\le c\le K$}\label{ssec:hz1}
\begin{proposition}\label{prop:hz1}
For $1\le c\le K$, $\operatorname{ord}_{\theta=-c}\kappa_{\mathsf H}\ge e_c$ (with $e_c$ as in \eqref{eq:hec}).
\end{proposition}
\begin{proof}
With $\theta=-c+\epsilon$: $x_a=c-a-\epsilon$, $x_a^0=c-a$, $W_0=\{c-K+1,\dots,c\}$; $\Psi_j$ has a simple zero at $\epsilon=0$ iff $(j+\theta)_K$ does iff $c-K+1\le j\le c$, i.e.\ (as $c\le K$) iff $j\in O:=\{1,\dots,\min(c,s)\}$. So $\mathsf H=\Lambda_\epsilon\tilde{\mathsf H}\Lambda_\epsilon$ with $\Lambda_\epsilon=\operatorname{diag}(\epsilon\text{ on }O,1)$ and $\tilde{\mathsf H}=\tilde\Psi A\tilde\Psi$ analytic. For $j\in O$, $a=c-j$ collides: $j-x_a=\epsilon$, and $\epsilon^2\nu_j\to\varrho_a^0$ exactly. The \emph{far} nodes are $F=\{c-K+1,\dots,0\}$, $|F|=K-c$; nodes of $W_0$ above $s$ lie in $Z$ (since $c\le K\le s+2m$) and have $\varrho_a=O(\epsilon)$; $\varrho^0>0$ on $O\cup F$ (there $\Pi_Z>0$).

\emph{Case $c<s$.} Here $J_1=O=\{1..c\}$, $J_2=\{c+1..s\}\ni g$. At $\theta=-c$, $\Omega_a$ has a double zero for $a<c$ (both Pochhammers contain the factor $a-c+i=0$ for $i=c-a\in[1,s]$), cancelling the double pole of the Cauchy factors, and $\Omega_a>0$ for $a\ge c$; the columns $j>c$ of $C_G$ have $y_a+j=a-c+j>0$ for all $a\ge0$. Hence $\tilde{\mathsf H}_{22}(0)=\tilde\Psi\bigl(\sum_{a\ge c}\Omega_ac_ac_a^\top\bigr)_{>c,>c}\tilde\Psi$ is positive definite (Cauchy columns $1/(a-c+j)$, $j\in(c,s]$, $a\in[c,K)$, have full rank $s-c\le K-c$). Lemma~\ref{lem:hsch}: exactly $c$ eigenvalues are $O(\epsilon^2)$ and their eigenvectors have $u_{ig}=O(\epsilon)$ ($g\in J_2$): $\sum\operatorname{ord}u_{ig}\ge c$.
\emph{Case $c\ge s$.} Then $O=G$, $\mathsf H=\epsilon^2\tilde{\mathsf H}$: all $s$ eigenvalues are $O(\epsilon^2)$, and $\sum\operatorname{ord}u_{ig}\ge0$.
\emph{The cluster.} Put $\mathcal N(\epsilon)=\epsilon^2\langle\cdot,\cdot\rangle_\nu$, $\mathcal R(\epsilon)=\langle\cdot,\cdot\rangle_\varrho$ on $\mathrm{Poly}_{<s}$, both analytic at $0$. $\mathcal R_0(P)\ge\sum_{O\cup F}\varrho^0P^2$ is positive definite ($|O\cup F|=\min(c,s)+K-c\ge s$), and $\mathcal N_0(P)=\sum_{j\in O}\varrho^0_{c-j}P(j)^2\le\mathcal R_0(P)$ with equality iff $P|_F=0$. With the analytic Cholesky factorisation $\mathcal R=L^\top L$, $\mathcal M=L^{-\top}\mathcal NL^{-1}$ is symmetric analytic, $\mathcal M_0$ has maximal eigenvalue $1$ with eigenspace $\mathcal U_0=\{P:P|_F=0\}$, of dimension $s-(K-c)=c-m$ (if $c\ge m$). The eigenvalues $\mu$ of $\mathcal M$ and $\lambda$ of $\mathsf H$ satisfy $\lambda=\epsilon^2/\mu$.

\emph{Residue identity.} For $P,Q\in\mathcal U_0$, $P^\top(\mathcal N_1-\mathcal R_1)Q=0$. Indeed with $F_0=PQ$ (double zeros on $F$), $W_0(x)=\prod_b(x-x_b^0)$ and $\mathcal Y=\Pi_ZF_0/W_0^2$: at a collision node $x_a^0=j$, $\frac{d}{d\epsilon}[\epsilon^2\nu_jF_0(j)-\varrho_aF_0(x_a)]_0=(\Pi_ZF_0/\prod_{b\ne a}(x-x_b^0)^2)'(j)=\operatorname{Res}_{x=j}\mathcal Y$ (double pole); at a $Z$-node $z=x_a^0$: $-\frac d{d\epsilon}\varrho_aF_0(x_a)=\Pi_Z'(z)F_0(z)/\prod_{b\ne a}(z-x_b^0)^2=\operatorname{Res}_z\mathcal Y$ (simple pole); far nodes and non-colliding $\nu$-nodes contribute $O(\epsilon^2)$ and $\mathcal Y$ has no pole there. So the sum is the sum of all finite residues of $\mathcal Y$, which vanishes because $\deg\mathcal Y\le2m+2(s-1)-2K=-2$. Since $\mathcal N_0(P,\cdot)=\mathcal R_0(P,\cdot)$ for $P\in\mathcal U_0$, this is $P_{\mathcal U_0}\mathcal M_1P_{\mathcal U_0}=0$. Lemma~\ref{lem:hclu} gives $c-m$ eigenvalues $\mu=1+O(\epsilon^2)$, i.e.\ $\lambda=\epsilon^2(1+O(\epsilon^2))$, in a window of width $O(\epsilon^4)$ (these are among the $\min(c,s)$ small eigenvalues).

\emph{Counting.} By Lemma~\ref{lem:hRell}: for $c<s$, $\operatorname{ord}\kappa_{\mathsf H}\ge c+2\bigl[\binom c2-\binom{c-m}2\bigr]+4\binom{c-m}2=c+c(c-1)+(c-m)(c-m-1)$; for $s\le c\le K$, $\ge s(s-1)+(c-m)(c-m-1)$ (pairs among the $s$ small eigenvalues give $2$, pairs inside the cluster $4$). This is \eqref{eq:hec}.
\end{proof}

\subsection{Zeros at $\theta=-c$ for $K<c\le2K-1$}\label{ssec:hz2}
Write $c=K+d'$, $1\le d'\le K-1$, $\epsilon=\theta+c$. Now $W_0=\{d'+1,\dots,c\}$ splits into $O=W_0\cap V=\{d'+1,\dots,s\}$ ($|O|=o$), $Z_c=W_0\cap Z$ and $F=(s+2m,c]$, $|F|=(d'-m)_+$; the free $\nu$-nodes are $Q=V\setminus O=\{1,\dots,\min(d',s)\}$, and $\Psi_j$ vanishes at $\epsilon=0$ exactly for $j\in O$. Put $\mathcal K_0=\{P\in\mathrm{Poly}_{<s}:P|_{O\cup F}=0\}$, $\dim\mathcal K_0=r_0:=s-o-|F|$ ($o+|F|\le s$), $n:=o+r_0=s-|F|$. Since $x_a-x_b$ does not depend on $\epsilon$,
\begin{equation}\label{eq:hexact}
 A(\epsilon)(P,Q)=\sum_ac_a\,\mathcal G(x_a^0-\epsilon),\quad\mathcal G=\Pi_ZPQ,\ c_a=\prod_{b\ne a}(x_a^0-x_b^0)^{-2},
\end{equation}
with $A=A_0+\epsilon A_1+\epsilon^2A_2+\dots$ (Taylor coefficients of this exact form).
\begin{lemma}[residue facts]\label{lem:hres}
For $P_0\in\mathcal K_0$ and arbitrary $R\in\mathrm{Poly}_{<s}$: $A_0(P_0,R)=0$ and $A_1(P_0,R)=0$.
\end{lemma}
\begin{proof}
$A_0$: $P_0$ vanishes on $O\cup F$ and $\varrho^0=0$ on $Z_c$. $A_1=-\sum_ac_a\mathcal G'(x_a^0)$ with $\mathcal G=\Pi_ZP_0R$. The function $\Phi=\Pi_ZP_0R/W_0^2$ has at most a simple pole at each node of $W_0=O\cup Z_c\cup F$ ($P_0$ cancels one factor at $O\cup F$, $\Pi_Z$ at $Z_c$) with residue $c_a\mathcal G'(x_a^0)$, and $\deg\Phi\le2m+2(s-1)-2K=-2$; the residues sum to $0$.
\end{proof}
\begin{proposition}[L1]\label{prop:hL1}
Exactly $n$ eigenvalues of $\mathsf H(\epsilon)$ are $O(\epsilon^2)$; the others tend to nonzero limits, and their eigenvectors have components $O(\epsilon)$ on $O\oplus\mathcal Q_0$ (defined in the proof). If $d'\ge s$ then $o=0$ and $\operatorname{ord}\kappa_{\mathsf H}\ge n(n-1)$.
\end{proposition}
\begin{proof}
$\mathsf H=\Lambda_\epsilon\tilde{\mathsf H}\Lambda_\epsilon$ ($\Lambda_\epsilon=\epsilon$ on $O$); $\tilde{\mathsf H}(0)\succeq0$ since $A_0\succeq0$. Work in value coordinates $f\leftrightarrow P_f$. The kernel $\mathcal Q_0\subseteq\R^Q$ of $\tilde{\mathsf H}_{QQ}(0)$ is $\tilde\Psi(0)^{-1}\mathcal K_0$: for $f$ supported on $Q$, $P_f$ vanishes on $O$, and $f^\top A_0f=\sum_{O\cup F}\varrho^0P_f^2=0$ iff $P_f|_F=0$. As $\tilde{\mathsf H}(0)\succeq0$, the rows of $\mathcal Q_0$ in $\tilde{\mathsf H}(0)$ vanish. For $f,f'\in\mathcal Q_0$, expanding $(\tilde\Psi(\epsilon)f)^\top A(\epsilon)(\tilde\Psi(\epsilon)f')$ to first order and using Lemma~\ref{lem:hres} (with $P_0=P_{\tilde\Psi(0)f}\in\mathcal K_0$) gives $\tilde{\mathsf H}_{\mathcal Q_0\mathcal Q_0}=O(\epsilon^2)$. So $\mathsf H=\Gamma\hat{\mathsf H}\Gamma$ with $\Gamma=\epsilon$ on $O\oplus\mathcal Q_0$, $1$ on $\mathcal Q_1=\mathcal Q_0^\perp\cap\R^Q$, $\hat{\mathsf H}$ analytic and $\hat{\mathsf H}_{\mathcal Q_1\mathcal Q_1}(0)=\tilde{\mathsf H}_{\mathcal Q_1\mathcal Q_1}(0)$ positive definite. Apply Lemma~\ref{lem:hsch}: $|O|+\dim\mathcal Q_0=o+r_0=n$ small eigenvalues. If $d'\ge s$, $O=\emptyset$ and all pairs of small eigenvalues have $\operatorname{ord}\ge2$: $2\binom n2=n(n-1)$.
\end{proof}
For $d'<s$ two more facts are needed (here $g=s\in O$). Fix a complement $\mathcal K_0^c=C_1\oplus C_2$ of $\mathcal K_0$ with $C_1=\{P:P|_F=0\}/\mathcal K_0$, $C_2=\{P:P|_O=0\}/\mathcal K_0$ ($\dim C_1=o$, $\dim C_2=|F|$), let $D_\epsilon(P_\perp+P_0)=P_\perp+\epsilon^{-1}P_0$ and $\mathcal A=A\circ D_\epsilon$, $\mathcal N=\epsilon^2N\circ D_\epsilon$ (analytic by Lemma~\ref{lem:hres}). At $\epsilon=0$, in the value coordinates on $O\cup F$ for $\mathcal K_0^c$,
\[
 \mathcal A_0=\operatorname{diag}(\varrho^0_O,\varrho^0_F)\oplus A_2|_{\mathcal K_0},\qquad\mathcal N_0=\operatorname{diag}(\varrho^0_O,0)\oplus N_Q|_{\mathcal K_0},\qquad N_Q(P_0)=\sum_{j\in Q}\nu_j^0P_0(j)^2 ,
\]
the cross terms vanishing by Lemma~\ref{lem:hres}(b). Thus the pencil $(\mathcal N,\mathcal A+\mathcal N)$ has at $\epsilon=0$: eigenvalue $\mu'=\frac12$ on $C_1$, $0$ on $C_2$ and the eigenvalues of $N_Q(A_2+N_Q)^{-1}$ on $\mathcal K_0$.

\textbf{Gap condition (G$'$)}: $A_2+N_Q$ is positive definite on $\mathcal K_0$ and no eigenvalue $\tau$ of $A_2N_Q^{-1}$ equals $1$.
\begin{lemma}[L2, L3 under (G$'$)]\label{lem:hL23}
Assume (G$'$). Then (L2) $o$ of the small eigenvalues of $\mathsf H$ lie in a window of width $O(\epsilon^4)$; (L3) $\sum_i\operatorname{ord}u_{ig}\ge d'$.
\end{lemma}
\begin{proof}
(L2) For $\epsilon\ne0$ the pencils $(\mathcal N,\mathcal A)$ and $(\mathcal N,\mathcal A+\mathcal N)$ have the same eigenvectors and $\mu=\mu'/(1-\mu')$. Under (G$'$), $\mathcal A_0+\mathcal N_0\succ0$ so $\mathcal M'=L'^{-\top}\mathcal NL'^{-1}$ is analytic with eigenvalue $\frac12$ of eigenspace exactly $C_1$ (isolated: the $\mathcal K_0$-eigenvalues are $>0$ and $\ne\frac12$ by $\tau\ne1$). For $P,Q\in C_1$, $D_\epsilon$ is the identity and the first-order term is $\frac12P^\top(\mathcal N_1-\mathcal A_1)Q$, which vanishes by the residue identity of \S\ref{ssec:hz1} (verbatim: $PQ$ has double zeros on $F$, double poles of $\Pi_ZPQ/W_0^2$ at $O$, simple poles at $Z_c$). Lemma~\ref{lem:hclu}: $o$ eigenvalues $\mu'=\frac12+O(\epsilon^2)$, hence $\mu=1+O(\epsilon^2)$ and $\lambda=\epsilon^2/\mu$ lies in a window of width $O(\epsilon^4)$. (L3) The eigenvectors of $\mathsf H$ correspond to $\nu$-orthonormal eigenpolynomials $\hat P$ and $u_{ig}=\pm\sqrt{\nu_g}\,\hat P(g)\,\omega'(g)/\text{(const)}$. For the $r_0$ branches whose limit pencil eigenvalue lies in $\mathcal K_0$'s spectrum (separated from $\{0,\frac12\}$ by (G$'$)), the group projector shows $\hat P\in\mathcal K_0+O(\epsilon)$; as $\hat P_0(g)=0$ ($g\in O$), $\hat P_\perp(g)=O(\epsilon)$ and $\mathcal N_0\ge c>0$ there, so $u_{ig}=\sqrt{\nu_g}\epsilon\hat P_\perp(g)/\mathcal N(\hat P)^{1/2}=O(\epsilon)$ ($\nu_g\sim\epsilon^{-2}$). The $|F|$ large-eigenvalue branches have $u_{ig}=O(\epsilon)$ by Proposition~\ref{prop:hL1}. Together $|F|+r_0=d'$ branches. (To avoid ambiguity at coinciding eigenvalues, one applies the argument to the spectral projector of each isolated group: if $\Pi(0)e_g=0$ then $\sum_{\rm group}u_{ig}^2=\|\Pi(\epsilon)e_g\|^2=O(\epsilon^2)$ and every branch of the group has $\operatorname{ord}u_{ig}\ge1$.)
\end{proof}
\textbf{Removing (G$'$): deformation of the free $\nu$-nodes.} (G$'$) is true for $K\le18$ with a margin shrinking like $3^{-K}$ (exact computation), but we do not need it. Keep the $\varrho$-side and the collision nodes $O$ and replace the free $\nu$-nodes $j\in Q=\{1..d'\}$ by parameters $y_1,\dots,y_{d'}$: $\nu_j:=\Pi_Z(y_j)/W_\theta(y_j)^2$, $A_{jj'}:=\sum_a\varrho_aL_j(x_a)L_{j'}(x_a)$ ($L_j$ the Lagrange basis on all $s$ nodes $Q\cup O$), $h_n:=((\mathcal N^{-1}A)^n)_{gg}$ and $\mathbb K(\theta;y):=\det[h_{i+j}]_{i,j<s}$. Then: (i) $\mathbb K\in\Q[y][1/\Delta][\theta]$ with $\Delta=\prod_{i<j}(y_i-y_j)\prod_{i,o\in O}(y_i-o)\prod_i\Pi_Z(y_i)$ (the denominators of the Lagrange basis and of $\nu_j^{-1}=W_\theta(y_j)^2/\Pi_Z(y_j)$; the numerators are polynomial in $\theta$ since $x_a-x_b$ is constant); (ii) at $y=(1,\dots,d')$, $\Delta\ne0$ and $\mathbb K=\kappa_{\mathsf H}^2$ (the Hankel matrix of the Krylov sequence is $\mathcal K^\top\mathcal K$ with $\mathcal K$ the Krylov matrix of the symmetric matrix $\mathcal N^{-1/2}A\mathcal N^{-1/2}$, whose powers have the same $gg$-entries as those of $\mathcal N^{-1}A$); (iii) the argument above (L1, L2, L3 and Lemma~\ref{lem:hRell}) uses only: the $\varrho$-side facts (Lemma~\ref{lem:hres}, exact form \eqref{eq:hexact}), the collision data ($\epsilon^2\nu_j\to\varrho^0$ for $j\in O$, $\Lambda_\epsilon=\epsilon$ on $O$), and for the free nodes only that $y_j\notin W_0$, $\Pi_Z(y_j)>0$, distinctness; all these are open conditions. So for $y$ in a nonempty open set $\mathcal Y\subset\R^{d'}$ on which (G$'$) holds we get $\operatorname{ord}_{\theta=-c}\mathbb K(\theta;y)\ge2e_c$. (G$'$) holds on a nonempty open set: let $\mathcal Z_c=W_0\cap Z$ ($|\mathcal Z_c|=m+\min(d',m)\ge r_0$, $d'\ge r_0$) and take $y_i=z_i\pm\delta$ for $r_0$ distinct $z_i\in\mathcal Z_c$ (sign chosen so that $\Pi_Z(y_i)>0$) and the other $y_j$ generic below $1$. Then $\nu_i^0=\Pi_Z(y_i)/W_0(y_i)^2\sim|\Pi_Z'(z_i)|/(W_0'(z_i)^2\delta)\to\infty$, and $P_0\mapsto(P_0(z_i))_{i<r_0}$ is injective on $\mathcal K_0=\Pi_{O\cup F}\mathrm{Poly}_{<r_0}$ (as $\Pi_{O\cup F}(z_i)\ne0$), so $N_Q\ge(c'/\delta)I-O(1)$ on $\mathcal K_0$ while $A_2$ does not depend on $y$; hence $A_2+N_Q\succ0$ and $|\tau|=O(\delta)<1$. Finally, the Taylor coefficients of order $<2e_c$ of $\mathbb K$ at $\theta=-c$ are elements of $\Q[y][1/\Delta]$ which vanish on the open set $\mathcal Y$; their numerators are polynomials vanishing on a nonempty open subset of $\R^{d'}$, hence identically, in particular at $y=(1,\dots,d')$, where $\Delta\ne0$. Hence $\operatorname{ord}_{\theta=-c}\kappa_{\mathsf H}^2\ge2e_c$.
\begin{proposition}\label{prop:hz2}
For $K<c\le2K-1$, $\operatorname{ord}_{\theta=-c}\kappa_{\mathsf H}\ge e_c$.
\end{proposition}
\begin{proof}
If $d'\ge s$: Proposition~\ref{prop:hL1}. If $d'<s$: by Lemma~\ref{lem:hRell}, $\operatorname{ord}\kappa_{\mathsf H}\ge2\binom n2+2\binom o2+d'$ (all pairs of small eigenvalues $\ge2$; the $o$-cluster $\ge4$; $\sum\operatorname{ord}u_{ig}\ge d'$), which is \eqref{eq:hec}; the deformation argument removes the need for (G$'$).
\end{proof}
(Numerical confirmation: exact evaluation of the Rellich counts for all $2\le K\le300$, $0\le m\le K-2$, $1\le c\le2K-1$ reproduces \eqref{eq:hec}; and the exact orders $\operatorname{ord}_{\theta=-c}\mathbb K(\theta;y)=2e_c$ at random rational $y$ were verified for $22$ triples $(K,m,c)$.)

\subsection{Degree and leading coefficient}\label{ssec:hlead}
Let $h_d(\varpi)$ denote the squared norm of the monic orthogonal polynomial $p_d$ of degree $d$ for a positive weight $\varpi$ on finitely many nodes (with at least $s$ nodes), and put the two weights
\begin{gather*}
 \varpi_a=\frac1{(a!\,(K-1-a)!)^2}\quad(a\in\{0,\dots,K-1\}),\\ v_j=\frac1{\omega'(j)^2\Pi_Z(j)}=\frac1{((j-1)!)^2(s-j)!\,(s+2m-j)!}\quad(j\in\{1,\dots,s\}).
\end{gather*}
\begin{theorem}\label{thm:hlead}
As $\theta\to+\infty$ the eigenvalues of $\mathsf H$ satisfy $\lambda^{(d)}=h_d(v)h_d(\varpi)\,\theta^{4K-2-4d}(1+O(1/\theta))$, $d=0,\dots,s-1$ (distinct orders), and the unit eigenvector of $\lambda^{(d)}$ has $u_{g}^2\to v_s\,p_d^{(v)}(s)^2/h_d(v)$. Consequently $\deg_\theta\kappa_{\mathsf H}=D$ exactly and the leading coefficient $\varkappa_{\mathsf H}$ satisfies
\[
 \varkappa_{\mathsf H}^2=\prod_{d<s}\frac{v_sp_d^{(v)}(s)^2}{h_d(v)}\ \prod_{d<s}\bigl(h_d(v)h_d(\varpi)\bigr)^{2(s-1-d)} .
\]
\end{theorem}
\begin{proof}
(1) \emph{Factorisation.} $\Psi_j=\theta^K\tilde\psi_j(1+O(1/\theta))$ with $\tilde\psi_j^2=v_j$ (indeed $\Psi_j^{-2}\sim\theta^{-2K}\omega'(j)^2\Pi_Z(j)$); $\varrho_a=\theta^{2m}\varpi_a(1+O(1/\theta))$. Expanding $y_a^{p}=(\theta+a)^p$, $\mathbb M=R\,\mathbb M_c R^\top$ with $R=D_0Z_0D_0^{-1}$, $D_0=\operatorname{diag}(\theta^p)$, $Z_0=[\binom pn]$, and $\mathbb M_c=[\sum_a\varrho_aa^{p+q}]=\theta^{2m}(L\Delta_wL^\top+O(1/\theta))$, where $L\Delta_wL^\top$ is the $LDL^\top$ factorisation of the moment matrix of $\varpi$, $(\Delta_w)_d=h_d(\varpi)>0$. Writing $\mathcal E=V_s\mathcal N$ ($\mathcal N J$ upper triangular with diagonal $\pm1$, $J$ the reversal), one finds $\mathsf H/\theta^{2m+2K}=B(I+F)B^\top$, $F=O(1/\theta)$, $B=\tilde\psi V_sD'YD_0^{-1}\Delta_w^{1/2}$, with $D'=\operatorname{diag}(\theta^{s-1-r})$, $Y=U_\theta JZ_0L_\theta=\operatorname{diag}(\pm1)JZ_0+O(1/\theta)$ ($U_\theta=D'^{-1}\mathcal NJD'$ and $L_\theta=D_0^{-1}LD_0$ have entries $O(1/\theta)$ off the diagonal). With $I+F=(I+F_1)(I+F_1)^\top$ the eigenvalues of $\mathsf H/\theta^{2m+2K}$ are $\sigma_i(B(I+F_1))^2$, and $\sigma_i(B(I+F_1))=\sigma_i(B)(1+O(1/\theta))$.
(2) \emph{Compounds.} $\Lambda^k(B(I+F_1))$: the scaling $\Lambda^k(D')\otimes\Lambda^k(D_0^{-1})$ is uniquely maximal at the index sets $\{0..k-1\}$, every other pair being smaller by a factor $\ge\theta$; $\det[\binom pn]_{p=s-k..s-1,\,n<k}=\pm1$. Hence $\Lambda^kB(I+F_1)=\theta^{e_k'}\bigl[\pm\bigl(\wedge_{r<k}\tilde\psi V_se_r\bigr)\otimes\prod_{d<k}h_d(\varpi)^{1/2}e_{J_0}^\top+O(1/\theta)\bigr]$, and $\|\wedge_{r<k}\tilde\psi V_se_r\|^2=\det(V_k^\top\tilde\psi^2V_k)=\prod_{d<k}h_d(v)$ (Gram determinant of the moments of $v$). Taking norms and ratios gives $\lambda^{(d)}$; the exponent of $\theta$ is $2(s-1-2d)+2m+2K=4K-2-4d$.
(3) \emph{Eigenvectors.} The compounds are rank one at leading order with scale gaps $\theta^2$, so (Wedin) the top-$k$ eigenspace of $\mathsf H$ converges to $\operatorname{span}(\tilde\psi V_k)$ [apply this to $\Lambda^k(B(I+F_1))$, whose left singular spaces are the eigenspaces of $\mathsf H$]. The $d$-th eigenvector tends to $\tilde\psi\,p_d^{(v)}(j)/h_d(v)^{1/2}$ (Gram--Schmidt of $\tilde\psi j^r$), hence the limit of $u_g^2$. No limit vanishes: if $p_d(s)=0$, put $q=p_d/(x-s)$; $q\ne0$ at some node $j<s$, and $0=\langle p_d,q\rangle_v=\sum_jv_j(j-s)q(j)^2<0$.
(4) By Proposition~\ref{prop:hH}'s symmetric form, $\kappa_{\mathsf H}=\pm\prod_iu_{ig}\prod_{i<i'}(\lambda_{i'}-\lambda_i)$ for the (real symmetric) $\mathsf H$ at each large $\theta$; the orders of $\lambda^{(d)}$ are strictly separated, so $\prod_{i<i'}|\lambda_i-\lambda_{i'}|=\prod_d(\lambda^{(d)})^{s-1-d}(1+O(\theta^{-4}))$, of degree $\sum_d(4K-2-4d)(s-1-d)=\frac{s(s-1)(4s+6m+1)}3=D$.
\end{proof}
\begin{proposition}[prime factors of $\varkappa$]\label{prop:hprimes}
$\varkappa_{\mathsf H}^2$, and hence $\varkappa=\varkappa_R=\Lambda_g^s(\det\Lambda)^{-1}\varkappa_{\mathsf H}$, $\varkappa_L$, are rational numbers whose prime factors are all $\le2K-1$.
\end{proposition}
\begin{proof}
Both weights are Hahn weights with negative parameters: using $\binom{\alpha+x}x\binom{\beta+N-x}{N-x}=(-1)^N\binom{-\alpha-1}x\binom{-\beta-1}{N-x}$, $\varpi_a=\pm((K-1)!)^{-2}\binom{K-1}a^2$ is the Hahn weight with $\alpha=\beta=-K$, $N=K-1$ (in the variable $x=a$), and $v_j=\pm\frac1{(s+2m-1)!(s-1)!}\binom{s+2m-1}{j-1}\binom{s-1}{s-j}$ the Hahn weight with $\alpha=-s-2m$, $\beta=-s$, $N=s-1$ (variable $x=j-1$, the endpoint $j=s$ being $x=N$). In both cases $\alpha+\beta+1=-2K+1$. We use the Hahn formulas of \cite[\S9.5]{KLS} at these negative parameters, with the following justification. For fixed $N$ and $n$, the orthogonality relation (with its norm $h_n$), the leading coefficient and the endpoint value $Q_n(N)=(-1)^n(\beta+1)_n/(\alpha+1)_n$ are identities between rational functions of $(\alpha,\beta)$ (the weight $\binom{\alpha+x}x\binom{\beta+N-x}{N-x}$ is a polynomial in $(\alpha,\beta)$); they hold on the open set $\alpha,\beta>-1$, hence identically, hence at every point at which no denominator vanishes. Here the quantities that must be nonzero for $0\le n\le s-1$ (denominators, and the factors of $h_n$ and of the leading coefficient) are the factors of $(\alpha+1)_k$, $(\beta+1)_k$, $(-N)_k$ ($k\le n$), $2n+\alpha+\beta+1=2n-2K+1$ (odd), $(n+\alpha+\beta+1)_n$ and $(n+\alpha+\beta+1)_{N+1}$ (last factor $n-K\le-1$ resp.\ $n-2K+s\le-2m-1$), all nonzero integers of absolute value $\le2K-1$ (the largest is $|n+\alpha+\beta+1|$ at $n=0$). Since both weights are positive, the monic orthogonal polynomials are unique and equal the Hahn polynomials divided by their leading coefficients, which are nonzero. All remaining constants are factorials of integers $\le K+m-1\le2K-3$ and $\Lambda_g^2=\Pi_X(s)/\Pi_Y(s)$, a ratio of integers of absolute value $<K+m$. This gives the claim.
\end{proof}
(Checked in exact arithmetic: the formula for $\varkappa_{R}^2,\varkappa_L^2$ agrees with the leading coefficient obtained by interpolation of $k_R,k_L$ at $D+2$ points, for $20$ pairs $(K,m)$ up to $(8,6)$, including $(3,1)$: $\varkappa_R=\frac12,\varkappa_L=\frac38$; and all prime factors are $\le2K-3$ for $K\le18$.)

\subsection{The sum $\sum_ce_c=D$}\label{ssec:hsum}
\begin{lemma}\label{lem:hsum}
$\sum_{c=1}^{2K-1}e_c=D$.
\end{lemma}
\begin{proof}
For $c\le K$: $\sum_{c<s}c=\frac{s(s-1)}2$; $\sum_{c\le K}T_2(\min(c,s))=\sum_{c\le s}c(c-1)+(K-s)s(s-1)=\frac{(s-1)s(s+1)}3+ms(s-1)$; $\sum_{c=m+1}^KT_2(c-m)=\sum_{n=1}^sn(n-1)=\frac{(s-1)s(s+1)}3$. For $c=K+d'$, $d'=1..K-1$: $n=s$ for $d'\le m$ and $n=K-d'$ for $d'>m$, so $\sum T_2(n)=ms(s-1)+\sum_{t=1}^{s-1}t(t-1)=ms(s-1)+\frac{(s-2)(s-1)s}3$; $\sum T_2(o)=\sum_{o=1}^{s-1}o(o-1)=\frac{(s-2)(s-1)s}3$; $\sum_{d'<s}d'=\frac{s(s-1)}2$. The total is $s(s-1)\bigl[1+2m+\frac{2(s+1)+2(s-2)}3\bigr]=\frac{s(s-1)(4s+6m+1)}3=D$.
\end{proof}

\subsection{Proof of Theorem~\ref{thm:hK}}
By Lemma~\ref{lem:hpoly} and Theorem~\ref{thm:hlead}, $k$ is a polynomial of degree exactly $D$ with leading coefficient $\varkappa$ (a rational with primes $\le2K-1$, Proposition~\ref{prop:hprimes}). By Propositions~\ref{prop:hz1}, \ref{prop:hz2} (applied to $\kappa_{\mathsf H}$, and $k_{R,L}=\text{const}\cdot\kappa_{\mathsf H}$ by Proposition~\ref{prop:hH}) $\prod_c(\theta+c)^{e_c}$ divides $k$, and $\sum e_c=D$ (Lemma~\ref{lem:hsum}) $=\deg k$. Hence $k=\varkappa\prod_c(\theta+c)^{e_c}$.

\begin{lemma}[from $\Q$ to $\F_\ell$]\label{lem:hbridge}
Let $\ell=4K-1$ be prime. Then $\hat\Gamma(0)\in M_K(\Z_{(\ell)})$ with $\det\hat\Gamma(0)\in\Z_{(\ell)}^\times$, $T(0)\in M_s(\Z_{(\ell)})$, $k_R(0),k_L(0)\in\Z_{(\ell)}$, and the reductions modulo $\ell$ are the Krylov determinants $k_R,k_L$ computed over $\F_\ell$ from $\bar T=(\bar{\hat\Gamma}^{-1})_{GG}$ in \S\ref{sec:hnv}.
\end{lemma}
\begin{proof}
$\det\hat\Gamma(0)=\det C_{G\cup X}\det C_{G\cup Y}$, with Cauchy factors $a'-a\le K-1$, $j'-j\le K+m-1\le2K-3$ and $a+j\le2K+m-1\le3K-3<\ell$; all are $\ell$-units. The entries $1/(a+j)$ are $\ell$-integral, so $\hat\Gamma(0)$ and its inverse (adjugate over a unit) are $\ell$-integral; reduction is a ring homomorphism and commutes with products, submatrices, inverses of unit-determinant matrices and determinants.
\end{proof}
By Lemma~\ref{lem:hbridge} and the product formula, $k_R(0)=\varkappa\prod_{c\le2K-1}c^{e_c}$ is an $\ell$-unit because $\varkappa$ and all $c\le2K-1<\ell$ are prime to $\ell$ ($2K\le\ell-1$). This is Theorem~\ref{thm:hK}. \qed

\medskip\noindent\emph{Proof of Proposition~\ref{prop:hNV}.} Theorem~\ref{thm:hNV} together with Theorem~\ref{thm:hK} (for $0\le m\le K-2$) and the trivial case $m=K-1$ ($k_R=k_L=1$), for every prime $\ell\equiv3\pmod 4$, $\ell\ge7$ (for $\ell=7$, $K=2$ and $m\in\{0,1\}$). \qed

\section{The case $\ell\mid\den$ and the block $R_{X'Y'}$}\label{sec:hl3}

This section proves Proposition~\ref{prop:hL3} and the statement $\det\Rk[X',Y']\ne0$ used in Lemma~\ref{lem:hschur}. Notation and hypotheses are those of \S\ref{sec:hnv}: $\ell\equiv3\ (4)$, $K=(\ell+1)/4$, $1\le m\le K-1$, $s=K-m$, a geometry $\mathcal G=\mathcal G(r_1,m,e)$ with (H); in $j$-coordinates $J_{X'}=s+\{2,4,\dots,2m\}$, $J_{Y'}=s+\{1,3,\dots,2m-1\}$. The new ingredient is a closed-form determinant.

\begin{lemma}[high--high block mod $\ell$]\label{lem:hhh}
For distinct high nodes $x,y$: $\ \Rk(x,y)\equiv\frac1{16}\sum_{a<K}\frac1{(a+j_x)(a+j_y)}=\frac1{16}W_{j_x}^\top W_{j_y}\pmod\ell$. Hence $\Rk[X',Y']\equiv\frac1{16}W_{X'}^\top W_{Y'}$.
\end{lemma}
\begin{proof}
Let $K\le a<b\le\ell-1$. For $k\in(a,b]$, $K<k\le\ell-1$ and $4k-1\equiv4(k-K)$ (as $4K=\ell+1$), so $\frac1{16k(4k-1)}\equiv\frac1{64k(k-K)}=\frac1{64K}\bigl(\frac1{k-K}-\frac1k\bigr)$, all of $k,k-K,b-a$ lying in $[1,\ell-1]$. Summing, $64K(b-a)\Rk(a,b)\equiv(H_{b-K}-H_b)-(H_{a-K}-H_a)$. By $H_n\equiv H_{\ell-1-n}$ ($1/(\ell-i)\equiv-1/i$, $H_{\ell-1}\equiv0$) with $j=\ell-b$: $H_b\equiv H_{j-1}$, $H_{b-K}\equiv H_{j+K-1}$, so $H_{b-K}-H_b\equiv F(j_b):=\sum_{a'<K}\frac1{a'+j_b}$. As $b-a=j_a-j_b$, $\Rk(a,b)\equiv\frac{F(j_b)-F(j_a)}{64K(j_a-j_b)}=\frac1{64K}\sum_{a'}\frac1{(a'+j_a)(a'+j_b)}$ and $64K=16(\ell+1)\equiv16$.
\end{proof}

For $\theta$ in a field of characteristic $0$ (or an integer with all entries invertible mod $\ell$) put $X_\theta=\theta+\{2,4,\dots,2m\}$, $Y_\theta=\theta+\{1,3,\dots,2m-1\}$ and $\Delta_{K,m}(\theta)=\det\bigl(C_{X_\theta}^\top C_{Y_\theta}\bigr)=\det\bigl[\sum_{a<K}\frac1{(a+x)(a+y)}\bigr]_{x\in X_\theta,y\in Y_\theta}$. Let $V(Z)=\prod_{z<z'}(z'-z)$, $\operatorname{sf}(n)=\prod_{k=1}^nk!$ and $\mathsf H_{K,m}=\det\bigl[\binom K{p+q+1}\bigr]_{0\le p,q<m}$.
\begin{proposition}[the block determinant]\label{prop:hDelta}
\begin{multline*}
 \Delta_{K,m}(\theta)=(-1)^{m(m-1)/2}\,V(2,4,\dots,2m)\,V(1,3,\dots,2m-1)\\ \times\operatorname{sf}(m-1)^2\,\mathsf H_{K,m}\Big/\prod_{i=1}^m(\theta+i)_m(\theta+K+i)_m ,
\end{multline*}
and $(-1)^{m(m-1)/2}\operatorname{sf}(m-1)^2\mathsf H_{K,m}=\sum_{S\subset[0,K),|S|=m}V(S)^2=\prod_{n=0}^{m-1}h_n$, $h_n=\frac{(n!)^4}{(2n)!(2n+1)!}\prod_{i=-n}^n(K+i)$.
\end{proposition}
\begin{proof}
\emph{(1) Cauchy--Binet.} $\Delta=\sum_{S}\det C_X[S]\det C_Y[S]$ over $m$-subsets $S\subset[0,K)$. By Cauchy's formula $\det C_Z[S]=V(S)V(Z)/\prod_{a\in S,z\in Z}(a+z)$, and $\prod_{x\in X}(a+x)\prod_{y\in Y}(a+y)=(a+\theta+1)_{2m}$. So $\Delta=V(X)V(Y)\Phi$, $\Phi=\sum_SV(S)^2\prod_{a\in S}\varrho(a)$, $\varrho(a)=1/(a+\theta+1)_{2m}$.
\emph{(2) Hankel form.} By Cauchy--Binet again $\Phi=\det[\langle a^p,a^q\rangle]_{p,q<m}$ with $\langle f,g\rangle=\sum_{a=0}^{K-1}fg\varrho$, and a unitriangular change of the monic bases gives $\Phi=\det[\langle\phi_p,\psi_q\rangle]$, $\phi_p(a)=(a+\theta+1)_p$, $\psi_q(a)=(a+\theta+2m-q+1)_q$.
\emph{(3) Telescoping.} $\phi_p\psi_q\varrho=1/(a+\theta+p+1)_L$, $L=2m-p-q\ge2$, and $1/(c)_L=\frac1{L-1}\bigl[\frac1{(c)_{L-1}}-\frac1{(c+1)_{L-1}}\bigr]$ telescopes the sum over $a$. With $u_p=\theta+p+1$, $v_q=\theta+2m-q$, $f(z)=(z)_K$: $\langle\phi_p,\psi_q\rangle=\frac{\Gamma(u_p)}{\Gamma(v_q+K)}\frac{f(v_q)-f(u_p)}{v_q-u_p}$. So $\Phi=\prod_p\Gamma(u_p)\big/\prod_q\Gamma(v_q+K)\cdot\det[f[u_p,v_q]]$ (divided differences).
\emph{(4) Differences.} Reverse the order of rows and of columns (two signs cancel): rows $u=\theta+m-p$, columns $v=\theta+m+1+q$. Replace row $p$ by $(-1)^p\nabla_u^p$ and column $q$ by $\Delta_v^q$ (unitriangular, leading coefficient $1$). For unit spacing $\nabla^pg(x)=p!\,g[x-p,\dots,x]$, $\Delta^qh(y)=q!\,h[y,\dots,y+q]$, and iterated divided differences compose: the new $(p,q)$ entry is $(-1)^pp!q!\,f[\theta+m-p,\dots,\theta+m+1+q]=(-1)^pp!q!\,\Delta^nf(\theta+m-p)/n!$, $n=p+q+1$. Since $\Delta(z)_k=k(z+1)_{k-1}$, $\Delta^n(z)_K=\frac{K!}{(K-n)!}(z+n)_{K-n}$, so the entry is $(-1)^pp!q!\binom Kn\frac{\Gamma(\theta+m+K-p)}{\Gamma(\theta+m+q+1)}$ (also correct for $n>K$, both sides vanishing). Taking out the row and column factors, $\det[f[u,v]]=(-1)^{m(m-1)/2}\operatorname{sf}(m-1)^2\mathsf H_{K,m}\prod_p\Gamma(\theta+m+K-p)/\prod_q\Gamma(\theta+m+q+1)$.
\emph{(5) Collecting} the Gamma factors: $\prod_p\Gamma(\theta+p+1)/\prod_q\Gamma(\theta+m+q+1)=\prod_{i=1}^m1/(\theta+i)_m$ and $\prod_p\Gamma(\theta+m+K-p)/\prod_q\Gamma(\theta+2m-q+K)=\prod_{i=1}^m1/(\theta+K+i)_m$. This is the identity. For the constant, multiply by $\theta^{2m^2}$ and let $\theta\to\infty$: each Cauchy--Binet term is $V(X)V(Y)V(S)^2\theta^{-2m^2}(1+O(1/\theta))$, giving $\sum_SV(S)^2=\det[\sum_aa^{p+q}]$ (Hankel determinant of the moments of the uniform weight on $\{0..K-1\}$), which equals $\prod_nh_n$ with $h_n$ the squared norm of the monic orthogonal (discrete Chebyshev, Hahn $\alpha=\beta=0$) polynomial of degree $n$ \cite[\S9.5]{KLS}.
\end{proof}
(Checked exactly for $K\le12$, $1\le m\le K+1$ and several non-integer $\theta$ ($540$ cases); the factorisation of $\Delta(\theta)$ over $\Q[\theta]$ was also computed independently.)

\begin{corollary}\label{cor:hRXY}
Let $\theta=s=K-m$, $1\le m\le K-1$. Then $\Delta_{K,m}(s)$, computed in $\F_\ell$ from the Cauchy entries, is nonzero; hence $\det(W_{X'}^\top W_{Y'})=\pm\Delta_{K,m}(s)\in\F_\ell^\times$ and $\det\Rk[X',Y']\equiv16^{-m}\det(W_{X'}^\top W_{Y'})\not\equiv0$. In particular $\det\Rk[X',Y']\neq0$, $v_\ell(\det\Rk[X',Y'])=0$ and $\Rk[X',Y']^{-1}\in M_m(\Z_{(\ell)})$.
\end{corollary}
\begin{proof}
The identity of Proposition~\ref{prop:hDelta} holds in $\Z_{(\ell)}$ at the integer $\theta=s$ (all entries $\frac1{(a+x)(a+y)}$ have $a+x\le K-1+s+2m=2K+m-1<\ell$ and are $\ell$-integral). The numerator is, by the second part, a product of integers $2(i'-i)<2m$, $k!$ with $k\le2m-1$, and $K+i$ with $|i|\le m-1$, all in $[1,\ell-1]$ (as $m\le K<\frac{\ell+1}2$), divided by factorials $(2n)!(2n+1)!$ with $2n+1\le2m-1<\ell$; the denominator is a product of integers in $[s+1,s+K+2m-1]\subseteq[1,3K-2]$, $3K-2<\ell$. So $\Delta$ is an $\ell$-unit. Finally use Lemma~\ref{lem:hhh} (no diagonal entries occur in $X'\times Y'$).
\end{proof}
(Direct exact computation of $\det(C_X^\top C_Y)$ modulo $\ell$ for every prime $\ell\equiv3\ (4)$, $\ell<200$, all $m$ and $s\in\{K-m-1,K-m,K-m+1\}$, and of the factor ranges for $\ell<3000$ ($223{,}320$ cases) gave no failure.)

\begin{proof}[Proof of Proposition~\ref{prop:hL3}]
Let $\mathcal G=\mathcal G(r_1,m,e)$, $n_C=|C|$, $x_0=u'/\ell^{e_0}$ with $u'=\ell^{e_0}x_0\in\Z_{(\ell)}^\times$ ($e_0\ne2$). By Lemma~\ref{lem:hschur} and Corollary~\ref{cor:hRXY} it suffices to show $P(x_0)\ne0$ for the monic polynomial $P(X)=\det(XI+M)$, $M=\Rk[C,C]-\Rk[C,Y']\Rk[X',Y']^{-1}\Rk[X',C]$. By Lemma~\ref{lem:hstruct}, $v_\ell(\Rk[C,C])\ge-2$ and $v_\ell(\Rk[C,Y']),v_\ell(\Rk[X',C])\ge-1$ (no diagonal entries occur, since $C,X',Y'$ are disjoint; low--high entries have valuation $-1$, the rest $\ge0$). By Corollary~\ref{cor:hRXY} $\Rk[X',Y']^{-1}$ is $\ell$-integral, so $\ell^2M\in M(\Z_{(\ell)})$.

\emph{$e_0\ge3$.} $\ell^{e_0n_C}P(u'\ell^{-e_0})=\det(u'I+\ell^{e_0-2}\cdot\ell^2M)\equiv u'^{\,n_C}\not\equiv0$.

\emph{$e_0=1$ (asymmetric scaling).} Let $N=\operatorname{diag}(\ell I_{C\sqcup X'})\Fu(u'/\ell)\operatorname{diag}(I_L,\ell I_H,I_{Y'})$ ($L$ low, $H=C_{\rm high}$): scale all rows by $\ell$ and the $H$ columns by $\ell$, leaving the columns of $L$ and $Y'$ unscaled. By Lemmas~\ref{lem:hstruct} and~\ref{lem:hhh}, $N\in M(\Z_{(\ell)})$ and modulo $\ell$ (rows $L,H,X'$; columns $L,H,Y'$)
\[
 N\equiv\begin{pmatrix}u'I&0&-\frac14W_{Y'}\\ \ast&-I&0\\-\frac14W_{X'}^\top&0&0\end{pmatrix}
\]
(the entries $\ell\Rk[H,Y']$ and $\ell\Rk[X',Y']$ are $\equiv0$ because off-diagonal high--high entries are integral; $\ell^2\Rk[L,H]\equiv\ell^2\Rk[X',H]\equiv0$). The $H$ columns vanish outside the $H$ rows; expanding, $\det N\equiv\pm\det\begin{psmallmatrix}u'I_K&-\frac14W_{Y'}\\-\frac14W_{X'}^\top&0\end{psmallmatrix}=\pm16^{-m}u'^{\,K-m}\det(W_{X'}^\top W_{Y'})\ne0$ (Schur complement of $u'I_K$; Corollary~\ref{cor:hRXY}). Since $\det N=\ell^{(n_C+m)+\#H}\det\Fu(u'/\ell)$, $v_\ell(\det\Fu(x_0))=-(2n_C-K+m)$.
\end{proof}
\begin{remark}[Newton polygon; not used]
Write $P(X)=\sum a_kX^k$ for a member, $k_0=K-m$, $v=2(r-K+m)$ (with $r=|C|$) and $b_k=v_\ell(a_k)+v$. Then $b_k\ge2(k-k_0)$ for all $k$ (the coefficient $a_k$ is $\pm$ a sum of principal $(r-k)$-minors of $M$ and $\ell^2M$ is integral), $b_k\ge0$ for all $k$ with $b_{k_0}=0$ (this is Proposition~\ref{prop:hred}: $\ell^{2(\#C_{\rm high}+m)}\det\Fu(X)\in\Z_{(\ell)}[X]$, whose reduction $\pm16^{-K}B(H^{(2)}_{K-1}-16X)$ has exact degree $K-m$ because its leading coefficient is $\pm\det(W_{X'}^\top W_{Y'})\ne0$; this is the place where the bound $b_k\ge0$ for $k<k_0$ must be taken from Proposition~\ref{prop:hred}, not from the $e_0=1$ argument), and $b_r=2(r-k_0)$. Hence the Newton polygon of $P$ has the slope-$2$ segment from $(k_0,0)$ to $(r,2(r-k_0))$ on $[k_0,r]$, which is another way to see Proposition~\ref{prop:hL3}: the minimum of $b_k-e_0k$ is unique at $k_0$ for $e_0=1$ and at $r$ for $e_0\ge3$, and all points tie for $e_0=2$.
\end{remark}
\begin{remark}
For $e_0=2$ the same computation mixes $u'I$ with the block $-I$ of $H$ and an exceptional residue of $u'$ can occur (as in Lemma~\ref{lem:3}); it is avoided by the choice of $\ell$ with $v_\ell(q)\ne2$. The scaling of Lemma~\ref{lem:3} cannot be used for $m\ge1$: scaling the $Y'$ columns by $\ell$ kills the entries $\ell\Rk[L,Y']$ which carry the rank.
\end{remark}

\section{Proof of Theorem~\ref{thm:hybrid}, status, and reproducibility}\label{sec:hfinal}

\subsection{The proof}
\begin{proof}[Proof of Theorem~\ref{thm:hybrid}]
Propositions~\ref{prop:hC} (\S\ref{sec:hcontent}), \ref{prop:hD} (\S\ref{sec:hdecay}), \ref{prop:hNV} (\S\ref{sec:hnv}, \S\ref{sec:hkrylov}) and~\ref{prop:hL3} (\S\ref{sec:hl3}) give the inputs of Theorem~\ref{thm:hmeasure}. With the certified numbers
\begin{gather*}
 D_F\ge3.90626950661,\qquad c_{\rm hyb}\le7.65406474415,\\ \Lambda_F=\tfrac{313}{105}+\tfrac{554}{840}+\tfrac{118/21+2\log2}{64}-\tfrac{129}{4480}=3.72114001606\ (\le3.7211401),
\end{gather*}
we get $\eta_{\rm hyb}=D_F-\Lambda_F\ge3.9062695-3.7211401=0.1851294$ (indeed $0.18512949$) and, rounding every constant in the direction that increases the bound,
\begin{multline*}
 c_{\rm hyb}\Bigl(\frac1{\eta_{\rm hyb}}+\frac23\Bigr)\le7.6540648\Bigl(\frac1{0.1851294}+\frac23\Bigr)\\=7.6540648\times6.06829\ldots=46.4471\ldots<46.45 .
\end{multline*}
(With the unrounded values the exponent is $46.44709$.) By Theorem~\ref{thm:hmeasure}, every $\nu>46.4471$ admits only finitely many $p/q$ with $|\lambda-p/q|<q^{-\nu}$; in particular $\nu=46.45$ does. \qedhere
\end{proof}
Note that the rounding is upward: the stated bound $46.45$ is larger than the computed $46.4471$. The sensitivity is $\partial\mu/\partial D_F=-223$ and $\partial\mu/\partial c_{\rm hyb}=6.07$.

\begin{remark}[a second certificate]
The certificate measure optimised for $D-0.04c$ (file \texttt{sig\_h500w.json}) gives $D_F\ge3.9045139$, $c\le7.5928685$ and the worse exponent $46.468$; the optimal weight on $c$ is about $0.027$, so a re-optimisation could gain only marginally.
\end{remark}

\subsection{What was proved by computer, and what was checked numerically}\label{ssec:hcomp}
\begin{itemize}[leftmargin=*]
\item \textbf{Computer-assisted (interval arithmetic).} The numbers $m'$, $S_{\rm cert}$, $g_{\min}$, $I_\rho(\sigma)$, $I(\nu_{1/8})$ in \S\ref{ssec:hcert} (script \texttt{hcert.py}, Arb); everything else in the deduction is pen-and-paper, with explicit constants where stated.
\item \textbf{Exact checks (rational arithmetic), not used in the proofs.} Reduction mod $\ell$ (Proposition~\ref{prop:hred}) including the constant, at every residue of $X$ for $\ell=11,\dots,43$ and several geometries; the product formula of Theorem~\ref{thm:hK} (exact polynomial interpolation at $D+1$ points for all $m$ at $K=11,\dots,14$, up to $D=3458$; $K\le10$ via the closed form of the exponents); $k_R(0),k_L(0)\bmod\ell$ for all $\ell\equiv3\ (4)$, $11\le\ell<800$ and every $m$, and $\ell<3000$ in the $\alpha$-window; end-to-end: $\det\Fu(x_0)\neq0$ for one of the two members for every residue class of $x_0$ and random rationals, $\ell=7,11,19,43,59,83,107,131$, no failure; Lemma~E and the content: $v_p(\den_F)=\tau_p-m$ exactly at every clean prime for $r=24,48,56,72$ (formal $X$), and for the member with the extra node ($e=1$) at $r=48,64,80,96$ (the script \texttt{member2\_content.py}); independently, $v_p(\den_F)$ equals the bound of Theorem~\ref{thm:hE} in all $27$ cases $(r,m,e,p)$ of Remark~\ref{rem:hmember2} (\texttt{check\_hE\_member2.py}); the bound \eqref{eq:hB} with $C_0=0$ for all odd $q$ at $21$ pairs $(r,m)$ (\texttt{check\_hB.py}); Proposition~\ref{prop:hL3} in $3086$ valuation checks.
\item \textbf{Cross-checks of $c_{\rm hyb}$ and $D_F$.} Exact values of $-r^{-2}\log|\det\mathcal T_t|$ extrapolate to $D\approx3.914$ (the certificate loses $0.008$), see Remark in \S\ref{ssec:hcert}.
\end{itemize}

\subsection{Status of the individual steps, and limitations}\label{ssec:htodo}
We list the points at which the written proof rests on something other than the audited arguments of \S\ref{sec:formal}.
\begin{enumerate}[leftmargin=*,label=\textbf{S\arabic*.}]
\item \emph{Constants of $\eta(r)$.} The $o(r^2)$ term in Theorem~\ref{thm:hN} is explicit, \eqref{eq:heta}: $\eta(r)\le(15\log r+K_0)/r$ with $K_0\le25{,}100$, for $r\ge2^{30}$. The Lipschitz bounds in $\alpha$ (Lemma~\ref{lem:hB}), the treatment of the coefficient discrepancy $e/r$ (the function $G^{(m)}$), the constant $C_E$ and the extra node $b$ are proved in \S\ref{sec:hdecay}; the dominant constant $4\sum_\ell(m'-\gamma_\ell)^+$ is read off the certificate data. These constants were re-derived by hand after the audits; they affect only the $o(1)$.
\item \emph{Member~2 (the extra common node).} The content saving for the geometry with the extra common node $b$ is proved in Theorem~\ref{thm:hE} (case $e=1$); this proof is new in this text. It was re-derived line by line by a further independent check (Remark~\ref{rem:hmember2}, which also explains why the extra node changes the content by only $O(r)$), and it is confirmed by two independent exact computations, in which the bound of Theorem~\ref{thm:hE} is attained with equality. For the decay the extra node is handled by the perturbation argument of \S\ref{sec:hdecay}.
\item \emph{Per-prime bound \eqref{eq:hB}.} Proved by hand, with $C_0=0$; this agrees with the exact evaluation.
\item \emph{Classical external inputs not proved here:} the norms $h_n$ of the discrete Chebyshev polynomials and the Hahn polynomial formulas at negative parameters, cited from \cite[\S9.5]{KLS}; the latter are used through the rational-identity (continuation) argument of \S\ref{ssec:hlead}: for fixed $N,n$ they are identities between rational functions of $(\alpha,\beta)$, valid on $\alpha,\beta>-1$ and hence at every point without vanishing denominators, which is the case for our parameters. Also Rellich's theorem and first-order degenerate perturbation theory (Kato), and the two-sided prime number theorem in progressions.
\item \emph{The zero orders for $K<c\le2K-1$.} The gap condition (G$'$) of \S\ref{ssec:hz2} is not assumed: the deformation of the free $\nu$-nodes removes it. The deformation argument was audited (with a correction of the denominator $\Delta$ by the factors $\prod(y_i-o)$, included above); in addition the counts of Rellich orders were verified by computer for $K\le300$, and the formulas (not the proof) of Proposition~\ref{prop:hprimes} were verified for $K\le18$ and in $20$ exact cases.
\item \emph{The certificate.} The copy of \texttt{hcert.py} in \texttt{scripts-hybrid/} has been re-run and reproduces $D_F\ge3.906269506610271$ and $c_{\rm hyb}\le7.654064744147046$ exactly (\S\ref{ssec:hcert}).
\end{enumerate}

\paragraph{Limitations.} The following is not done in this paper: the hybrid for the mirror constant $\lambda'=G/2+5\pi^2/96$ (primes $\ell\equiv1\pmod4$), for which Theorem~\ref{thm:mirror} ($57.85$) is unchanged; values $\alpha\ne\frac18$ (the saving formula holds only for $\frac18\le\alpha\le\frac16$, the tropical law as stated only for $\alpha\le\frac18$); the savings from non-clean primes and from primes $p<\frac58r$, which probably exist but are neither proved nor used; and a Lean formalization of any part of Part~\ref{part:hybrid}.

\subsection{Reproducing the hybrid computations}\label{ssec:hrepro2}
The directory \texttt{scripts-hybrid/} (in the repository next to \texttt{paper-lambda/scripts/}) contains: \texttt{hcert.py} (the Arb certificate of \S\ref{ssec:hcert}); \texttt{sig\_h500.json} (the certificate measure), \texttt{sig\_h500w.json} (second measure), \texttt{sigma300.json} (negative control); \texttt{expected\_h500.log}, \texttt{expected\_ctl0.log} (expected output); \texttt{member2\_content.py} (exact content check for the extra-node member), \texttt{check\_hE\_member2.py} (independent exact check of Theorem~\ref{thm:hE} for both members), \texttt{check\_hB.py} (the bound \eqref{eq:hB}); and a \texttt{README} with the commands of \S\ref{ssec:hrepro}. The value $D_F\ge3.906269506610$, $c\le7.654064744147$ appears in the last lines of \texttt{expected\_h500.log}. The directory also contains the exact checks written for this part: \texttt{check\_hE.py} (Theorem~\ref{thm:hE} with formal $X$ for admissible $(r,m,e)$, $e\in\{0,1\}$: $27$ primes, $0$ failures, several with equality), \texttt{check\_k.py} (the product formula of Theorem~\ref{thm:hK} with the exponents \eqref{eq:hec}, exact, for seven pairs $(K,m)$), \texttt{sum\_ec.py} ($\sum e_c=D$ for $K\le200$ and agreement of \eqref{eq:hec} with the min-formula of our earlier derivation for $K<80$) and \texttt{bsum.py} (the sum of the per-prime bounds of Proposition~\ref{prop:hdual}: $3.6497,3.7227,3.7418,3.7458$ at $r=10^3,10^4,10^5,4\cdot10^5$, converging to $\Lambda_{\rm trop}=3.7499$).


%%% --------------------------------------------------------------------
\section{Formalization status}\label{sec:formal}

We state this carefully. \textbf{The hybrid theorem (Theorem~\ref{thm:hybrid}, $\mu(\lambda)\le46.45$) and everything in Part~\ref{part:hybrid} are NOT formalized in Lean.} The Lean results cover only the $57.85$-line results of \S\ref{sec:matrix}--\S\ref{sec:mirror}, namely the weaker exponents $77.28$ ($\lambda$), $78$ ($\lambda'$) and $79$ ($G+5\pi^2/48$) below. The companion paper \cite{Litsey2026} on Catalan\'s constant is fully formalized in Lean~4, and so are the irrationality of $\lambda$ and the bound $\mu(\lambda)\le77.28$ proved here; Theorem~A, Lemma~CC and the bound $57.85$ are not. Each of the Lean results below compiles without \texttt{sorry} and depends only on the axioms \texttt{propext}, \texttt{Classical.choice} and \texttt{Quot.sound}. The library uses the Lean module system; a full build of the $\lambda$ library takes about $14$ minutes, and the kernel-checked numerical certificate about one CPU-hour.

\begin{itemize}[leftmargin=*]
\item \textbf{Formalized (Lean~4 and Mathlib): nonvanishing.} The $\ell$-adic nonvanishing theorem, Theorem~\ref{thm:nonvanish}: for a prime $\ell\equiv3\pmod 4$, $\ell\ge7$, and $r\in\Q$ with $\ell\nmid\den r$, one of $P_{\ell-1}(r),P_{\ell-2}(r),P_{(3\ell-1)/4}(r)$ is nonzero. The theorem is \lean{CatalanLambda.lambda_nonvanishing}. The formalized proof is organized slightly differently from \S\ref{sec:nonvanishing}: it works in the subring $\Z_{(\ell)}$ with the reduction map to $\Z/\ell$, uses the telescoping $\frac1k-\frac1{k+\ell-K}=-\frac1{k(4k-1)}$ directly, and absorbs the constant $-H^{(2)}_m$ into a single constant. It does not formalize the valuation statement $v_\ell=-2(h-K)$.
\item \textbf{Formalized: the prime number theorem in arithmetic progressions.} The theorem for every modulus, used in Theorems~\ref{thm:content} and~\ref{thm:measure}, as \lean{CatalanLambda.pnt_AP}, together with the upper bound \lean{CatalanLambda.theta_mod4_upper}. It is derived from the vendored Wiener--Ikehara theorem and Mathlib's results on Dirichlet $L$-functions.
\item \textbf{Formalized: the content bound.} The integrality estimate (Lemma~\ref{lem:integrality}) as \lean{CatalanLambda.integrality}, and the denominator bound of Theorem~\ref{thm:content} in the form \lean{CatalanLambda.log_Dden_le_eventually}.
\item \textbf{Formalized: the L\"owner form.} The off-diagonal and diagonal formulas for the matrices $T_h$ as L\"owner matrices (Proposition~\ref{prop:lowner}), as \lean{CatalanLambda.Decay.lowner_offdiag} and \lean{CatalanLambda.Decay.lowner_diag}, and the identity $\lambda=G/2-5\pi^2/96$ as \lean{CatalanLambda.Decay.lamS_eq}.
\item \textbf{Formalized: the contour identity.} Proposition~\ref{prop:contour} as \lean{CatalanLambda.Decay.contour_identity}.
\item \textbf{Formalized: Andr\'eief and Cauchy.} The Andr\'eief-type determinant bound (Proposition~\ref{prop:andreief}) as \lean{CatalanLambda.Decay.tent_andreief_bound}, and the Cauchy determinant as \lean{CatalanLambda.Decay.cauchy_det}.
\item \textbf{Formalized: the kernel $k_\rho$, Theorem~\ref{thm:N} and the deduction.} The positive-definiteness lemma for the kernel $k_\rho$ (Lemmas~\ref{lem:F}, \ref{lem:G}, \ref{lem:K} and Corollary~\ref{cor:Kp}); the contour identity for a general ray angle (including the direction $(-1,3)$ used here); Theorem~\ref{thm:N} (node removal $N_\rho$), as \lean{CatalanLambda.Decay.theoremN}, together with \lean{CatalanLambda.Decay.window_liminf_P}; Lemma~\ref{lem:3} as \lean{CatalanLambda.Measure.lemma3_pq}; and the deduction of the irrationality and of the bound of Remark~\ref{rem:77} from the certificate facts, as \lean{CatalanLambda.Measure.lambda_irrational} and \lean{CatalanLambda.Measure.lambda_measure_one}.
\item \textbf{Formalized: the numerical certificate.} The structure \lean{CatalanLambda.Cert.certFacts} collecting the numerical facts of \S\ref{ssec:cert} and \S\ref{ssec:cert2} (the certified values of $D_{\rm cert}$, $m^\prime$, $S_{\rm cert}$, $g_{\min}$) is inhabited by a kernel-checked proof: about $27.6$k ray boxes, $2.8$k segment cells and $512$ line cells, the bound on $I_\rho$ and the tail estimate. The proof uses a far-field first-order Taylor model with verified rational tables for the logarithm and the arctangent, and takes about one CPU-hour in the kernel. Lean works with the short rational constants $D_{\rm cert}=31111/10000$ and $c_{\rm cert}=68726/10000$ obtained by rounding the certified values, which is why it proves $77.28$ rather than $77.13$.
\item \textbf{Formalized: the final theorems.} \lean{CatalanLambda.lambda_irrational_final} ($\lambda$ is irrational), \lean{CatalanLambda.catalan_sub_irrational} ($G-5\pi^2/48$ is irrational), \lean{CatalanLambda.lambda_measure_final} ($\mu(\lambda)\le31682445459/409965000\approx77.28$) and \lean{CatalanLambda.lambda_measure_78_catalan} (the round exponent $78$), all unconditional. The statements are in \texttt{ChallengeLambda.lean}, which needs only Mathlib and writes $G$ as the series $\sum_{j\ge0}(-1)^j/(2j+1)^2$; the proofs are in \texttt{SolutionLambda.lean}.
\item \textbf{Formalized: the mirror theorem (weaker exponent).} The irrationality of $\lambda'=G/2+5\pi^2/96$ and of $G+5\pi^2/48$, and the measure with the same exponent $31682445459/409965000\approx77.28$ (hence $78$; and $79$ for $G+5\pi^2/48$) are formalized in \lean{CatalanLambda/Mirror/}, with the three-$h$ nonvanishing for $\ell\equiv1\pmod4$ (no Theorem~A), the mirror contour identity, and the \emph{same} kernel-checked certificate \lean{CatalanLambda.Cert.certFacts}. The final statements are \lean{CatalanLambda.Mirror.lambda_mirror_measure_78_catalan}, \lean{CatalanLambda.Mirror.catalan_add_irrational} and \lean{CatalanLambda.Mirror.catalan_add_measure_79}, with challenge statements \lean{CatalanLambdaChallenge.lambdaMirror_irrationalityMeasure_78} and \lean{CatalanLambdaChallenge.catalan_add_irrationalityMeasure_79}; axioms \texttt{propext}, \texttt{Classical.choice}, \texttt{Quot.sound}. The bound $57.85$ for $\lambda'$ (via Theorem~A) is not formalized.
\item \textbf{Not formalized.} Theorem~A, Lemma~CC and Theorem~\ref{thm:nonvanish2} (\S\ref{sec:krylov}), the deduction of the bound $57.85$ in Theorem~\ref{thm:measure}, and \emph{all of Part~\ref{part:hybrid}} (the hybrid theorem $\mu(\lambda)\le46.45$ with its tropical content bound, Lemma~E, the two-family Andr\'eief estimates, the bordered Krylov determinants and the $\ell\mid\den$ analogue; the Arb certificate \texttt{hcert.py} is not re-proved in Lean either). The Lean statement is the weaker bound $77.28$, obtained with the simpler choice of $h$ (Remark~\ref{rem:77}) and the rounded constants.
\item \textbf{Computer-assisted.} Theorem~\ref{thm:decay} and Theorem~\ref{thm:N} were established by interval-arithmetic computations in Arb (\S\ref{ssec:cert} and \S\ref{ssec:cert2}), carried out by the Python script \texttt{kcert.py} on the certificate measure \texttt{sigma300.json} (produced from a float optimum by the dyadic rounding script \texttt{roundsig.py}; \texttt{mkdata.py} re-derives the data in exact rationals; \texttt{scan.py} is the supplementary content scan). The script was re-run by an independent audit agent, which reproduced the numbers of \S\ref{ssec:cert} and sampled the certified cell bounds against floating-point values at about $5\times10^5$ ray points, $4\times10^4$ segment points and $10^4$ points of the line; but this sampling is not a proof; the certificate facts themselves are re-proved in Lean (the item on the numerical certificate above), by a separate kernel-checked computation. The float construction of $\sigma$ need not be correct since the certificate is checked as it stands. The scripts are archived in the directory \texttt{paper-lambda/scripts/} of \FormalizationLink.
\end{itemize}
The Lean sources are in the directory \texttt{CatalanLambda/} of the repository \FormalizationLink, with the files \texttt{ChallengeLambda.lean} and \texttt{SolutionLambda.lean} at its top level.

All other claims were audited by independent Claude (Anthropic) agents who re-derived the statements by hand and re-ran the numerical checks: the content bound; Lemmas~\ref{lem:1}, \ref{lem:2} and Theorem~\ref{thm:nonvanish} (exact characteristic polynomials at $h=8,12,16,20$ and all admissible $\ell$; $\gcd=1$ for $\ell=7,\dots,383$; no violation of the tropical bound against the exact $v_p(\den P_h)$ at $h=32,64$); Lemmas~\ref{lem:F}, \ref{lem:G}, \ref{lem:K} and Corollary~\ref{cor:Kp} (by hand, with a random numerical test of the inequality in 2000 trials); the contour identity for the direction $(-1,3)$ (by hand and to $10^{-17}$ numerically); Lemma~\ref{lem:B} and Theorem~\ref{thm:N}, checked line by line against the circle-smoothing version; the certificate re-run (reproducing $D_{\rm cert}\ge3.11137905$, $c\le6.87179932$); Lemma~\ref{lem:3} ($1044$ exact cases); and the simpler deduction of Remark~\ref{rem:77} ($4/(3\eta)+1=11.2228582$, $77.12123<77.13$). The new material of \S\ref{sec:krylov} and Theorem~\ref{thm:measure} was audited by a further independent Claude agent (full pass): Theorem~A (including the sign) was checked in exact rational arithmetic at $90$ random rational pairs $(\theta,s)$, including non-integer and negative $s$, for $n\le10$, and against the direct Cauchy Gram matrices ($\theta=0$, $s\le2n+1$, $n\le10$); $u(K,K)$ modulo $\ell$ is nonzero and equals the closed form for all $86$ primes $\ell\equiv3\pmod4$ with $7\le\ell<1000$; the resultant identity (R) was checked exactly; the characteristic polynomials of $N_K$, $N_{K+1}$ are coprime for $39$ primes $\ell<400$; Lemma~\ref{lem:2} was checked for $\ell=7,11,19,23$ and $h=\ell-K,\ell-K-1$; and the numbers $6.8717994\,(1/0.1304266+0.75)=57.84095$, $62.9948$ (without Lemma~\ref{lem:3}) and $77.1213$ were recomputed. The formula of Theorem~A was also verified exactly for every $n\le12$. Exact computation of $u(K,K)$ modulo $\ell$ for all $\ell\equiv3\pmod 4$, $7\le\ell<4500$, gave no zero. The deduction of Theorem~\ref{thm:measure} is written out in \S\ref{sec:measure} from the audited arithmetic. The audit listed only cosmetic issues, which are corrected in this text (notably the factor $4/h$ in $\eta$).

The mirror section \S\ref{sec:mirror} was audited by a separate Claude agent (full pass, no fatal issues; four cosmetic points, corrected here). Lemmas~\ref{lem:mstruct}--\ref{lem:m2} were re-derived by hand (Schur complement with high block $-I$); $N_0$ was checked against $-(\lambda'I+\bar T_h)$ to $10^{-51}$ at $h=9$; the reflection $\bar F=\tilde F+\bar E$ was checked by hand and to $10^{-22}$, and the contour identity independently to $2\cdot10^{-17}$, with the atom set compared to the poles right of $x^*$; the valuation $2(h-K)$ and degree $K$ were checked for all $131$ admissible $(\ell,h)$ with $\ell=13,17,29,37$; the characteristic polynomials of $N_K,N_{K+1}$ are coprime over $\mathbb F_\ell$ for all $58$ primes $\ell\equiv1\pmod4$, $13\le\ell<700$; Lemma~\ref{lem:m3} deviates only for $e=2$, $u\equiv1$; $v_p(\den\bar P_h)\le\sum_e\delta(p^e)$ for all $p\ge5$ at $h=16,32,64,128$; and $6.8717994\,(1/0.13042662+0.75)=57.84094$, $77.12126$ were recomputed. This is not a substitute for human refereeing.

The material of Part~\ref{part:hybrid} was audited by the following independent Claude agents (each re-derived the statements by hand and re-ran exact computations with its own code): the reduction modulo $\ell$, the bordered Lemma~CC analogue and the selection of the pair of members (exact constant checked at every residue of $X$ for $\ell=11,\dots,43$, $0$ mismatches; end-to-end exact nonvanishing for $\ell=7,\dots,131$, $0$ failures; corrected the member pair to $\{r_1,r_1+1\}$, $r_1=\ell-K-m$); the closed form of the bordered Krylov determinants (the constant $\varkappa$: $20$ cases by interpolation, no product assumption; $\ell$-integrality bridge; $\sum e_c=D$, which we now prove by direct summation; the zero orders: Rellich counts for $K\le300$, $17{,}954{,}950$ cases, the product formula verified exactly for all $m$ at $K=11,\dots,14$, and the deformation argument, with the correction of the denominator $\Delta$ by the factors $\prod(y_i-o)$); Lemma~E (hand check, exact $X$-formal checks at $r=24,48,56,72$, specialised checks at $X\in\{\frac2{13},-7,5\}$, and the observation that the statement fails at specialised $X\equiv H^{(2)}_{K_p-1}/16$, Remark~\ref{rem:hformal}); the tropical law (the duals D1--D6 re-derived and checked against certified Hungarian optima for $1300$ prime powers, the integral by quadrature to $10^{-15}$, the two-sided PNT paragraph of Lemma~\ref{lem:hS}); the decay certificate (the Arb output reproduced exactly: $D_F\ge3.906269506610$, $c\le7.654064744147$; negative control $\alpha=0$ reproduces $3.1113790$/$6.8717994$; closed forms compared with quadrature; the window falsification test of the line $D_F-\beta c_{\rm hyb}$ against exact profiles for $r\le128$); and the $\ell\mid\den$ analogue ($3086$ exact valuation checks, $0$ failures). The audits found only cosmetic points, all corrected here, with the exceptions listed in \S\ref{ssec:htodo}.

%%% --------------------------------------------------------------------
\section*{Acknowledgments and provenance}
The irrationality of Catalan's constant, and the Lean formalization on which the author's
earlier work builds, are due to OpenAI \cite{OpenAI2026}. The mathematics of this paper
---the matrices $T_h$, the content, decay, nonvanishing and node-removal arguments, the numerical
certificates and the text---were derived by Claude (Anthropic) agents working at the request
of the author, who posed the problem, directed the work, and is responsible for this
submission; independent Claude agents performed the audits described in \S\ref{sec:formal}. The same applies to Part~\ref{part:hybrid}: the rectangular hybrid frame, the content, decay, bordered nonvanishing and $\ell\mid\den$ arguments and the numerical certificate $\mu(\lambda)\le46.45$ were derived by Claude agents at the request of the author, and audited by further Claude agents; the hybrid theorem is not formalized in Lean.
The Lean formalization was likewise written by Claude agents working with the author.
The repository infrastructure, vendored formalizations and certificate tooling of the companion project, which this work reuses, were developed with OpenAI Codex as described in \cite{Litsey2026}.
We use Lean~4 \cite{Lean4}, Mathlib \cite{Mathlib}, the PNT+ project \cite{PNTPlus} (whose Wiener--Ikehara theorem we use to formalize the prime number theorem in progressions), and Arb \cite{Arb} for interval arithmetic.
The use of AI in this work is disclosed here in the same spirit as in \cite{Litsey2026}.

\begin{thebibliography}{99}

\bibitem{Arb}
F.~Johansson,
\emph{Arb: efficient arbitrary-precision midpoint-radius interval arithmetic},
IEEE Trans.\ Comput.\ \textbf{66} (2017), 1281--1292.

\bibitem{Carreira2026}
Marcos Costa Santos Carreira,
\emph{General designs and an irrationality exponent for Catalan's constant},
preprint v1 (October 9, 2026),
\href{https://doi.org/10.5281/zenodo.23326325}{doi:10.5281/zenodo.23326325}.

\bibitem{Davenport}
H.~Davenport,
\emph{Multiplicative Number Theory}, 3rd ed., Graduate Texts in Mathematics 74, Springer, 2000.

\bibitem{Forrester2018Andreief}
P.~J.~Forrester,
\emph{Meet Andr\'eief, Bordeaux 1886, and Andreev, Kharkov 1882--83},
preprint (2018), \href{https://arxiv.org/abs/1806.10411}{arXiv:1806.10411}.

\bibitem{Jha2026}
R.~K.~Jha,
\emph{A finite irrationality measure for Catalan's constant},
preprint v1, October 8, 2026,
\href{https://doi.org/10.5281/zenodo.23237269}{doi:10.5281/zenodo.23237269}.

\bibitem{KLS}
R.~Koekoek, P.~A.~Lesky and R.~F.~Swarttouw,
\emph{Hypergeometric Orthogonal Polynomials and Their $q$-Analogues}, Springer Monographs in Mathematics, Springer, 2010.

\bibitem{Litsey2026}
S.~Litsey,
\emph{An irrationality measure for Catalan's constant},
preprint, October 16, 2026; formalization at \FormalizationLink.

\bibitem{Mathlib}
The mathlib Community,
\emph{The Lean mathematical library},
in: Proceedings of the 9th ACM SIGPLAN International Conference on Certified Programs and
Proofs (CPP 2020), 367--381.

\bibitem{Lean4}
L.~de~Moura and S.~Ullrich,
\emph{The Lean~4 theorem prover and programming language},
in: Automated Deduction -- CADE 28, Lecture Notes in Computer Science 12699, Springer, 2021,
625--635.

\bibitem{MV}
H.~L.~Montgomery and R.~C.~Vaughan,
\emph{Multiplicative Number Theory I: Classical Theory}, Cambridge Studies in Advanced Mathematics 97, Cambridge University Press, 2007.

\bibitem{OpenAI2026}
OpenAI,
\emph{Catalan's constant is irrational},
preprint (September 24, 2026), with a Lean formalization, available at
\href{https://github.com/openai/math/blob/fd4aeeb2ee4fc729c18d98444fed42fd0529eeeb/preprints/Catalans-constant-is-irrational-September-24-2026/paper.pdf}{the original paper PDF}
and
\href{https://github.com/openai/math/tree/fd4aeeb2ee4fc729c18d98444fed42fd0529eeeb/lean/OAI/NumberTheory/Catalan}{the Lean formalization}
(repository commit \texttt{fd4aeeb}, October 7, 2026).

\bibitem{PNTPlus}
The PrimeNumberTheoremAnd contributors,
\emph{The PNT+ project: a Lean formalization of the prime number theorem
and related results},
\url{https://github.com/AlexKontorovich/PrimeNumberTheoremAnd}.

\bibitem{SaffTotik}
E.~B.~Saff and V.~Totik,
\emph{Logarithmic Potentials with External Fields}, Grundlehren der mathematischen Wissenschaften 316, Springer, 1997.

\bibitem{Sun2026}
Z.-W.~Sun,
\emph{Catalan's constant is irrational},
preprint (2026), \href{https://arxiv.org/abs/2609.04176}{arXiv:2609.04176}.

\bibitem{Zudilin}
W.~Zudilin,
\emph{A determinantal approach to irrationality},
Constr.\ Approx.\ \textbf{45} (2017), 301--310; \href{https://arxiv.org/abs/1507.05697}{arXiv:1507.05697} (2015).

\bibitem{Wachs2026}
D.~Wachs,
\emph{A note on a recent claimed proof of the irrationality of Catalan's constant},
preprint (2026), \href{https://arxiv.org/abs/2609.22339}{arXiv:2609.22339}.

\end{thebibliography}

\end{document}
